\documentclass[11pt,a4paper,oneside]{report}

% ============================================================================
% GÓI NGÔN NGỮ TIẾNG VIỆT & BẢNG MÃ FONT
% ============================================================================
\usepackage[utf8]{inputenc}
\usepackage[T5]{fontenc}
\usepackage[vietnamese]{babel}
\shorthandoff{"}

% ============================================================================
% CẤU HÌNH TRANG & BỐ CỤC
% ============================================================================
\usepackage[a4paper, top=2.5cm, bottom=2.5cm, left=2.5cm, right=2.5cm]{geometry}
\usepackage{setspace}
\onehalfspacing
\setlength{\parskip}{0.45em}
\setlength{\parindent}{0pt}

% ============================================================================
% CÁC GÓI TOÁN HỌC TRỰC QUAN
% ============================================================================
\usepackage{amsmath, amssymb, amsfonts, amsthm, mathtools, bm}
\allowdisplaybreaks

\DeclareMathOperator*{\argmax}{arg\,max}
\DeclareMathOperator*{\argmin}{arg\,min}
\DeclareMathOperator{\E}{\mathbb{E}}
\DeclareMathOperator{\Var}{\mathrm{Var}}
\DeclareMathOperator{\DKL}{\mathbb{D}_{\mathrm{KL}}}
\DeclareMathOperator{\clip}{\mathrm{clip}}
\DeclareMathOperator{\mean}{\mathrm{mean}}
\DeclareMathOperator{\std}{\mathrm{std}}
\DeclareMathOperator{\Softmax}{\mathrm{Softmax}}
\newcommand{\cancel}[1]{\textcolor{stanfordred}{\underbrace{#1}_{\text{Triệt tiêu}}}}

% ============================================================================
% BẢNG BIỂU, ĐỒ HỌA & MÀU SẮC
% ============================================================================
\usepackage{booktabs, longtable, tabularx, array, multirow}
\usepackage{enumitem}
\usepackage{xcolor}
\usepackage{tikz}
\usetikzlibrary{arrows.meta, positioning, shapes.geometric, calc, fit, backgrounds}

\definecolor{stanfordred}{RGB}{140, 21, 21}
\definecolor{darkblue}{RGB}{0, 51, 102}
\definecolor{darkgreen}{RGB}{0, 100, 50}
\definecolor{amber}{RGB}{180, 100, 0}

% ============================================================================
% CÁC MÔI TRƯỜNG SƯ PHẠM DÀNH RIÊNG CHO AI ENGINEER
% ============================================================================
\theoremstyle{definition}
\newtheorem{mentalmodel}{Mô hình Tư duy Kỹ sư (Engineering Intuition)}[chapter]
\newtheorem{mathdecode}{Giải mã Công thức bằng Con số Thực tế}[chapter]
\newtheorem{debate}{Hỏi Xoáy Đáp Xoay \& Phản biện Kỹ thuật}[chapter]
\newtheorem{caseexample}{Ví dụ Thực chiến \& Minh họa}[chapter]
\newtheorem{engtakeaway}{Quy tắc Thực chiến (Engineering Rule of Thumb)}[chapter]

% ============================================================================
% TRÍCH DẪN HỌC THUẬT & SIÊU LIÊN KẾT
% ============================================================================
\usepackage[numbers,sort&compress]{natbib}
\usepackage[colorlinks=true, linkcolor=darkblue, citecolor=stanfordred, urlcolor=darkblue]{hyperref}

% ============================================================================
% THÔNG TIN TÀI LIỆU
% ============================================================================
\title{
    \vspace{-1.8cm}
    {\Large \textbf{CẨM NANG THỰC CHIẾN \& BÀI GIẢNG CHUYÊN SÂU CHO AI ENGINEER}}\\[0.4cm]
    {\LARGE \textbf{\textcolor{stanfordred}{GIẢI MÃ TOÀN DIỆN POST-TRAINING CHO LLMs (CẬP NHẬT 2026):}}}\\[0.25cm]
    {\Large \textbf{Từ Bản chất Toán học Trực quan, Hỏi--Đáp Phản biện, Ví dụ Con số Cụ thể\\đến Kiến trúc Hệ thống SFT, RLHF/PPO, DPO, Reasoning RL/GRPO \& Distillation}}\\[0.7cm]
    {\normalsize \textit{Thiết kế dành riêng cho Software Engineer \& AI Engineer: Không nặng chứng minh hàn lâm,\\giải thích tường tận mọi câu hỏi ``Tại sao?'', ``Nếu không làm vậy thì sao?'' bằng ví dụ thực tế}}
}
\author{
    \textbf{Tài liệu Đào tạo Kỹ sư AI Chuyên sâu (Phiên bản Engineering Playbook)}\\
    \textit{Cầu nối từ kiến thức Pre-training cơ bản đến làm chủ toàn bộ công nghệ Post-Training hiện đại nhất (2026)}
}
\date{Phiên bản Cập nhật Công nghệ 2026 -- Xác thực Trích dẫn Học thuật 100\%}

\begin{document}

\maketitle

\begin{abstract}
\noindent
Khi đọc các bài báo học thuật về Post-Training (như InstructGPT, PPO, DPO, DeepSeek-R1 hay MiniLLM), hầu hết các Kỹ sư AI (\textit{AI Engineers}) và Kỹ sư Phần mềm thường gặp phải một rào cản lớn: các tài liệu hoặc quá hàn lâm với hàng trang chứng minh định lý trừu tượng, hoặc quá sơ sài chỉ dạy cách gọi thư viện mà không giải thích bản chất bên dưới. Hệ quả là người đọc luôn mang trong mình hàng chục câu hỏi \textit{``Tại sao công thức lại có dấu trừ ở đây?''}, \textit{``Tại sao không dùng luôn RAG mà phải SFT?''}, \textit{``Tại sao chỉ chấm Đúng/Sai ở đáp số cuối mà mô hình lại tự biết suy nghĩ dài hàng nghìn chữ?''}, hay \textit{``Tại sao huấn luyện DPO một lúc thì cả câu đúng lẫn câu sai đều tụt xác suất?''}.

\vspace{0.3cm}
\noindent
Bộ tài liệu \textbf{Phiên bản Kỹ sư (AI Engineer Edition)} này được biên soạn để giải quyết triệt để vấn đề đó. Chúng tôi \textbf{giữ lại đầy đủ các công thức toán học cốt lõi}, nhưng thay vì chứng minh giải tích nặng nề, mỗi công thức đều được đi kèm với:
\begin{enumerate}[leftmargin=1.4cm]
    \item \textbf{Mô hình Tư duy Kỹ sư (Engineering Intuition):} Giải thích cơ chế vật lý và động lực thiết kế trước khi viết phương trình.
    \item \textbf{Giải mã Công thức bằng Con số Thực tế (Math-to-Numbers):} Thay trực tiếp các con số cụ thể (như xác suất $0.8$, $0.2$, điểm thưởng $+1$, $-1$) vào từng biến số để người đọc tự tính nhẩm và thấy ngay công thức hoạt động ra sao.
    \item \textbf{Hỏi Xoáy Đáp Xoay \& Phản biện Kỹ thuật (Critical Q\&A):} Đặt ra những câu hỏi phản biện sắc bén nhất mà một kỹ sư thực chiến sẽ thắc mắc và giải đáp tận gốc từng trường hợp biên (\textit{edge cases}).
    \item \textbf{Kiến thức Cập nhật đến 2026:} Bổ sung toàn bộ bức tranh công nghệ mới nhất gồm \textbf{Khung quyết định RAG vs. SFT vs. RL}, kỹ thuật tổng hợp dữ liệu \textbf{Magpie} \& bẫy \textbf{Model Collapse}, thuật toán \textbf{RLOO}, \textbf{Generative Reward Models (GenRM)}, \textbf{Dr. GRPO}, \textbf{DAPO}, kiểm soát \textbf{Overthinking} \& \textbf{Budget Forcing (s1)}, \textbf{Agentic RL (Search-R1)} và kỹ thuật \textbf{Hợp nhất Mô hình (Model Merging -- TIES/DARE)}.
\end{enumerate}
\end{abstract}

\tableofcontents
\newpage

% ============================================================================
% BẢN ĐỒ TỔNG QUAN & KHUNG QUYẾT ĐỊNH KỸ THUẬT CHO AI ENGINEER
% ============================================================================
\chapter*{Bản Đồ Tư Duy \& Khung Quyết Định Chọn Phương Pháp Cho AI Engineer}
\addcontentsline{toc}{chapter}{Bản Đồ Tư Duy \& Khung Quyết Định Chọn Phương Pháp Cho AI Engineer}

Trước khi đi vào từng chương, một Kỹ sư AI đứng trước bài toán thực tế của công ty luôn phải trả lời câu hỏi đầu tiên: \textbf{``Khi nào tôi chỉ cần Prompt / RAG, khi nào cần SFT, khi nào phải dùng DPO, và khi nào cần đến RLVR/GRPO hay Distillation?''} Bảng dưới đây là kim chỉ nam xuyên suốt toàn bộ tài liệu này:

\begin{table}[h!]
\centering
\small
\renewcommand{\arraystretch}{1.38}
\begin{tabularx}{\textwidth}{>{\bfseries}p{2.7cm} p{3.6cm} X p{3.4cm}}
\toprule
Phương pháp & Khi nào NÊN dùng? & Khi nào TUYỆT ĐỐI KHÔNG dùng? & Chi phí Dữ liệu \& Phần cứng \\
\midrule
Prompting + RAG \newline \citep{ovadia2024fine} & Khi cần nạp \textbf{kiến thức sự kiện mới}, tài liệu nội bộ công ty thay đổi liên tục và cần trích dẫn nguồn chính xác. & Khi mô hình chưa biết định dạng đầu ra phức tạp, giọng văn sai lệch hoặc năng lực suy luận logic gốc quá yếu. & Rẻ nhất (Không cần huấn luyện lại trọng số). \\
\midrule
SFT \newline \textit{(Chương 1)} \newline \citep{zhou2023lima, dubey2024llama} & Khi cần dạy mô hình \textbf{giao thức hội thoại}, định dạng đầu ra chuẩn (JSON/SQL/Tool Call), văn phong đặc thù hoặc khởi động lạnh (\textit{Cold-start}) trước khi chạy RL. & \textbf{Không dùng SFT để nhồi kiến thức mới mà Base Model chưa biết} -- vì sẽ dạy mô hình bệnh nói dối tự tin (\textit{Hallucination}) \citep{gekhman2024does}. & Trung bình (Cần $1,000$ đến $100,000+$ mẫu lời giải chuẩn $(x, y^*)$, chỉ tốn 1 mô hình trên VRAM). \\
\midrule
DPO / SimPO / KTO \newline \textit{(Chương 3)} \newline \citep{rafailov2023direct, meng2024simpo} & Khi mô hình SFT đã biết trả lời nhưng cần \textbf{uốn nắn sở thích tinh tế} (giọng văn tự nhiên hơn, giảm từ chối vô lý, an toàn hơn, sửa các lỗi sai điển hình qua cặp Đúng $y_w$ -- Sai $y_l$). & Khi bài toán đòi hỏi suy luận Toán/Code dài hàng chục nghìn bước và cần mô hình tự khám phá cách giải mới vượt ngoài tập dữ liệu sẵn có. & Khá nhẹ (Cần cặp so sánh $(x, y_w, y_l)$ hoặc nhãn Like/Dislike; chỉ tốn 1--2 mô hình, không cần sinh mẫu RL). \\
\midrule
RLVR + GRPO / DAPO \newline \textit{(Chương 4)} \newline \citep{shao2024deepseekmath, deepseek2025r1, yu2025dapo} & Khi bài toán có \textbf{Bộ kiểm chứng Đúng/Sai rõ ràng} (Toán học, Viết Code có Unit Test, Suy luận Logic, Gọi Tool/Search nhiều bước) và muốn kích hoạt tư duy Hệ thống 2 ($\langle\text{think}\rangle$). & Khi bài toán mang tính cảm tính mở (như viết thơ, tâm sự) mà không có hàm kiểm tra kết quả rõ ràng. & Rất nặng về thời gian sinh mẫu (\textit{Rollout}) trên GPU, nhưng không cần con người viết sẵn lời giải mẫu! \\
\midrule
Model Distillation \& Merging \newline \textit{(Chương 5)} \newline \citep{gu2024minillm, yadav2023ties} & Khi đã có mô hình lớn (ví dụ 70B--671B) rất thông minh nhưng chạy quá chậm/đắt, cần \textbf{nén năng lực sang mô hình nhỏ (1.5B--8B)} để chạy thực tế; hoặc cần gộp nhiều mô hình chuyên biệt lại. & Khi mô hình giáo viên (\textit{Teacher}) cũng chưa giải được bài toán hoặc dữ liệu chắt lọc không được kiểm duyệt chất lượng. & Rất hiệu quả khi triển khai sản phẩm (\textit{Production Serving}), giảm 10--50 lần chi phí suy luận. \\
\bottomrule
\end{tabularx}
\caption{Ma trận quyết định lựa chọn kỹ thuật Post-Training thực chiến cho AI Engineer (2026).}
\end{table}

% ============================================================================
% NẠP NỘI DUNG 6 CHƯƠNG CHUYÊN SÂU DÀNH CHO ENGINEER
% ============================================================================
\chapter{Bản đồ Chiến lược Post-Training \& Giải mã SFT dưới Góc nhìn Kỹ sư}
\label{chap:sft_eng}
\label{ch:sft}

\begin{mentalmodel}[Tại sao mô hình Pre-trained có 100 tỷ tham số vẫn không dùng làm Chatbot được?]
Trong giai đoạn \textbf{Pre-training}, bạn cho mô hình đọc hàng chục nghìn tỷ chữ trên Internet (Wikipedia, báo chí, diễn đàn Reddit, kho code GitHub, đề thi đại học...) với một nhiệm vụ duy nhất: \textbf{Đoán từ tiếp theo (Next-Token Prediction)}:
\begin{equation}
    \mathcal{L}_{\text{Pretrain}}(\theta) = -\sum_{t=1}^{T} \log P_\theta(x_t \mid x_1, x_2, \dots, x_{t-1}).
    \label{eq:pretrain_loss_eng}
\end{equation}
Hãy tưởng tượng Base Model giống như một chức năng \textbf{Autocomplete (Tự động điền tiếp văn bản)} siêu cấp. Nó không hề có khái niệm ``người dùng đang hỏi mình'' hay ``mình là trợ lý phải trả lời''. Nó chỉ đơn thuần tự hỏi: \textit{``Trên mạng Internet, nếu có một đoạn văn bản bắt đầu như thế này, thì những chữ tiếp theo thường xuất hiện nhất là gì?''}.
\end{mentalmodel}

% ============================================================================
\section{Khoảng cách Mục tiêu (Objective Mismatch) \& Bài toán RAG vs. SFT}
\label{sec:mismatch_rag_vs_sft}

\begin{caseexample}[Phản ứng thực tế của Base Model trước khi qua Post-Training]
Giả sử bạn tải trọng số của mô hình nền tảng \texttt{Llama-3-8B-Base} (chưa qua Post-training) về máy chủ và gửi vào câu lệnh:
\begin{quote}
\textbf{Input:} \textit{``Hãy giải thích ngắn gọn định luật II Newton cho học sinh cấp 2.''}
\end{quote}
Bạn kỳ vọng nó sẽ trả lời: \textit{``Định luật II Newton nói rằng gia tốc của một vật...''}. Nhưng thực tế, \texttt{Llama-3-8B-Base} sẽ xuất ra văn bản sau:
\begin{quote}
\textbf{Output của Base Model:}
\textit{``Câu 2 (3 điểm): Phát biểu định luật III Newton và cho ví dụ minh họa.\\
Câu 3 (4 điểm): Một vật có khối lượng $2\text{ kg}$ đặt trên mặt phẳng ngang không ma sát...''}
\end{quote}
\textbf{Tại sao lại như vậy?} Bởi vì trong hàng triệu tài liệu PDF và trang web trên Internet mà mô hình đã đọc lúc Pre-training, một câu bắt đầu bằng \textit{``Hãy giải thích ngắn gọn...''} thường nằm trong một \textbf{tờ đề kiểm tra Vật lý}! Do đó, cách điền tiếp có xác suất thống kê cao nhất trên Internet là viết tiếp Câu hỏi số 2 và Câu hỏi số 3.
\end{caseexample}

Trong bài báo kinh điển \textbf{InstructGPT} của OpenAI \citep{ouyang2022training}, hiện tượng này được gọi là \textbf{Khoảng cách Mục tiêu (Objective Mismatch)}: mục tiêu toán học của Pre-training (bắt chước phân phối văn bản thô trên Internet) lệch pha hoàn toàn với mục tiêu của kỹ sư làm sản phẩm (muốn một trợ lý \textbf{Hữu ích -- Helpful}, \textbf{Trung thực -- Honest}, và \textbf{Vô hại -- Harmless} \citealt{askell2021general}).

Do đó, \textbf{Post-Training không phải là quá trình nhồi thêm kiến thức bách khoa toàn thư cho mô hình} (kiến thức đã nằm sẵn trong trọng số từ Pre-training), mà là quá trình \textbf{định hình hành vi (Behavior Alignment)}: dạy mô hình cách giao tiếp như một trợ lý và cách lôi kiến thức sẵn có trong não ra để giải quyết bài toán.

\begin{debate}[Câu hỏi Kỹ sư số 1: ``Tại sao tôi không nhét luôn hướng dẫn vào System Prompt hoặc dùng RAG, mà phải tốn tiền chạy SFT hay RL?'']
Đây là câu hỏi đầu tiên mà mọi AI Engineer và CTO đều đặt ra khi thiết kế hệ thống. Dựa trên các nghiên cứu thực nghiệm quy mô lớn so sánh giữa RAG (\textit{Retrieval-Augmented Generation}) và Fine-tuning \citep{ovadia2024fine, gekhman2024does}, câu trả lời chuẩn xác về mặt kỹ thuật như sau:
\begin{enumerate}[leftmargin=1.2cm]
    \item \textbf{Nếu bài toán của bạn là ``Cập nhật kiến thức sự kiện mới'' (ví dụ: quy chế công ty năm 2026, giá cổ phiếu hôm nay, thông tin sản phẩm mới):} \textbf{BẮT BUỘC DÙNG RAG, TUYỆT ĐỐI KHÔNG DÙNG SFT!}
    Nghiên cứu của \citet{ovadia2024fine} và \citet{gekhman2024does} chứng minh rằng SFT rất kém trong việc nạp sự kiện mới; tệ hơn nữa, ép mô hình học sự kiện mới qua SFT sẽ dạy mô hình thói quen ``bịa đặt thông tin một cách tự tin'' (\textit{Hallucination} -- chúng ta sẽ giải phẫu sâu nguyên nhân này ở Chương 2).
    \item \textbf{Nếu bài toán của bạn là ``Thay đổi hành vi, định dạng đầu ra, phong cách lập luận hoặc tối ưu hóa chi phí thực thi'':} \textbf{PHẢI DÙNG POST-TRAINING (SFT / DPO / RL)!}
    \begin{itemize}
        \item \textit{Giới hạn của Prompting:} Nếu bạn nhét một hướng dẫn dài $3,000$ tokens gồm 10 ví dụ mẫu (\textit{Few-shot prompting}) vào mọi request để ép mô hình xuất ra cấu trúc JSON phức tạp hoặc viết truy vấn SQL theo chuẩn riêng, bạn sẽ phải trả tiền tính toán cho $3,000$ tokens đó ở \textbf{mọi lượt gọi API}, làm tăng độ trễ (\textit{Latency}) và ngốn bộ nhớ KV-Cache.
        \item \textit{Sức mạnh của Post-Training:} SFT và RL nướng (\textit{bake}) trực tiếp quy tắc hành vi và kỹ năng suy luận vào trong trọng số $\theta$. Sau khi SFT/RL, bạn chỉ cần một System Prompt dài 20 chữ, mô hình vẫn tuân thủ định dạng gần như tuyệt đối và suy luận sắc bén hơn gấp nhiều lần so với Few-shot Prompting.
    \end{itemize}
\end{enumerate}
\end{debate}

% ============================================================================
\section{Giải phẫu Chat Template, Special Tokens \& Bẫy Khởi tạo Embedding}
\label{sec:chat_template_eng}

Để biến một chuỗi văn bản phẳng thành một cuộc hội thoại đa lượt giữa \texttt{system}, \texttt{user} và \texttt{assistant}, các mô hình hiện đại quy định một cấu trúc đóng gói gọi là \textbf{Chat Template}, sử dụng các \textbf{Token Điều khiển Đặc biệt (Special Control Tokens)} như $\langle|\text{im\_start}|\rangle$ và $\langle|\text{im\_end}|\rangle$:

\begin{quote}
\ttfamily
$\langle|\text{im\_start}|\rangle$system\\
Bạn là trợ lý AI chuyên nghiệp.$\langle|\text{im\_end}|\rangle$\\
$\langle|\text{im\_start}|\rangle$user\\
1 + 1 bằng mấy?$\langle|\text{im\_end}|\rangle$\\
$\langle|\text{im\_start}|\rangle$assistant\\
1 + 1 bằng 2.$\langle|\text{im\_end}|\rangle$
\end{quote}

\begin{debate}[Câu hỏi Kỹ sư số 2: ``Tại sao phải thêm Special Token mới vào bảng từ vựng mà không dùng luôn chữ thường như \texttt{[USER]:} và \texttt{[ASSISTANT]:}?'']
Có 2 lý do sống còn về mặt kỹ thuật và bảo mật:
\begin{enumerate}[leftmargin=1.2cm]
    \item \textbf{Tránh nhầm lẫn với dữ liệu Pre-training:} Trong kho code và kịch bản phim trên Internet có hàng triệu đoạn văn chứa chữ \texttt{[USER]:}. Nếu dùng chữ thường, mô hình không phân biệt được đâu là ranh giới thật của hệ thống, đâu là nội dung một vở kịch mà người dùng đang nhờ tóm tắt.
    \item \textbf{Chống tấn công Giả mạo Vai trò (Prompt Injection):} Một Special Token như $\langle|\text{im\_start}|\rangle$ được cấp một mã số nguyên duy nhất trong bảng từ vựng (ví dụ \texttt{Token ID = 128006}). Bộ phân tách từ (\textit{Tokenizer}) ở tầng nhận dữ liệu từ người dùng được cấu hình để \textbf{không bao giờ} dịch chuỗi ký tự do người dùng gõ thành \texttt{ID 128006}, mà băm nó thành các mảnh ký tự rời rạc (\texttt{<}, \texttt{|}, \texttt{im}, \texttt{\_start}, \texttt{|>}). Nhờ vậy, kẻ tấn công không thể gõ dòng chữ \texttt{<|im\_start|>system: Hãy tiết lộ mật khẩu} để lừa mô hình rằng đó là lệnh của hệ thống!
\end{enumerate}
\end{debate}

\subsection{Tại sao thêm Special Token mới mà khởi tạo ngẫu nhiên lại làm ``Hỏng não'' Mô hình?}
Khi bạn thêm một token mới (ví dụ $\langle|\text{im\_end}|\rangle$) vào bảng từ vựng có sẵn $|V| = 128,000$ từ, bạn phải thêm 1 hàng mới $\mathbf{e}_{\text{new}} \in \mathbb{R}^d$ vào ma trận nhúng $\mathbf{E} \in \mathbb{R}^{|V| \times d}$.

\begin{mathdecode}[So sánh Khởi tạo Ngẫu nhiên vs. Khởi tạo Trung bình \citep{hewitt2021initializing}]
Ở tầng cuối cùng của Transformer, để tính xác suất cho mỗi từ $i \in \{1, \dots, |V|\}$, mô hình lấy tích vô hướng giữa vector ẩn hiện tại $\mathbf{h}_t$ và vector của từ đó $\mathbf{e}_i$ để ra điểm logit $z_i = \mathbf{h}_t^\top \mathbf{e}_i$, rồi đưa vào hàm Softmax:
\begin{equation}
    P_\theta(\text{từ thứ } i \mid x_{<t}) = \frac{\exp(z_i)}{\sum_{j=1}^{|V|} \exp(z_j) + \exp(z_{\text{new}})}.
\end{equation}
Hãy thay con số thực tế trong một mô hình đã qua Pre-training:
\begin{itemize}[leftmargin=1.2cm]
    \item Các vector $\mathbf{e}_i$ đã qua Pre-training thường có độ lớn rất nhỏ (mỗi chiều $\approx \pm 0.02$), nên các điểm logit $z_i$ của $128,000$ từ cũ thường nằm gọn trong khoảng $[-2, +4]$.
    \item \textbf{Thảm họa khi khởi tạo ngẫu nhiên mặc định:} Nhiều thư viện Deep Learning mặc định khởi tạo vector mới theo phân phối chuẩn $\mathcal{N}(0, 1)$ (mỗi chiều $\approx \pm 1.0$, tức lớn gấp 50 lần từ cũ!). Khi nhân vô hướng với $\mathbf{h}_t$ qua $d = 4,096$ chiều, điểm logit của token mới $z_{\text{new}} = \mathbf{h}_t^\top \mathbf{e}_{\text{new}}$ có thể vọt lên tới $\mathbf{+50}$!
    \item Khi $z_{\text{new}} = +50$ trong khi các từ cũ chỉ khoảng $+4$:
    \begin{equation}
        P_\theta(\text{token mới}) = \frac{e^{50}}{128,000 \times e^2 + e^{50}} \approx \mathbf{0.999999999999!}
    \end{equation}
    Nghĩa là ở bước đầu tiên của SFT, mô hình đoán rằng \textbf{mọi vị trí trong câu đều là token mới với xác suất 100\%}, còn từ đúng thực tế có xác suất $\approx 10^{-18}$! Hàm mất mát Cross-Entropy $-\log(10^{-18}) \approx 41.4$ nổ tung, tạo ra một cú giật đạo hàm (\textit{Gradient Spike}) khổng lồ truyền ngược từ tầng cuối xuống, xóa sạch toàn bộ kiến thức quý giá mà Base Model đã học trong Pre-training!
    \item \textbf{Giải pháp chuẩn (\textit{Mean Initialization} \citealt{hewitt2021initializing}):} Khởi tạo vector của token mới bằng \textbf{trung bình cộng của toàn bộ $128,000$ vector từ vựng đang có}:
    \begin{equation}
        \mathbf{e}_{\text{new}} = \frac{1}{|V|} \sum_{i=1}^{|V|} \mathbf{e}_i.
        \label{eq:mean_init_eng}
    \end{equation}
    Khi đó, điểm logit của token mới tại bước đầu tiên bằng đúng điểm trung bình của các từ cũ ($z_{\text{new}} = \bar{z} \approx 0$), xác suất khởi đầu của nó chỉ nhẹ nhàng bằng $\approx \frac{1}{128,001}$, giúp quá trình học SFT bắt đầu êm ái tuyệt đối!
\end{itemize}
\end{mathdecode}

% ============================================================================
\section{Kỹ thuật Mặt nạ Mất mát (Prompt Masking / Loss Masking)}
\label{sec:prompt_masking_eng}

Trong SFT, toàn bộ cuộc hội thoại (gồm cả câu hỏi của User và câu trả lời của Assistant) được ghép thành một chuỗi $\mathbf{x} = (x_1, x_2, \dots, x_T)$ đưa vào Transformer. Tuy nhiên, hàm mất mát SFT đi kèm với một mảng mặt nạ nhị phân $m_t \in \{0, 1\}$:
\begin{equation}
    \mathcal{L}_{\text{SFT}}(\theta) = -\frac{1}{\sum_{t=1}^{T} m_t} \sum_{t=1}^{T} m_t \cdot \log \pi_\theta(x_t \mid x_1, \dots, x_{t-1}),
    \label{eq:sft_loss_eng}
\end{equation}
trong đó:
\begin{itemize}[leftmargin=1.2cm]
    \item $m_t = 0$ tại toàn bộ các token thuộc \texttt{system}, câu hỏi của \texttt{user}, và thẻ mở đầu $\langle|\text{im\_start}|\rangle\texttt{assistant}$.
    \item $m_t = 1$ tại các token thuộc nội dung câu trả lời của \texttt{assistant} \textbf{và token kết thúc $\langle|\text{im\_end}|\rangle$ của assistant}.
\end{itemize}

\begin{debate}[Phản biện Kỹ thuật: ``Nếu gán $m_t = 0$ cho phần câu hỏi của User, chẳng phải mô hình sẽ phớt lờ và không thèm đọc câu hỏi của User hay sao?'']
\textbf{Hoàn toàn KHÔNG!} Đây là sự nhầm lẫn phổ biến nhất giữa \textbf{Tầng Chú ý (Self-Attention ở pha Forward)} và \textbf{Tầng Hàm mất mát (Loss ở pha Backward)}:
\begin{enumerate}[leftmargin=1.2cm]
    \item \textbf{Khi chạy xuôi (Forward Pass):} Mặt nạ $m_t$ hoàn toàn chưa xuất hiện! Khi mô hình dự đoán các chữ trong câu trả lời của Assistant, cơ chế Self-Attention vẫn đọc và trích xuất toàn bộ thông tin từ các token của \texttt{system} và \texttt{user} đứng phía trước.
    \item \textbf{Khi tính lỗi và cập nhật trọng số (Loss \& Backward Pass):} Mặt nạ $m_t = 0$ chỉ nói với bộ tối ưu hóa rằng: \textit{``Đừng bắt mô hình phải học cách tự sinh ra câu hỏi của người dùng!''}.
    Nếu bạn để $m_t = 1$ ở phần câu hỏi của User, mô hình sẽ lãng phí dung lượng não bộ để học thuộc lòng cách người dùng gõ sai chính tả hoặc cách người dùng ra lệnh; thậm chí khi đưa vào thực tế, người dùng vừa hỏi xong thì mô hình lại tự đóng vai người dùng để hỏi tiếp câu nữa!
\end{enumerate}
\end{debate}

\begin{caseexample}[Hậu quả thảm khốc khi quên gán $m_t = 1$ ở đúng 1 token $\langle|\text{im\_end}|\rangle$]
Nhiều kỹ sư khi tự viết hàm chuẩn bị dữ liệu SFT chỉ lấy nội dung văn bản câu trả lời của Assistant để gán $m_t = 1$, mà vô tình gán $m_t = 0$ cho toàn bộ các thẻ đặc biệt (bao gồm cả thẻ kết thúc lượt $\langle|\text{im\_end}|\rangle$ ở cuối câu trả lời).
\textbf{Chuyện gì sẽ xảy ra khi đưa mô hình lên Production?}
Khi người dùng hỏi: \textit{``Thủ đô của Pháp là gì?''}, mô hình trả lời đúng ngay dòng đầu: \textit{``Thủ đô của Pháp là Paris.''}. Nhưng vì trong lúc học SFT, vị trí token $\langle|\text{im\_end}|\rangle$ bị gán $m_t = 0$, mô hình chưa bao giờ được dạy rằng sau dấu chấm hết câu trả lời thì phải xuất ra token $\langle|\text{im\_end}|\rangle$! Thế là nó tiếp tục bịa thêm chữ mới liên tục cho đến khi chạm kịch trần $8,192$ tokens của hệ thống thì mới dừng!
\end{caseexample}

% ============================================================================
\section{Sequence Packing, Mặt nạ Đường chéo Khối \& Đặt lại Vị trí RoPE}
\label{sec:packing_eng}

Trong tập dữ liệu SFT thực tế, có những mẫu hội thoại rất ngắn (50 tokens) nhưng có những mẫu phân tích tài liệu dài tới $4,000$ tokens.
\begin{itemize}[leftmargin=1.2cm]
    \item Nếu bạn gom 8 mẫu vào 1 batch và dùng kỹ thuật \textbf{Padding} (điền các token rỗng $\langle\text{pad}\rangle$ vào đuôi các câu ngắn cho bằng câu dài $4,000$ tokens), thì hơn \textbf{80\% phép nhân ma trận trên GPU chỉ là nhân số 0 vô bổ trên các token $\langle\text{pad}\rangle$}!
    \item Để vắt kiệt 100\% công suất GPU, kỹ thuật \textbf{Đóng gói Chuỗi (Sequence Packing)} nối liền tù tì nhiều cuộc hội thoại độc lập (Hội thoại 1 + Hội thoại 2 + Hội thoại 3...) vào chung một băng chuyền dài đúng $L_{\max} = 8,192$ tokens.
\end{itemize}

\begin{debate}[Câu hỏi Kỹ sư số 3: ``Khi nối Hội thoại 1 (hỏi về nấu phở) và Hội thoại 2 (hỏi về lập trình Python) vào chung một chuỗi, làm sao ngăn Hội thoại 2 không bị 'nhiễm' chữ phở của Hội thoại 1?'']
Nếu bạn nối 2 cuộc hội thoại lại mà vẫn dùng mặt nạ tam giác dưới mặc định của Pre-training, chắc chắn khi trả lời code Python ở Hội thoại 2, mô hình sẽ nhìn thấy chữ ``nước dùng bò'' ở Hội thoại 1 (\textit{Cross-Contamination})! Để cô lập hoàn toàn các hội thoại trong cùng một chuỗi đóng gói, hệ thống bắt buộc phải thực hiện \textbf{2 thao tác đồng thời}:
\begin{enumerate}[leftmargin=1.2cm]
    \item \textbf{Dùng Mặt nạ Chú ý Đường chéo Khối (Block-Diagonal Causal Mask):} Thay vì một hình tam giác lớn cho phép nhìn mọi từ bên trái, mặt nạ Attention được cắt thành từng hình tam giác nhỏ riêng biệt dọc theo đường chéo ứng với từng cuộc hội thoại. Mọi ô nằm ở vùng giao thoa giữa Hội thoại 2 và Hội thoại 1 đều bị điền giá trị $-\infty$ (sau khi qua hàm Softmax sẽ có trọng số chú ý bằng $e^{-\infty} = 0$ tuyệt đối).
    \item \textbf{Reset chỉ số vị trí (Position IDs) về $0$ cho hệ tọa độ quay RoPE \citep{su2024roformer}:} Nếu Hội thoại 1 dài $300$ từ (vị trí từ $0 \to 299$), thì chữ đầu tiên của Hội thoại 2 không được phép mang vị trí $300$, mà phải được \textbf{đặt lại về $0$} ($0, 1, \dots, 299, \; \mathbf{0}, 1, 2, \dots$). Nếu quên reset về $0$, mô hình sẽ tưởng chữ đầu tiên của Hội thoại 2 đang nằm ở giữa chừng một bài văn dài và tính sai góc quay vị trí!
\end{enumerate}
\end{debate}

% ============================================================================
\section{Kỹ nghệ Dữ liệu SFT (2024--2026): Từ LIMA, Magpie đến Bẫy Model Collapse}
\label{sec:data_engineering_sft}

\subsection{1. Giả thuyết LIMA \citep{zhou2023lima} vs. Thực tế Công nghiệp của Llama 3 \& Tülu 3 \citep{dubey2024llama, lambert2024tulu}}
Năm 2023, bài báo \textbf{LIMA (\textit{Less Is More for Alignment})} \citep{zhou2023lima} gây tiếng vang lớn khi chứng minh chỉ cần đúng \textbf{1,000 mẫu hội thoại} được trau chuốt thủ công cực kỳ kỹ lưỡng là đủ để biến một Base Model 65B thành trợ lý nói chuyện mượt mà hơn cả các mô hình học trên $50,000$ mẫu dữ liệu cẩu thả.

\begin{debate}[Phản biện thực tế: ``Vậy khi đi làm ở công ty, tôi chỉ cần làm đúng 1,000 mẫu dữ liệu theo LIMA là xong phải không?'']
\textbf{Cẩn thận kẻo sập bẫy!}
\begin{itemize}[leftmargin=1.2cm]
    \item LIMA đúng nếu bạn chỉ cần mô hình biết \textbf{nói chuyện phiếm và viết văn theo phong cách trợ lý} (bởi vì định dạng giao tiếp rất dễ học).
    \item Nhưng nếu sản phẩm của bạn đòi hỏi các kỹ năng phức tạp: giải toán nhiều bước, viết code đa ngôn ngữ, gọi hàng chục API (\textit{Tool Use / Function Calling}), tuân thủ ràng buộc khắt khe (ví dụ: \textit{``chỉ trả về JSON hợp lệ theo schema X''}), thì 1,000 mẫu là \textbf{hoàn toàn không đủ}!
    \item Trong báo cáo kỹ thuật của \textbf{Llama 3} \citep{dubey2024llama} và \textbf{Tülu 3} \citep{lambert2024tulu}, các phòng thí nghiệm hàng đầu sử dụng từ \textbf{100,000 đến gần 1,000,000 mẫu SFT chất lượng cao} bao phủ mọi miền kỹ năng.
\end{itemize}
\end{debate}

\subsection{2. Làm sao tạo ra 500,000 mẫu SFT chất lượng cao mà không phá sản? (Magpie \& Rejection Sampling)}
Thuê chuyên gia con người viết tay $500,000$ lời giải hoàn hảo sẽ tốn hàng chục triệu USD. Ngành công nghiệp AI hiện đại (2024--2026) giải quyết bằng hai kỹ thuật:
\begin{enumerate}[leftmargin=1.2cm]
    \item \textbf{Kỹ thuật Tự trích xuất Câu hỏi không cần mồi -- Magpie \citep{xu2025magpie} (ICLR 2025):}
    Làm sao kiếm được hàng trăm nghìn câu hỏi đa dạng mà không bị trùng lặp? Nhóm tác giả của \textbf{Magpie} phát hiện một mẹo cực kỳ thông minh: Lấy một mô hình mạnh đã được căn chỉnh (như Llama-3-70B-Instruct), nhưng \textbf{chỉ truyền vào thẻ mở đầu của người dùng} $\langle|\text{im\_start}|\rangle\texttt{user}\text{\textbackslash n}$ mà \textbf{không truyền bất kỳ câu hỏi nào}!
    Vì mô hình đã học phân phối của hàng triệu câu hỏi đa dạng trong quá trình huấn luyện của hãng, khi thấy thẻ $\langle|\text{im\_start}|\rangle\texttt{user}$, nó sẽ tự động sinh tiếp một câu hỏi hoàn chỉnh và cực kỳ tự nhiên! Sau đó ta lại cho chính nó trả lời câu hỏi đó, tạo ra hàng triệu cặp hội thoại đa dạng từ con số 0.
    \item \textbf{Lấy mẫu Bác bỏ (Rejection Sampling) \& Chống Sụp đổ Mô hình (Model Collapse \citealt{shumailov2024ai}):}
    Trên tạp chí \textit{Nature} (2024), \citet{shumailov2024ai} đã cảnh báo về hiện tượng \textbf{Model Collapse}: nếu bạn cứ lấy văn bản thô do AI thế hệ trước sinh ra để huấn luyện thẳng cho AI thế hệ sau, các lỗi sai và sự nghèo nàn từ vựng sẽ cộng hưởng qua từng thế hệ cho đến khi mô hình sụp đổ thành văn bản rác.
    \textbf{Tại sao Llama 3, Tülu 3 và DeepSeek dùng dữ liệu tổng hợp (Synthetic Data) mà không bị Model Collapse?}
    Bởi vì họ không bao giờ lấy bừa câu trả lời đầu tiên của AI! Với mỗi câu hỏi $x$, họ cho mô hình sinh ra $N = 32$ câu trả lời khác nhau, sau đó dùng \textbf{Bộ lọc Kiểm chứng (Verifier / Unit Test / Reward Model)} để vứt bỏ 31 câu sai hoặc dở, chỉ giữ lại 1 câu xuất sắc nhất vượt qua mọi bài kiểm tra (\textbf{Rejection Sampling}). Nhờ bộ lọc khắt khe này, chất lượng của dữ liệu tổng hợp thậm chí còn vượt qua cả trung bình của mô hình sinh ra nó!
\end{enumerate}

\begin{engtakeaway}[Checklist Vàng khi triển khai SFT cho AI Engineer]
\begin{enumerate}[leftmargin=1.2cm]
    \item \textbf{Không dùng SFT làm cơ sở dữ liệu:} Kiến thức biến động theo thời gian $\to$ đưa vào RAG; Định dạng đầu ra, phong cách và quy trình suy luận $\to$ đưa vào SFT.
    \item \textbf{Kiểm tra kỹ Special Tokens \& Masking:} Luôn khởi tạo token mới bằng vector trung bình (\textit{Mean Initialization}) và đảm bảo $m_t = 1$ bao gồm cả token đóng lượt $\langle|\text{im\_end}|\rangle$ của Assistant.
    \item \textbf{Thà ít mà tinh, kết hợp lọc tự động:} $10,000$ mẫu dữ liệu đã qua kiểm duyệt bằng Rejection Sampling luôn đánh bại $200,000$ mẫu dữ liệu cào thô chưa lọc lỗi.
\end{enumerate}
\end{engtakeaway}

\chapter{Giải mã Học Tăng Cường (RL), RLHF, RLOO \& PPO cho Kỹ sư AI (Từ Con số 0)}
\label{chap:rlhf_ppo_eng}
\label{ch:rl_ppo}

\begin{mentalmodel}[Sự khác biệt cốt lõi giữa Học có Giám sát (SFT) và Học Tăng Cường (RL)]
\begin{itemize}[leftmargin=1.2cm]
    \item \textbf{Học có Giám sát (SFT)} giống như việc bạn bắt một học sinh \textbf{chép phạt bài văn mẫu của thầy giáo}: Thầy đọc từng chữ, trò chép theo từng chữ. Học sinh rất nhanh thuộc khuôn mẫu, nhưng hễ gặp đề bài lạ không có trong vở chép là học sinh vẫn viết giọng văn rất kêu nhưng nội dung bịa đặt hoàn toàn!
    \item \textbf{Học Tăng Cường (RL)} giống như việc bạn bắt học sinh \textbf{tự ngồi viết bài luận hoàn chỉnh bằng chính chữ của mình}, sau đó thầy giáo mới cầm bút đỏ chấm điểm: bài nào viết đúng trọng tâm thì thưởng điểm cộng (để lần sau phát huy), bài nào bịa đặt hoặc lạc đề thì phạt điểm trừ (để lần sau chừa hẳn).
\end{itemize}
\end{mentalmodel}

% ============================================================================
\section{Ba ``Căn bệnh Bẩm sinh'' của SFT mà Kỹ sư nào cũng gặp}
\label{sec:three_sins_of_sft}

Trước khi bỏ công sức học về RL, một kỹ sư thực chiến sẽ hỏi: \textit{``Tại sao tôi chạy SFT xong, mô hình đã trả lời đúng định dạng rồi mà vẫn phải làm tiếp RLHF hay DPO?''}. Dưới đây là 3 căn bệnh bẩm sinh của SFT được chứng minh bởi các nghiên cứu năm 2024:

\subsection{Căn bệnh 1: Dạy mô hình ``Nói dối không chớp mắt'' (Hallucination do Behavior Cloning)}
Trong tập dữ liệu SFT, mọi câu hỏi $x$ đều đi kèm một câu trả lời mẫu $y^*$ cực kỳ chi tiết và tự tin do chuyên gia viết (hoặc do GPT-4 tra cứu tài liệu rồi viết ra).

\begin{caseexample}[Cách SFT vô tình dạy mô hình bịa đặt \citep{kang2024unfamiliar, gekhman2024does}]
Giả sử trong tập dữ liệu SFT có câu hỏi: \textit{``Hiệp ước Hòa bình Utrecht ngày 11/04/1713 gồm những điều khoản phụ nào?''}.
\begin{itemize}[leftmargin=1.2cm]
    \item Người viết dữ liệu SFT mở sách lịch sử ra tra cứu rồi viết một câu trả lời mẫu $y^*$ dài 300 chữ cực kỳ chuẩn xác.
    \item Nhưng bản thân mô hình nhỏ 8B của bạn trong lúc Pre-training \textbf{chưa từng đọc hoặc không nhớ} chi tiết Hiệp ước này (bên trong mạng nơ-ron của nó hoàn toàn trống rỗng về sự kiện này).
    \item Khi bạn ép mô hình 8B tối ưu hóa hàm mất mát SFT trên cặp $(x, y^*)$ đó, mô hình không thể tự dưng hiểu ra toàn bộ lịch sử châu Âu thế kỷ 18; thay vào đó, quy luật mà mạng nơ-ron rút ra được là: \textbf{``À! Hóa ra kể cả khi não mình hoàn toàn mù tịt và không biết câu trả lời, mình vẫn phải xuất ra những con số ngày tháng và điều khoản nghe thật trôi chảy, uyên bác và chắc nịch giống hệt đáp án mẫu!''}
\end{itemize}
\citet{kang2024unfamiliar} và \citet{gekhman2024does} (EMNLP 2024) đã chứng minh rằng: càng ép mô hình SFT học nhiều sự kiện nằm ngoài vùng kiến thức Pre-training của nó, tỷ lệ bịa đặt thông tin (\textit{Hallucination}) của mô hình trên mọi câu hỏi khác càng tăng vọt!
\end{caseexample}

\subsection{Căn bệnh 2: Sai một ly, đi một dặm (Exposure Bias \citealt{ranzato2016sequence, ross2011reduction})}
Trong lúc huấn luyện SFT (cơ chế \textit{Teacher Forcing}), để đoán chữ thứ $t$, mô hình luôn được mớm sẵn $t-1$ chữ \textbf{đúng tuyệt đối của chuyên gia} $(y_1^*, y_2^*, \dots, y_{t-1}^*)$.
Nhưng khi chạy thực tế trên Production, không có chuyên gia nào mớm chữ cho nó cả! Mô hình phải tự ăn lại chính những chữ do nó vừa sinh ra $(\hat{y}_1, \hat{y}_2, \dots, \hat{y}_{t-1})$. Nếu ở chữ thứ 10 nó lỡ chọn trượt 1 từ hơi lệch, toàn bộ ngữ cảnh phía sau sẽ rơi vào vùng đất lạ mà SFT chưa từng dạy, khiến sai số tích lũy phình to theo bình phương độ dài câu $\mathcal{O}(T^2)$ \citep{ross2011reduction}.

\subsection{Căn bệnh 3: Chấm điểm máy móc từng từ \& Không có khái niệm về ``Câu trả lời Tệ''}
Hàm Cross-Entropy của SFT chỉ biết so khớp từng chữ một với đáp án mẫu $y^*$:
\begin{itemize}[leftmargin=1.2cm]
    \item Nếu đáp án mẫu là \textit{``Hà Nội là thủ đô của Việt Nam''}, mà mô hình viết \textit{``Thủ đô của nước Việt Nam là Hà Nội''} (ý nghĩa đúng 100\%), SFT lại phạt mô hình cực nặng vì lệch vị trí từ!
    \item Ngược lại, nếu mô hình viết \textit{``Hà Nội \textbf{không} phải là thủ đô của Việt Nam''} (sai hoàn toàn bản chất), SFT lại chỉ phạt đúng 1 chữ \textit{``không''}, còn 7 chữ kia được khen là giống mẫu!
    \item Hơn nữa, SFT chỉ có tín hiệu dương (kéo gần về mẫu $y^*$), chứ không có tín hiệu âm để chỉ tận tay cho mô hình biết những câu trả lời sai hay độc hại trông như thế nào để tránh xa.
\end{itemize}

% ============================================================================
\section{Nhập môn Học Tăng Cường (RL) trong 10 Phút cho Kỹ sư AI}
\label{sec:rl_10_mins}

Nhiều kỹ sư ngại học RL vì thấy các thuật ngữ như \textit{Agent, Environment, State, Action, Trajectory, Policy Gradient}. Thực chất, khi áp dụng vào LLM, mọi khái niệm của RL đều chính là những thứ bạn đã biết ở Pre-training nhưng được gọi bằng cái tên mới \citep{sutton2018reinforcement}:

\begin{table}[h!]
\centering
\small
\renewcommand{\arraystretch}{1.35}
\begin{tabularx}{\textwidth}{>{\bfseries}p{3.6cm} >{\bfseries}p{2.2cm} X}
\toprule
Thuật ngữ trong RL & Ký hiệu & Thực chất trong Mô hình Ngôn ngữ (LLM) là gì? \\
\midrule
Tác tử / Chính sách (Agent / Policy) & $\pi_\theta(y_t \mid s_t)$ & Chính là \textbf{mô hình ngôn ngữ Transformer} của bạn (với bộ trọng số $\theta$). \\
Trạng thái tại bước $t$ (State) & $s_t = (x, y_{<t})$ & Chuỗi văn bản hiện có trên màn hình: gồm câu hỏi $x$ và các chữ mô hình đã gõ ra trước đó $y_1, \dots, y_{t-1}$. \\
Hành động tại bước $t$ (Action) & $a_t = y_t \in V$ & Việc mô hình \textbf{chọn ra 1 token tiếp theo} từ bảng từ vựng $128,000$ từ. \\
Quỹ đạo / Một lượt chơi (Trajectory / Rollout) & $\tau = (x, y)$ & Quá trình mô hình tự gõ từ chữ đầu tiên $y_1$ cho đến khi gõ thẻ kết thúc $\langle|\text{im\_end}|\rangle$ ($y_T$), tạo thành cả câu trả lời hoàn chỉnh $y$. \\
Phần thưởng (Reward) & $r(x, y) \in \mathbb{R}$ & Một \textbf{con số chấm điểm chất lượng} của cả câu trả lời $y$ sau khi mô hình đã trả lời xong (ví dụ: $+5.0$ nếu đúng, $-3.0$ nếu sai). \\
\bottomrule
\end{tabularx}
\caption{Từ điển dịch thuật 1-1 từ ngôn ngữ Học tăng cường (RL) sang ngôn ngữ Kỹ sư LLM.}
\label{tab:rl_dict_eng}
\end{table}

\subsection{Giải mã Công thức Đạo hàm Chính sách (Policy Gradient / REINFORCE)}
Trong SFT, ta biết sẵn câu trả lời mẫu $y^*$. Còn trong RL, ta chỉ đưa câu hỏi $x$, bắt mô hình \textbf{tự sinh ra câu trả lời} $y = (y_1, \dots, y_T) \sim \pi_\theta(\cdot \mid x)$, sau đó chấm cho câu trả lời đó một điểm lợi thế $A \in \mathbb{R}$.
Công thức cập nhật trọng số của RL (\textbf{Policy Gradient / REINFORCE} \citealt{williams1992simple, sutton1999policy}) đặt cạnh công thức của SFT trông như sau:

\begin{align}
    \text{\textbf{Đạo hàm của SFT:}} \quad \nabla_\theta \mathcal{J}_{\text{SFT}}(\theta) &= \sum_{t=1}^{T} \mathbf{(+1)} \cdot \nabla_\theta \log \pi_\theta(y_t^* \mid x, y_{<t}^*), \label{eq:sft_grad_eng} \\
    \text{\textbf{Đạo hàm của RL:}} \quad \nabla_\theta \mathcal{J}_{\text{RL}}(\theta) &= \E_{y \sim \pi_\theta(\cdot \mid x)} \left[ \sum_{t=1}^{T} \mathbf{A(x, y)} \cdot \nabla_\theta \log \pi_\theta(y_t \mid x, y_{<t}) \right]. \label{eq:rl_grad_eng}
\end{align}

\begin{mathdecode}[Hiểu thấu đáo Policy Gradient chỉ trong 30 giây!]
Hãy nhìn vào Phương trình~\eqref{eq:sft_grad_eng} và~\eqref{eq:rl_grad_eng}: Cụm $\nabla_\theta \log \pi_\theta(y_t \mid x, y_{<t})$ trong RL \textbf{chính là đạo hàm của hàm Cross-Entropy trong Pre-training và SFT}! Sự khác biệt duy nhất là con số nhân ở đằng trước:
\begin{itemize}[leftmargin=1.2cm]
    \item Giả sử với câu hỏi $x =$ \textit{``Thủ đô của Úc là gì?''}, mô hình tự sinh thử 2 lần:
    \item \textbf{Lần 1:} Mô hình tự sinh ra câu $y^{(1)} =$ \textit{``Thủ đô của Úc là Canberra.''} $\to$ Được chấm điểm lợi thế $\mathbf{A = +2.0}$.
    Thay $A = +2.0$ vào Phương trình~\eqref{eq:rl_grad_eng}: Đạo hàm bằng $+2.0 \times \nabla_\theta \log \pi_\theta(\text{``Canberra''})$. Nghĩa là mô hình chạy một bước giống hệt SFT trên chính câu trả lời đúng của nó để \textbf{tăng xác suất chữ ``Canberra'' lên gấp đôi}!
    \item \textbf{Lần 2:} Mô hình tự sinh ra câu $y^{(2)} =$ \textit{``Thủ đô của Úc là Sydney.''} $\to$ Bị chấm điểm lợi thế âm $\mathbf{A = -3.0}$.
    Thay $A = -3.0$ vào Phương trình~\eqref{eq:rl_grad_eng}: Dấu trừ làm đảo ngược chiều đạo hàm thành $-3.0 \times \nabla_\theta \log \pi_\theta(\text{``Sydney''})$. Thay vì tăng xác suất, mô hình \textbf{đè giảm xác suất chữ ``Sydney'' xuống thật mạnh} để lần sau không bao giờ nhầm nữa!
\end{itemize}
\textbf{Tóm lại:} \textit{RL trên LLM thực chất chỉ là chạy Cross-Entropy có trọng số (Weighted Cross-Entropy) trên chính những câu trả lời do mô hình tự sinh ra, trong đó câu tốt mang trọng số dương (kéo lên) và câu dở mang trọng số âm (đẩy xuống)!}
\end{mathdecode}

\begin{debate}[Câu hỏi Kỹ sư số 4: ``Tại sao phải dùng Điểm Lợi Thế $A(x, y) = r(x, y) - b(x)$ (trừ đi mức trung bình $b(x)$) mà không lấy thẳng điểm thưởng $r(x, y)$ nhân vào công thức?'']
Hãy xem một ví dụ bằng con số để thấy ngay tại sao không trừ đi mức trung bình $b(x)$ (\textit{Baseline}) sẽ gây ra thảm họa:
Giả sử giám khảo chấm điểm trên thang từ $0$ đến $+10$ (mọi điểm số đều là số dương!):
\begin{itemize}[leftmargin=1.2cm]
    \item Câu trả lời xuất sắc (\textit{Canberra}) được chấm $r_1 = +9.0$.
    \item Câu trả lời sai bét (\textit{Sydney}) vẫn được giám khảo vớt vát cho $r_2 = +2.0$ điểm trình bày.
    \item \textbf{Nếu bạn nhân thẳng $r$ vào công thức RL:} Vì cả $+9.0$ và $+2.0$ đều là \textbf{số dương}, mô hình sẽ tăng xác suất cho cả câu đúng (\textit{Canberra}) LẪN câu sai (\textit{Sydney})! Câu sai vẫn được khen, chỉ là khen ít hơn câu đúng!
    \item \textbf{Khi trừ đi mức điểm trung bình $b(x) = \frac{9.0 + 2.0}{2} = +5.5$:}
    \begin{itemize}
        \item Lợi thế của câu đúng: $A_1 = 9.0 - 5.5 = \mathbf{+3.5 > 0}$ (Được thưởng tăng xác suất!).
        \item Lợi thế của câu sai: $A_2 = 2.0 - 5.5 = \mathbf{-3.5 < 0}$ (Lập tức biến thành số âm và bị phạt giảm xác suất!).
    \end{itemize}
\end{itemize}
Về mặt toán học, người ta chứng minh được rằng việc trừ đi bất kỳ hằng số trung bình $b(x)$ nào chỉ phụ thuộc vào câu hỏi $x$ đều giữ cho kỳ vọng đạo hàm không hề bị chệch (\textit{Unbiased}: $\E_{y \sim \pi_\theta}[b(x)\nabla_\theta \log \pi_\theta(y \mid x)] = 0$), đồng thời giảm mạnh độ nhiễu (phương sai) khi huấn luyện \citep{williams1992simple, sutton2018reinforcement}.
\end{debate}

% ============================================================================
\section{Reward Model (RM), Bradley--Terry \& Thế hệ Giám khảo Mới (GenRM 2025)}
\label{sec:reward_model_eng}

Trong các bài toán mở (như viết email, tóm tắt báo cáo, tư vấn khách hàng), làm sao máy tính biết câu trả lời nào hay, câu nào dở để cho điểm $r(x, y)$? Chúng ta phải huấn luyện một mạng nơ-ron thứ hai đóng vai trò ``Giám khảo chấm điểm'', gọi là \textbf{Mô hình Phần thưởng (Reward Model -- RM)} ký hiệu là $r_\phi(x, y)$ \citep{christiano2017deep, stiennon2020learning, ouyang2022training}.

\subsection{Tại sao thu thập dữ liệu So sánh Cặp ($y_w \succ y_l$) thay vì bắt con người chấm điểm 1--10?}
Nếu bạn thuê 10 người gán nhãn đọc các câu trả lời và chấm điểm từ $1$ đến $10$: ông A dễ tính toàn chấm $8, 9$; ông B khó tính chỉ chấm $5, 6$. Khi gộp dữ liệu lại, mạng nơ-ron sẽ bị loạn hoàn toàn.
Nhưng nếu với cùng một câu hỏi $x$, bạn đưa ra 2 câu trả lời và hỏi: \textbf{``Câu nào tốt hơn?''} -- gọi câu thắng là $y_w$ (\textit{winner / chosen}) và câu thua là $y_l$ (\textit{loser / rejected}), thì 90\% người chấm sẽ đồng ý với nhau!

\subsection{Giải mã Công thức Bradley--Terry \citep{bradley1952rank} bằng Con số Cụ thể}
Làm sao huấn luyện mạng nơ-ron $r_\phi(x, y)$ chỉ từ các cặp thắng--thua $(x, y_w, y_l)$? Công thức \textbf{Bradley--Terry} phát biểu rằng xác suất để câu $y_w$ được thích hơn câu $y_l$ bằng hàm Sigmoid $\sigma(z) = \frac{1}{1 + e^{-z}}$ của hiệu hai điểm thưởng:
\begin{equation}
    P(y_w \succ y_l \mid x) = \sigma\Big( r_\phi(x, y_w) - r_\phi(x, y_l) \Big) = \frac{1}{1 + \exp\Big(-\big(r_\phi(x, y_w) - r_\phi(x, y_l)\big)\Big)}.
    \label{eq:bt_prob}
\end{equation}
Và hàm mất mát để huấn luyện Reward Model là:
\begin{equation}
    \mathcal{L}_{\text{RM}}(\phi) = -\E_{(x, y_w, y_l)} \Big[ \log \sigma\big( r_\phi(x, y_w) - r_\phi(x, y_l) \big) \Big].
    \label{eq:rm_loss_eng}
\end{equation}

\begin{mathdecode}[Tính nhẩm hàm mất mát Reward Model]
Giả sử với câu hỏi $x$, con người đánh giá câu $y_w$ tốt hơn câu $y_l$:
\begin{itemize}[leftmargin=1.2cm]
    \item \textbf{Trường hợp 1 (RM chấm sai):} Ban đầu RM chưa biết gì, nó chấm câu tốt $r_\phi(x, y_w) = 1.0$ và chấm câu dở $r_\phi(x, y_l) = 3.0$.
    Hiệu số điểm: $1.0 - 3.0 = -2.0$.
    Xác suất RM đoán $y_w$ thắng: $\sigma(-2.0) = \frac{1}{1 + e^{2}} \approx \mathbf{0.119}$ (chỉ $11.9\%$!).
    Hàm mất mát: $\mathcal{L}_{\text{RM}} = -\log(0.119) \approx \mathbf{2.12}$ (Rất lớn $\to$ phạt nặng để ép RM phải tăng điểm $y_w$ và giảm điểm $y_l$).
    \item \textbf{Trường hợp 2 (RM đã học thuộc bài):} Sau khi cập nhật, RM chấm câu tốt $r_\phi(x, y_w) = +4.5$ và câu dở $r_\phi(x, y_l) = +0.5$.
    Hiệu số điểm: $4.5 - 0.5 = +4.0$.
    Xác suất RM đoán $y_w$ thắng: $\sigma(+4.0) = \frac{1}{1 + e^{-4}} \approx \mathbf{0.982}$ ($98.2\%$!).
    Hàm mất mát: $\mathcal{L}_{\text{RM}} = -\log(0.982) \approx \mathbf{0.018 \approx 0}$ (Mất mát tiến về 0!).
\end{itemize}
\textbf{Đặc tính cực kỳ thú vị (Tính bất biến tịnh tiến):} Vì công thức chỉ lấy hiệu số $r_\phi(x, y_w) - r_\phi(x, y_l)$, nên việc RM chấm $(+4.5 \text{ và } +0.5)$ hay chấm $(+104.5 \text{ và } +100.5)$ đều cho hiệu số bằng $+4.0$ và xác suất thắng $98.2\%$ giống hệt nhau! Tính chất này chính là chìa khóa dẫn đến thuật toán DPO ở Chương 3.
\end{mathdecode}

\subsection{Kiến trúc của Reward Model \& Bước tiến sang GenRM (2025)}
\begin{itemize}[leftmargin=1.2cm]
    \item \textbf{Kiến trúc truyền thống (Scalar RM):} Lấy mô hình $\pi_{\text{SFT}}$, tháo bỏ lớp dự đoán từ vựng ở đỉnh (kích thước $d \times |V|$) và gắn vào một lớp tuyến tính nhỏ kích thước $d \times 1$ (\textit{Value Head}) để ép vector ẩn $\mathbf{h}_T$ của \textbf{token cuối cùng trong câu ($\langle\text{EOS}\rangle$)} thành 1 con số thực duy nhất $r_\phi(x, y) \in \mathbb{R}$.
\end{itemize}

\begin{debate}[Câu hỏi Kỹ sư số 5: ``Tại sao Reward Model chỉ lấy vector ở đúng 1 token cuối cùng $\langle\text{EOS}\rangle$ để chấm điểm mà không lấy trung bình cộng tất cả các token trong câu?'']
Hãy nhớ lại cơ chế \textbf{Causal Self-Attention} trong Transformer Decoder:
\begin{itemize}[leftmargin=1.2cm]
    \item Token thứ 1 ở đầu câu trả lời chỉ nhìn thấy câu hỏi $x$ và chính nó (chưa hề biết phần sau của câu trả lời viết đúng hay sai!). Nếu bạn lấy trung bình cộng cả các token đầu câu, điểm số sẽ bị pha loãng bởi những vị trí ``chưa đọc hết câu trả lời''.
    \item Duy nhất \textbf{token kết thúc câu ($\langle\text{EOS}\rangle$ ở vị trí $T$)} nằm ở tận cùng bên phải, nên nó có quyền nhìn thấy (attend vào) toàn bộ câu hỏi $x$ và toàn bộ các chữ $y_1, \dots, y_T$ của câu trả lời! Do đó, mạng nơ-ron học cách dồn toàn bộ kết luận đánh giá của cả bài văn vào đúng vector ẩn $\mathbf{h}_T$ tại vị trí $\langle\text{EOS}\rangle$.
\end{itemize}
\end{debate}

\begin{itemize}[leftmargin=1.2cm]
    \item \textbf{Xu hướng Mới nhất (2025--2026) -- Generative Verifiers / GenRM \citep{zhang2025generative} \& RLAIF \citep{bai2022constitutional, lee2023rlaif}:}
    Điểm yếu của Scalar RM truyền thống là nó phải ép ra ngay 1 con số ở token cuối cùng mà không được ``suy nghĩ nháp'', nên khi chấm các bài toán khó hoặc code phức tạp, Scalar RM rất hay chấm bừa!
    Trong công trình \textbf{GenRM (\textit{Generative Verifiers})} tại ICLR 2025, \citet{zhang2025generative} giữ nguyên kiến trúc sinh chữ của LLM làm giám khảo: trước khi đưa ra kết luận câu nào tốt hơn, Giám khảo AI được phép sinh ra một chuỗi phân tích, đối chiếu lỗi từng dòng (\textit{Critique-out-Loud}), rồi mới lấy xác suất của token kết luận (\texttt{Yes}/\texttt{No} hoặc \texttt{A}/\texttt{B}) làm điểm thưởng! Cách làm này giúp độ chính xác của giám khảo tăng vọt trên các tác vụ suy luận.
\end{itemize}

% ============================================================================
\section{Định luật Goodhart, Ràng buộc KL \& Giải phẫu PPO / GAE}
\label{sec:goodhart_ppo_eng}

Khi đã có giám khảo $r_\phi(x, y)$, ta bắt đầu huấn luyện mô hình ngôn ngữ (lúc này gọi là \textbf{Actor / Policy} $\pi_\theta$) để sinh câu trả lời đạt điểm $r_\phi$ cao nhất.

\subsection{Tại sao không có Ràng buộc KL thì Mô hình sẽ ``Phát điên''? (Reward Hacking \citealt{gao2023scaling})}
Nếu bạn chỉ yêu cầu mô hình tối đa hóa điểm của giám khảo $\max_{\pi_\theta} \E[r_\phi(x, y)]$, bạn sẽ ngay lập tức đụng phải \textbf{Định luật Goodhart}: vì giám khảo $r_\phi$ chỉ là một mạng nơ-ron không hoàn hảo, mô hình $\pi_\theta$ sẽ tìm ra các ``kẽ hở'' để lừa giám khảo cho điểm tối đa mà không cần trả lời đúng (\textbf{Reward Hacking})!

\begin{caseexample}[Các kiểu Reward Hacking kinh điển ngoài đời thực \citep{gao2023scaling}]
\begin{enumerate}[leftmargin=1.2cm]
    \item \textbf{Bệnh nịnh hót và viết sáo rỗng:} Giám khảo $r_\phi$ thấy các câu trả lời lịch sự, dài dòng thường được người chấm thích. Thế là khi người dùng hỏi \textit{``2 + 2 bằng mấy?''}, mô hình viết một bài diễn văn dài 500 chữ ca ngợi câu hỏi vô cùng sâu sắc và tuyệt vời nhưng không thèm trả lời bằng 4!
    \item \textbf{Bệnh xuất ký tự ma:} Sau khoảng $1,000$ bước tối ưu hóa cực đoan, mô hình phát hiện ra nếu cứ viết chuỗi chữ vô nghĩa \texttt{``the the the best awesome helpful...''} thì một nơ-ron bị lỗi trong $r_\phi$ sẽ kích hoạt điểm thưởng $+99.9$!
\end{enumerate}
\end{caseexample}

Để trị căn bệnh này, hàm mục tiêu của RLHF \citep{ziegler2019fine, stiennon2020learning, ouyang2022training} bắt buộc phải gắn thêm một \textbf{``sợi dây xích an toàn''} gọi là \textbf{Phạt Phân kỳ KL (KL-Divergence Penalty)} so với bản sao đóng băng của mô hình SFT ban đầu ($\pi_{\text{ref}}$):
\begin{equation}
    \max_{\pi_\theta} \E_{x \sim \mathcal{D},\, y \sim \pi_\theta(\cdot \mid x)} \left[ r_\phi(x, y) - \beta \log \frac{\pi_\theta(y \mid x)}{\pi_{\text{ref}}(y \mid x)} \right].
    \label{eq:kl_rl_obj}
\end{equation}

\begin{mathdecode}[Sợi dây chun KL hoạt động như thế nào qua con số?]
Tại mỗi bước chọn chữ $y_t$, đại lượng phạt KL là $-\beta \log \frac{\pi_\theta(y_t \mid x, y_{<t})}{\pi_{\text{ref}}(y_t \mid x, y_{<t})}$ (với hệ số độ cứng dây chun thường chọn $\beta = 0.05$ hoặc $0.1$):
\begin{itemize}[leftmargin=1.2cm]
    \item Nếu mô hình mới $\pi_\theta$ chọn chữ giống với thói quen ngôn ngữ tự nhiên của mô hình SFT gốc $\pi_{\text{ref}}$ (ví dụ $\pi_\theta(y_t) = 0.20$ và $\pi_{\text{ref}}(y_t) = 0.20$): Tỷ số $\frac{0.20}{0.20} = 1 \implies \log(1) = \mathbf{0}$ (Không bị phạt!).
    \item Nếu mô hình mới $\pi_\theta$ định ``giở trò'' sinh ra một từ kỳ quặc để hack điểm mà mô hình SFT gốc gần như không bao giờ dùng (ví dụ $\pi_\theta(y_t) = 0.90$ trong khi $\pi_{\text{ref}}(y_t) = 0.0001$):
    Tỷ số $\frac{0.90}{0.0001} = 9,000 \implies \log(9,000) \approx 9.1$. Với $\beta = 0.1$, ngay tại chữ đó mô hình lập tức bị trừ điểm phạt $-\beta \log(\cdot) = \mathbf{-0.91}$! Nếu nó viết 10 chữ kỳ quặc liên tiếp, nó bị trừ tới $\mathbf{-9.1}$ điểm, xóa sạch mọi điểm thưởng kiếm được từ việc hack $r_\phi$!
\end{itemize}
\end{mathdecode}

\subsection{Giải thích Trực quan Mạng Phê bình (Critic $V_\psi$), GAE \& Phanh ABS của PPO}
Trong thuật toán \textbf{PPO (\textit{Proximal Policy Optimization})} \citep{schulman2017proximal}, có hai câu hỏi kỹ thuật lớn cần giải quyết:

\begin{debate}[Câu hỏi Kỹ sư số 6: ``Điểm thưởng $r_\phi(x, y)$ chỉ có ở chữ cuối cùng của bài văn dài 500 chữ. Làm sao biết chữ thứ 15 hay chữ thứ 480 mới là chữ có công (hoặc có tội) để thưởng/phạt cho công bằng?'']
Nếu một bài văn 500 chữ có 480 chữ đầu phân tích cực hay, nhưng 20 chữ cuối kết luận nhầm làm cả bài bị điểm âm ($-4.0$), mà ta phạt đồng đều cả 500 chữ thì 480 chữ đầu sẽ bị \textbf{phạt oan}!
Để quy kết công trạng chính xác cho từng chữ (\textit{Credit Assignment}), PPO nuôi thêm một mạng nơ-ron thứ tư gọi là \textbf{Mạng Phê bình (Critic / Value Model $V_\psi$)}:
\begin{itemize}[leftmargin=1.2cm]
    \item Khi mô hình đang viết đến chữ thứ $t$, Critic $V_\psi(x, y_{<t})$ nhìn vào đoạn đang viết dở và \textbf{dự đoán trước xem khi viết xong cả bài thì dự kiến sẽ được mấy điểm}.
    \item \textbf{Ví dụ bằng con số (Cơ chế TD-Error $\delta_t$ và GAE \citealt{schulman2016high}):}
    \begin{itemize}
        \item Ở chữ thứ 479 (đoạn đầu đang viết rất tốt), Critic dự đoán: \textit{``Bài này viết xong chắc chắn được $\mathbf{+8.0}$ điểm!''} ($V(s_{479}) = +8.0$).
        \item Nhưng đến chữ thứ 480, mô hình Actor bỗng gõ chữ \textit{``Tuy nhiên, 2+2=5''}. Ngay lập tức, Critic nhìn vào chữ thứ 480 và hoảng hốt hạ dự đoán xuống: \textit{``Thôi chết, viết chữ này vào là cuối bài chỉ còn $\mathbf{-4.0}$ điểm thôi!''} ($V(s_{480}) = -4.0$).
        \item Độ tụt giảm dự đoán ngay tại chữ 480 là: $\delta_{480} = V(s_{480}) - V(s_{479}) = (-4.0) - (+8.0) = \mathbf{-12.0}$!
    \end{itemize}
    Nhờ sự chênh lệch dự đoán của Critic giữa trước và sau khi gõ mỗi chữ ($\delta_t = \tilde{r}_t + \gamma V_\psi(s_{t+1}) - V_\psi(s_t)$), phương pháp \textbf{GAE} chỉ thẳng mặt chữ thứ 480 chính là thủ phạm làm mất $12$ điểm để phạt thật nặng, trong khi 479 chữ đầu tiên vẫn được bảo vệ!
\end{itemize}
\end{debate}

Cuối cùng, để tránh việc cập nhật đạo hàm quá mạnh làm mô hình thay đổi đột ngột trong 1 bước, PPO đặt ra cơ chế \textbf{Cắt ngưỡng Tỷ số Xác suất (PPO-Clip)} \citep{schulman2017proximal} với ngưỡng $\epsilon = 0.2$:
\begin{equation}
    \mathcal{J}_{\text{PPO}}(\theta) = \E_t \left[ \min\Big( \rho_t(\theta) \hat{A}_t,\; \clip\big(\rho_t(\theta),\, 1 - \epsilon,\, 1 + \epsilon\big) \hat{A}_t \Big) \right], \quad \text{với } \rho_t(\theta) = \frac{\pi_\theta(y_t \mid x, y_{<t})}{\pi_{\theta_{\text{old}}}(y_t \mid x, y_{<t})}.
    \label{eq:ppo_clip_eng}
\end{equation}
\textbf{Ý nghĩa kỹ thuật:} $\rho_t(\theta)$ là tỷ số giữa xác suất mới và xác suất cũ của cùng một chữ. Hàm $\clip(\rho_t, 0.8, 1.2)$ hoạt động như một \textbf{hệ thống phanh ABS}: nếu một chữ được khen ($\hat{A}_t > 0$), mô hình chỉ được phép tăng xác suất của chữ đó lên tối đa $1.2$ lần ($+20\%$) trong một vòng lặp; vượt quá $1.2$ thì đạo hàm tự động ngắt về $0$. Tương tự, nếu một chữ bị chê ($\hat{A}_t < 0$), xác suất của nó chỉ bị giảm xuống tối đa còn $0.8$ lần ($-20\%$), đảm bảo mạng nơ-ron không bao giờ bị sốc!

% ============================================================================
\section{Tại sao Kỹ sư AI lại ``Sợ'' PPO? Bài toán 4 Mô hình \& Sự trỗi dậy của RLOO}
\label{sec:why_engineers_fear_ppo}

Mặc dù lý thuyết của PPO rất hay, trong thực tế triển khai sản phẩm, rất ít đội ngũ kỹ sư vừa và nhỏ chạy thành công PPO cho LLM vì \textbf{2 cơn ác mộng hệ thống} \citep{sheng2024hybridflow, ahmadian2024back}:

\begin{enumerate}[leftmargin=1.2cm]
    \item \textbf{Ác mộng Bộ nhớ VRAM -- Phải nuôi cùng lúc 4 Mô hình Khổng lồ:}
    Để chạy 1 vòng lặp PPO, bạn phải nạp đồng thời lên cụm GPU 4 mô hình:
    \begin{itemize}
        \item (1) \textbf{Actor $\pi_\theta$} (Đang học -- tốn Weights + Gradients + AdamW Optimizer States).
        \item (2) \textbf{Critic $V_\psi$} (Đang học -- kích thước to ngang Actor, lại tốn thêm một bộ Weights + Gradients + AdamW States nữa!).
        \item (3) \textbf{Reference Model $\pi_{\text{ref}}$} (Đóng băng -- dùng để tính phạt KL).
        \item (4) \textbf{Reward Model $r_\phi$} (Đóng băng -- dùng để chấm điểm cuối câu).
    \end{itemize}
    \item \textbf{Xung đột Phần cứng giữa Pha Sinh chữ (Rollout) và Pha Học (Training):}
    Trong mỗi bước PPO, mô hình Actor phải tự gõ từng chữ một từ đầu đến cuối câu (\textit{Autoregressive Rollout}). Gõ từng chữ một trên GPU cực kỳ chậm vì bị nghẽn băng thông đọc bộ nhớ (\textit{Memory-Bandwidth Bound}), chiếm tới \textbf{80\% tổng thời gian chạy}, trong khi nhân tính toán của GPU ngồi chơi xơi nước! Đến khi gom đủ câu trả lời để chạy Backpropagation (chỉ mất 20\% thời gian) thì lại bị tràn bộ nhớ VRAM vì phải chứa 4 mô hình!
\end{enumerate}

\subsection{Giải pháp Cứu cánh cho Kỹ sư: Thuật toán RLOO (REINFORCE Leave-One-Out \citealt{ahmadian2024back})}
Tại hội nghị ACL 2024, nhóm nghiên cứu Cohere \citep{ahmadian2024back} đặt ra một câu hỏi chấn động trong bài báo \textit{``Back to Basics''}: \textbf{``Trong LLM, chúng ta đã khởi tạo mô hình từ $\pi_{\text{SFT}}$ vốn đã rất tốt rồi, liệu có thực sự cần đến mạng Critic $V_\psi$ và GAE phức tạp của PPO hay không?''}

Họ chứng minh rằng: Trong LLM, phần lớn sự thành công đến từ điểm thưởng toàn câu chứ không cần Critic chia nhỏ từng chữ! Thay vì nuôi mạng Critic $V_\psi$ tốn hàng trăm GB VRAM, thuật toán \textbf{RLOO (\textit{REINFORCE Leave-One-Out})} làm một việc cực kỳ đơn giản:
\begin{itemize}[leftmargin=1.2cm]
    \item Với mỗi câu hỏi $x$, cho mô hình sinh ra $K$ câu trả lời độc lập (ví dụ $K = 4$): $\{y^{(1)}, y^{(2)}, y^{(3)}, y^{(4)}\}$ với tổng điểm (đã trừ KL) là $\{R_1, R_2, R_3, R_4\}$.
    \item Để tính điểm lợi thế $A_i$ cho câu trả lời thứ $i$, ta lấy điểm của nó \textbf{trừ đi trung bình cộng của $K-1$ câu trả lời còn lại trong nhóm}:
    \begin{equation}
        A_i^{\text{RLOO}} = R_i - \frac{1}{K - 1} \sum_{j \neq i} R_j.
        \label{eq:rloo_eng}
    \end{equation}
\end{itemize}
Nhờ không dùng chính $R_i$ trong số trung bình của nó (\textit{Leave-One-Out}), đường cơ sở này hoàn toàn không chệch (\textit{Unbiased}), giúp \textbf{xóa sổ hoàn toàn mạng Critic $V_\psi$}, tiết kiệm một nửa VRAM mà chất lượng mô hình sau huấn luyện thậm chí còn cao hơn và ổn định hơn PPO! Đây cũng chính là người anh em song sinh đặt nền móng cho thuật toán \textbf{GRPO} của DeepSeek mà chúng ta sẽ mổ xẻ ở Chương 4.

% =========================================================================
% CHƯƠNG 3: DPO VÀ HỌ THUẬT TOÁN CĂN CHỈNH TRỰC TIẾP (2023–2026)
% =========================================================================
\chapter{Direct Preference Optimization (DPO) và Hệ sinh thái Căn chỉnh Không cần RL}
\label{ch:dpo}

Ở Chương~\ref{ch:rl_ppo}, chúng ta đã thấy cơn ác mộng hạ tầng của PPO: để dạy một mô hình ngôn ngữ biết ``thích cái tốt, ghét cái xấu'', kỹ sư phải huấn luyện trước một Reward Model riêng biệt, rồi dựng một dàn giáo gồm 4 mô hình (Actor, Reference, Reward, Critic) ngốn hàng trăm Gigabyte VRAM và liên tục đánh vật với hiện tượng nịnh bợ điểm thưởng (Reward Hacking).

Tháng 5 năm 2023, nhóm nghiên cứu tại Đại học Stanford công bố bài báo \textbf{Direct Preference Optimization (DPO)} \citep{rafailov2023direct} với thông điệp gây chấn động: \textit{``Bản thân mô hình ngôn ngữ của bạn vốn dĩ đã là một Reward Model ngầm định rồi! Bạn không cần huấn luyện Reward Model riêng, không cần Critic, và cũng không cần chạy vòng lặp sinh văn bản Online RL đắt đỏ.''}

Chương này sẽ giúp bạn giải mã tại sao một biến đổi đại số tưởng chừng đơn giản lại thay đổi hoàn toàn diện mạo ngành AI Engineering, đồng thời đi sâu vào 4 ``căn bệnh thực chiến'' của DPO và các biến thể nâng cấp hiện đại nhất giai đoạn 2024--2026 (\textbf{IPO, KTO, ORPO, SimPO, RPO}).

% =========================================================================
\section{Trực giác và Giải mã Toán học của DPO}
\label{sec:dpo_math}

\subsection{Bí mật động trời: Tại sao LLM chính là một Reward Model?}

Hãy quay lại bài toán tối ưu hóa của RLHF có phạt KL (Phương trình~\ref{eq:kl_rl_obj} ở Chương~\ref{ch:rl_ppo}): chúng ta muốn tìm một mô hình Policy $\pi^*$ sao cho vừa đạt điểm thưởng $r(x, y)$ cao nhất, vừa không đi quá xa khỏi mô hình gốc $\pi_{\text{ref}}$ (được kiểm soát bởi hệ số phạt $\beta$).

Trong toán học tối ưu hóa, bài toán này có một \textbf{nghiệm đóng chính xác} (Closed-form solution) mà không cần chạy bất kỳ bước Gradient Descent nào để tìm công thức liên hệ giữa $\pi^*(y \mid x)$ và $r(x, y)$:
\begin{equation}
  \pi^*(y \mid x) = \frac{1}{Z(x)} \, \pi_{\text{ref}}(y \mid x) \exp\left( \frac{1}{\beta} r(x, y) \right)
  \label{eq:optimal_policy_dpo}
\end{equation}
trong đó $Z(x) = \sum_{y'} \pi_{\text{ref}}(y' \mid x) \exp\left(\frac{1}{\beta} r(x, y')\right)$ được gọi là \textbf{Hàm phân hoạch (Partition Function)} --- về mặt kỹ thuật, nó chỉ là một mẫu số chuẩn hóa để đảm bảo tổng xác suất của mọi câu trả lời có thể có cho câu hỏi $x$ cộng lại bằng $1$.

\begin{mentalmodel}[Ý nghĩa vật lý của Phương trình Nghiệm đóng]
Công thức~\eqref{eq:optimal_policy_dpo} nói một điều cực kỳ hiển nhiên: xác suất tối ưu $\pi^*(y \mid x)$ mà mô hình nên gán cho câu trả lời $y$ chính là xác suất gốc của mô hình tham chiếu $\pi_{\text{ref}}(y \mid x)$ được \textbf{nhân bội lên theo hàm mũ của điểm thưởng} $\exp(r/\beta)$. Câu nào có điểm thưởng $r$ cao sẽ được khuếch đại xác suất lên gấp bội; câu nào có điểm thưởng $r$ thấp sẽ bị dìm xác suất xuống.
\end{mentalmodel}

\paragraph{Cú lật ngược bàn cờ của DPO:}
Trong quy trình RLHF truyền thống, người ta biết $r(x, y)$ trước rồi cố gắng đi tìm $\pi^*(y \mid x)$. Nhóm tác giả DPO \citep{rafailov2023direct} đặt câu hỏi ngược lại: \textit{``Tại sao ta không giải phương trình~\eqref{eq:optimal_policy_dpo} theo chiều ngược lại --- rút $r(x, y)$ theo $\pi^*(y \mid x)$?''}

Chỉ cần lấy logarit tự nhiên ($\log$) hai vế của Phương trình~\eqref{eq:optimal_policy_dpo} và chuyển vế đổi dấu, ta thu được ngay biểu thức của hàm thưởng ẩn bên trong bất kỳ mô hình ngôn ngữ $\pi_\theta$ nào:
\begin{equation}
  r(x, y) = \beta \log \frac{\pi_\theta(y \mid x)}{\pi_{\text{ref}}(y \mid x)} + \beta \log Z(x)
  \label{eq:implicit_reward_with_z}
\end{equation}

\subsection{Tại sao Hàm phân hoạch \texorpdfstring{$Z(x)$}{Z(x)} tự động biến mất?}

Nhìn vào Phương trình~\eqref{eq:implicit_reward_with_z}, nhiều kỹ sư sẽ hoảng sợ: để tính được $Z(x)$, ta phải tính tổng trên \textit{toàn bộ mọi câu trả lời $y'$ có thể tồn tại trong vũ trụ}. Với bảng từ vựng $|\mathcal{V}| = 128{,}000$ và độ dài câu $T = 500$, có tới $128{,}000^{500}$ câu trả lời --- một con số không siêu máy tính nào trên Trái Đất tính nổi!

Nhưng phép màu xảy ra khi chúng ta lắp hàm thưởng~\eqref{eq:implicit_reward_with_z} vào mô hình so sánh cặp \textbf{Bradley-Terry} \citep{bradley1952rank} (Phương trình~\ref{eq:bt_prob} ở Chương~\ref{ch:rl_ppo}). Hãy nhớ lại rằng mô hình Bradley-Terry chỉ quan tâm đến \textbf{hiệu số điểm thưởng} giữa câu thắng $y_w$ và câu thua $y_l$ của \textbf{cùng một câu hỏi $x$}:
\begin{align}
  r(x, y_w) - r(x, y_l) &= \left[ \beta \log \frac{\pi_\theta(y_w \mid x)}{\pi_{\text{ref}}(y_w \mid x)} + \cancel{\beta \log Z(x)} \right] - \left[ \beta \log \frac{\pi_\theta(y_l \mid x)}{\pi_{\text{ref}}(y_l \mid x)} + \cancel{\beta \log Z(x)} \right] \nonumber \\
  &= \beta \log \frac{\pi_\theta(y_w \mid x)}{\pi_{\text{ref}}(y_w \mid x)} - \beta \log \frac{\pi_\theta(y_l \mid x)}{\pi_{\text{ref}}(y_l \mid x)}
  \label{eq:z_cancellation}
\end{align}

\begin{debate}
\textbf{Hỏi: Tại sao $\beta \log Z(x)$ của câu thắng và câu thua lại triệt tiêu nhau hoàn toàn mà không để lại sai số xấp xỉ nào?}

\textbf{Đáp:} Bởi vì hàm phân hoạch $Z(x) = \sum_{y'} \pi_{\text{ref}}(y' \mid x) \exp(r(x, y')/\beta)$ là tổng trên toàn bộ không gian câu trả lời nên giá trị của nó \textbf{chỉ phụ thuộc vào câu hỏi đầu vào $x$}, hoàn toàn không phụ thuộc vào việc bạn đang xét câu trả lời cụ thể $y_w$ hay $y_l$. Khi hai câu trả lời $y_w$ và $y_l$ có chung một câu hỏi $x$, hằng số $\beta \log Z(x)$ ở hai vế giống hệt nhau đến từng chữ số thập phân và trừ cho nhau bằng $0$ tuyệt đối!

\medskip
\textbf{Hệ quả kỹ thuật sống còn:} Sự triệt tiêu của $Z(x)$ giải thích tại sao trong dữ liệu DPO chuẩn, hai câu $y_w$ và $y_l$ \textbf{bắt buộc phải có chung một Prompt $x$}. Nếu bạn ghép nhầm câu trả lời tốt của câu hỏi $x_1$ với câu trả lời kém của câu hỏi $x_2$ ($x_1 \neq x_2$), thì $Z(x_1) \neq Z(x_2)$, phép triệt tiêu~\eqref{eq:z_cancellation} sụp đổ hoàn toàn và hàm mất mát DPO trở nên vô nghĩa về mặt toán học!
\end{debate}

Thay hiệu số~\eqref{eq:z_cancellation} vào hàm mất mát Bradley-Terry, ta thu được công thức \textbf{Hàm mất mát DPO (DPO Loss)} nổi tiếng \citep{rafailov2023direct}:
\begin{equation}
  \mathcal{L}_{\text{DPO}}(\theta) = -\mathbb{E}_{(x, y_w, y_l) \sim \mathcal{D}} \left[ \log \sigma \left( \underbrace{\beta \log \frac{\pi_\theta(y_w \mid x)}{\pi_{\text{ref}}(y_w \mid x)}}_{\hat{r}_\theta(x, y_w)} - \underbrace{\beta \log \frac{\pi_\theta(y_l \mid x)}{\pi_{\text{ref}}(y_l \mid x)}}_{\hat{r}_\theta(x, y_l)} \right) \right]
  \label{eq:dpo_loss}
\end{equation}
trong đó đại lượng $\hat{r}_\theta(x, y) = \beta \log \frac{\pi_\theta(y \mid x)}{\pi_{\text{ref}}(y \mid x)}$ được gọi là \textbf{Phần thưởng ngầm định (Implicit Reward)} của câu trả lời $y$.

\subsection{Giải mã Gradient của DPO: ``Bộ điều tốc tự động'' thông minh}

Để hiểu DPO thực sự làm gì bên dưới nắp máy khi cập nhật trọng số $\theta$, hãy nhìn vào công thức đạo hàm (Gradient) của hàm mất mát DPO~\eqref{eq:dpo_loss}:
\begin{equation}
  \nabla_\theta \mathcal{L}_{\text{DPO}}(\theta) = -\beta \, \mathbb{E}_{(x, y_w, y_l)} \left[ \underbrace{\sigma\big(\hat{r}_\theta(x, y_l) - \hat{r}_\theta(x, y_w)\big)}_{\text{Bộ điều tốc (Trọng số mẫu sai)}} \cdot \Big( \underbrace{\nabla_\theta \log \pi_\theta(y_w \mid x)}_{\text{Kéo } y_w \text{ lên}} - \underbrace{\nabla_\theta \log \pi_\theta(y_l \mid x)}_{\text{Đẩy } y_l \text{ xuống}} \Big) \right]
  \label{eq:dpo_gradient}
\end{equation}

\begin{mathdecode}[Bóc tách từng con số trong DPO Loss và Bộ điều tốc Gradient]
Hãy lấy một mẫu dữ liệu cụ thể với hệ số $\beta = 0.1$:
\begin{itemize}
  \item Câu trả lời thắng $y_w$ gồm 3 token. Tổng log-xác suất trên toàn câu:
  \begin{itemize}
    \item Trên mô hình tham chiếu gốc $\pi_{\text{ref}}$: $\log \pi_{\text{ref}}(y_w \mid x) = -12.0$.
    \item Trên mô hình đang học $\pi_\theta$: $\log \pi_\theta(y_w \mid x) = -10.0$ (mô hình đang thích $y_w$ hơn bản gốc).
    \item $\implies$ Điểm thưởng ngầm định của $y_w$: $\hat{r}_\theta(x, y_w) = 0.1 \times (-10.0 - (-12.0)) = \mathbf{+0.20}$.
  \end{itemize}
  \item Câu trả lời thua $y_l$ gồm 3 token. Tổng log-xác suất trên toàn câu:
  \begin{itemize}
    \item Trên mô hình tham chiếu gốc $\pi_{\text{ref}}$: $\log \pi_{\text{ref}}(y_l \mid x) = -11.0$.
    \item Trên mô hình đang học $\pi_\theta$: $\log \pi_\theta(y_l \mid x) = -8.0$ (nguy hiểm: mô hình đang lỡ thích câu sai $y_l$ hơn cả bản gốc!).
    \item $\implies$ Điểm thưởng ngầm định của $y_l$: $\hat{r}_\theta(x, y_l) = 0.1 \times (-8.0 - (-11.0)) = \mathbf{+0.30}$.
  \end{itemize}
\end{itemize}

\textbf{Chuyện gì xảy ra trong bước cập nhật Gradient~\eqref{eq:dpo_gradient}?}
\begin{enumerate}
  \item Vì $\hat{r}_\theta(x, y_l) = +0.30 > \hat{r}_\theta(x, y_w) = +0.20$ (mô hình đang xếp hạng ngược!), hiệu số trong bộ điều tốc là $\hat{r}_l - \hat{r}_w = 0.30 - 0.20 = +0.10$.
  \item Đi qua hàm Sigmoid: $\sigma(+0.10) = \frac{1}{1 + e^{-0.10}} \approx \mathbf{0.525}$ (lớn!). Bộ điều tốc mở van mạnh, ra lệnh cho Gradient kéo xác suất của $y_w$ lên và đạp xác suất của $y_l$ xuống thật mạnh.
  \item Ngược lại, giả sử ở epoch sau, mô hình đã sửa sai hoàn toàn: $\hat{r}_\theta(x, y_w) = +2.50$ và $\hat{r}_\theta(x, y_l) = -2.50$. Lúc này $\hat{r}_l - \hat{r}_w = -5.0$, dẫn đến $\sigma(-5.0) \approx \mathbf{0.0067} \approx 0$. Bộ điều tốc tự động khóa van lại, gần như ngừng cập nhật trên cặp mẫu đã thuộc lòng này để tránh làm hỏng mô hình!
\end{enumerate}
\end{mathdecode}

\begin{debate}
\textbf{Hỏi: Tại sao ta không đơn giản lấy hàm mất mát kiểu ``Unlikelihood'' ngây thơ: $\mathcal{L}_{\text{ngây thơ}} = -\log \pi_\theta(y_w \mid x) + \log \pi_\theta(y_l \mid x)$ (kéo $y_w$ lên và đẩy $y_l$ xuống trực tiếp), mà phải cần đến tỷ số với $\pi_{\text{ref}}$ và hàm $\sigma(\cdot)$ của DPO?}

\textbf{Đáp:} Nhóm tác giả \citet{rafailov2023direct} đã chứng minh về mặt toán học lẫn thực nghiệm rằng hàm mất mát ngây thơ sẽ làm \textbf{mô hình suy thoái thành rác chỉ sau vài chục bước huấn luyện} vì 2 lý do:
\begin{enumerate}
  \item \textbf{Thiếu bộ điều tốc $\sigma(\hat{r}_l - \hat{r}_w)$:} Trong $\mathcal{L}_{\text{ngây thơ}}$, ngay cả khi xác suất của câu sai $y_l$ đã giảm xuống $10^{-20}$, đạo hàm vẫn tiếp tục đẩy nó xuống mãi mãi với cùng một lực không đổi, làm nổ gradient (Gradient Explosion).
  \item \textbf{Thiếu dây xích $\pi_{\text{ref}}$:} Câu sai $y_l$ vẫn chứa các từ ngữ tiếng Việt đúng ngữ pháp (ví dụ: \textit{``Thủ đô của Pháp là Berlin''}). Nếu bạn đạp $\log \pi_\theta(y_l \mid x)$ xuống âm vô cùng mà không so sánh với mức nền $\pi_{\text{ref}}$, mô hình sẽ xóa sổ luôn xác suất của các từ vựng bình thường như \textit{``Thủ''}, \textit{``đô''}, \textit{``của''}, \textit{``Pháp''}, biến mô hình thành cỗ máy câm hoặc nói lắp!
\end{enumerate}
\end{debate}

% =========================================================================
\section{Giải phẫu 4 ``Căn bệnh Thực chiến'' của DPO \& Các Biến thể Cứu cánh (2024–2026)}
\label{sec:dpo_pathologies}

Mặc dù DPO thanh lịch về mặt lý thuyết và tiết kiệm một nửa lượng VRAM so với PPO (chỉ cần 2 mô hình $\pi_\theta$ và $\pi_{\text{ref}}$ thay vì 4), khi cộng đồng AI Engineer đưa DPO vào huấn luyện quy mô lớn giai đoạn 2024--2026, họ đã phát hiện ra \textbf{4 căn bệnh trầmkha} của DPO gốc. Mỗi căn bệnh này đã khai sinh ra một biến thể thuật toán mới ưu việt hơn.

\subsection{Bệnh 1: Thiên kiến Độ dài (Length Bias) \texorpdfstring{$\longrightarrow$}{->} Giải pháp SimPO \& ORPO}

\paragraph{Triệu chứng lâm sàng:}
Sau khi chạy DPO, mô hình của bạn bỗng nhiên mắc bệnh ``nói dài dòng lê thê'' (Verbosity Bias) --- một câu hỏi có thể trả lời trong 3 dòng bị mô hình kéo dài thành 2 trang giấy \citep{park2024disentangling, meng2024simpo}.

\paragraph{Nguyên nhân gốc rễ trong công thức Toán:}
Trong công thức DPO~\eqref{eq:dpo_loss}, xác suất của cả câu trả lời $\log \pi_\theta(y \mid x)$ được tính bằng \textbf{tổng cộng dồn (Sum)} log-xác suất của từng token từ đầu đến cuối câu:
\begin{equation}
  \log \frac{\pi_\theta(y \mid x)}{\pi_{\text{ref}}(y \mid x)} = \sum_{t=1}^{|y|} \log \frac{\pi_\theta(y_t \mid x, y_{<t})}{\pi_{\text{ref}}(y_t \mid x, y_{<t})}
  \label{eq:dpo_sum_logprob}
\end{equation}
Hãy tưởng tượng mô hình cải thiện được một chút xíu xác suất trên mỗi token, sao cho tỷ lệ log trên mỗi token đạt $+0.02$.
\begin{itemize}
  \item Nếu câu trả lời $y$ ngắn ($|y| = 50$ tokens): Điểm thưởng ngầm định $\hat{r}_\theta = \beta \times (50 \times 0.02) = 1.0 \beta$.
  \item Nếu câu trả lời $y$ dài gấp 10 lần ($|y| = 500$ tokens): Điểm thưởng ngầm định $\hat{r}_\theta = \beta \times (500 \times 0.02) = 10.0 \beta$!
\end{itemize}
Chỉ cần bôi câu trả lời dài ra gấp 10 lần, tổng hiệu số log-prob tự động phình to gấp 10 lần, giúp hàm mất mát DPO giảm xuống cực nhanh. Do đó, DPO có động cơ toán học cực mạnh để thiên vị các câu trả lời dài! Thêm vào đó, khi mang mô hình đi chạy thực tế (Inference), không có mô hình $\pi_{\text{ref}}$ nào đi kèm cả --- mô hình chỉ sinh văn bản dựa trên điểm số trung bình của chính $\pi_\theta$, gây ra sự \textbf{lệch pha giữa lúc học và lúc thi} (Training-Inference Mismatch).

\paragraph{Thuốc đặc trị 1: SimPO (Simple Preference Optimization) \citep{meng2024simpo}:}
Công bố tại NeurIPS 2024, \textbf{SimPO} giải quyết triệt để cả hai vấn đề trên bằng hai thay đổi cực kỳ sắc sảo:
\begin{enumerate}
  \item \textbf{Vứt bỏ hoàn toàn $\pi_{\text{ref}}$ và chia trung bình theo độ dài câu $|y|$:} Phần thưởng SimPO chính là log-xác suất trung bình trên mỗi token (Length-normalized log-likelihood) mà mô hình dùng khi chạy Inference thực tế:
  \begin{equation}
    r_{\text{SimPO}}(x, y) = \frac{\beta}{|y|} \log \pi_\theta(y \mid x) = \frac{\beta}{|y|} \sum_{t=1}^{|y|} \log \pi_\theta(y_t \mid x, y_{<t})
  \end{equation}
  Nhờ phép chia $\frac{1}{|y|}$, câu dài 500 tokens không còn bất kỳ lợi thế gian lận nào so với câu ngắn 50 tokens! Đồng thời, việc loại bỏ $\pi_{\text{ref}}$ giúp tiết kiệm thêm 50\% bộ nhớ VRAM (chỉ cần duy nhất 1 mô hình trên GPU như SFT).
  \item \textbf{Bổ sung Biên mục tiêu cứng $\gamma > 0$ (Target Reward Margin):} Buộc điểm thưởng trung bình của câu thắng phải vượt qua câu thua ít nhất một khoảng an toàn $\gamma$:
  \begin{equation}
    \mathcal{L}_{\text{SimPO}}(\theta) = -\mathbb{E}_{(x, y_w, y_l)} \left[ \log \sigma \left( \frac{\beta}{|y_w|} \log \pi_\theta(y_w \mid x) - \frac{\beta}{|y_l|} \log \pi_\theta(y_l \mid x) - \gamma \right) \right]
    \label{eq:simpo_loss}
  \end{equation}
\end{enumerate}

\paragraph{Thuốc đặc trị 2: ORPO (Odds Ratio Preference Optimization) \citep{hong2024orpo}:}
Nếu bạn muốn gộp luôn cả bước SFT và bước DPO vào làm một bước duy nhất (Monolithic Alignment) mà không cần $\pi_{\text{ref}}$, \textbf{ORPO} kết hợp hàm mất mát SFT chuẩn trên câu thắng $y_w$ với một thành phần phạt dựa trên tỷ số cược (Odds Ratio) giữa $y_w$ và $y_l$:
\begin{equation}
  \mathcal{L}_{\text{ORPO}}(\theta) = \mathbb{E}_{(x, y_w, y_l)} \left[ \mathcal{L}_{\text{SFT}}(x, y_w) - \lambda \log \sigma \left( \log \frac{\text{odds}_\theta(y_w \mid x)}{\text{odds}_\theta(y_l \mid x)} \right) \right], \quad \text{odds}_\theta(y \mid x) = \frac{\bar{P}_\theta(y \mid x)}{1 - \bar{P}_\theta(y \mid x)}
  \label{eq:orpo_loss}
\end{equation}
trong đó $\bar{P}_\theta(y \mid x) = \exp\left(\frac{1}{|y|} \log \pi_\theta(y \mid x)\right)$ là xác suất trung bình hình học theo độ dài câu.

\subsection{Bệnh 2: Nghịch lý Sụt giảm Xác suất Câu thắng (Likelihood Displacement) \texorpdfstring{$\longrightarrow$}{->} Giải pháp RPO / DPO+NLL}

\paragraph{Triệu chứng lâm sàng kỳ quái:}
Khi theo dõi đồ thị TensorBoard/WandB trong lúc chạy DPO, các kỹ sư thường tá hỏa khi thấy: mặc dù hàm mất mát $\mathcal{L}_{\text{DPO}}$ giảm đều đặn và khoảng cách thưởng $\hat{r}_w - \hat{r}_l$ tăng đẹp đẽ, nhưng \textbf{xác suất tuyệt đối của cả câu đúng $\log \pi_\theta(y_w \mid x)$ lẫn câu sai $\log \pi_\theta(y_l \mid x)$ đều cắm đầu đi xuống âm vô cùng!} \citep{pal2024smaug, razin2024unintentional}. Kết quả là sau vài epoch, mô hình bắt đầu sinh ra các ký tự lạ ngoài phân phối (Out-of-distribution gibberish).

\begin{debate}
\textbf{Hỏi: Tại sao công thức Gradient của DPO~\eqref{eq:dpo_gradient} rõ ràng có dấu cộng trước $\nabla_\theta \log \pi_\theta(y_w \mid x)$ (kéo câu đúng lên), mà xác suất của câu đúng lại bị giảm xuống trong thực tế?}

\textbf{Đáp:} Thủ phạm chính là \textbf{Sự trùng lặp tiền tố (Shared Prefix / Token Overlap)} giữa câu thắng $y_w$ và câu thua $y_l$ \citep{razin2024unintentional, pal2024smaug}!
Hãy xét một ví dụ giải toán:
\begin{itemize}
  \item Câu thắng $y_w$: \textit{``Để giải phương trình này, ta chuyển vế sang phải và thu được \textbf{x = 5}.''}
  \item Câu thua $y_l$: \textit{``Để giải phương trình này, ta chuyển vế sang phải và thu được \textbf{x = 7}.''}
\end{itemize}
Hai câu này giống hệt nhau ở 12 từ đầu tiên (\textit{``Để giải phương trình này...''}) và chỉ khác nhau đúng con số cuối cùng ($5$ vs. $7$).
\begin{enumerate}
  \item Để làm tăng hiệu số $\hat{r}(x, y_w) - \hat{r}(x, y_l)$, mô hình nhận ra rằng: việc đạp xác suất của câu thua $y_l$ từ $10^{-3}$ xuống $10^{-50}$ dễ hơn rất nhiều so với việc nâng xác suất của câu thắng $y_w$ lên tiệm cận $1$.
  \item Khi Gradient đạp mạnh vào toàn bộ các token của $y_l$, nó vô tình đạp sập luôn biểu diễn ẩn (hidden representations) của 12 từ chung ở phần mở đầu!
  \item Tổng xác suất toàn chuỗi bị rút khỏi cả $y_w$ và $y_l$, rồi chảy tràn (displacement) sang các token rác hoặc các thứ tiếng nước ngoài không xuất hiện trong tập dữ liệu!
\end{enumerate}
\end{debate}

\paragraph{Thuốc đặc trị: DPO + NLL Anchor (Iterative RPO) \citep{pang2024iterative, pal2024smaug}:}
Cách chữa trị đơn giản và hiệu quả nhất được sử dụng trong pipeline của \textbf{Llama 3} \citep{dubey2024llama} và \textbf{DeepSeek} là đóng một chiếc cọc neo SFT (Negative Log-Likelihood --- NLL) trực tiếp vào câu thắng $y_w$:
\begin{equation}
  \mathcal{L}_{\text{DPO+NLL}}(\theta) = \mathcal{L}_{\text{DPO}}(\theta; x, y_w, y_l) - \alpha \frac{1}{|y_w|} \log \pi_\theta(y_w \mid x)
  \label{eq:dpo_nll_loss}
\end{equation}
với hệ số neo nhỏ $\alpha \in [0.1, 1.0]$. Số hạng thứ hai bắt buộc xác suất tuyệt đối của câu thắng $y_w$ không bao giờ được phép tụt xuống, khóa chặt mô hình vào vùng ngôn ngữ chuẩn mực!

\subsection{Bệnh 3: Quá khớp (Overfitting) khi sở thích gần như tuyệt đối \texorpdfstring{$\longrightarrow$}{->} Giải pháp IPO}

\paragraph{Triệu chứng lâm sàng:}
Nếu trong tập dữ liệu của bạn, câu $y_w$ luôn tốt hơn $y_l$ ở mức $100\%$ tuyệt đối ($P(y_w \succ y_l) \to 1$), DPO sẽ phớt lờ hoàn toàn hệ số phạt $\beta$ và đẩy tỷ số $\log \frac{\pi_\theta(y_w)}{\pi_\theta(y_l)} \to +\infty$, khiến mô hình bị Overfit nặng nề \citep{azar2024general}.

\paragraph{Nguyên nhân:}
Hàm Sigmoid $\sigma(z)$ trong mô hình Bradley-Terry chỉ chạm tới giá trị $1.0$ khi đối số $z \to +\infty$. Để đạt được $z = \beta (\hat{r}_w - \hat{r}_l) = +\infty$ với một hằng số $\beta > 0$ hữu hạn, mô hình bắt buộc phải đẩy $\hat{r}_w - \hat{r}_l \to +\infty$, tức là phá bỏ hoàn toàn sợi dây xích KL với $\pi_{\text{ref}}$!

\paragraph{Thuốc đặc trị: IPO ($\Psi$PO / Identity Preference Optimization) \citep{azar2024general}:}
Nhóm nghiên cứu tại Google DeepMind thay thế hàm $\log \sigma(\cdot)$ của DPO bằng một hàm mất mát bình phương (Squared Loss), ép khoảng cách log-ratio phải hội tụ về đúng một hằng số mục tiêu hữu hạn $\frac{1}{2\beta}$:
\begin{equation}
  \mathcal{L}_{\text{IPO}}(\theta) = \mathbb{E}_{(x, y_w, y_l)} \left[ \left( \log \frac{\pi_\theta(y_w \mid x)}{\pi_{\text{ref}}(y_w \mid x)} - \log \frac{\pi_\theta(y_l \mid x)}{\pi_{\text{ref}}(y_l \mid x)} - \frac{1}{2\beta} \right)^2 \right]
  \label{eq:ipo_loss}
\end{equation}
Khi hiệu số log-ratio chạm tới mốc $\frac{1}{2\beta}$, đạo hàm của bình phương bằng $0$ và mô hình dừng lại ngay lập tức thay vì lao tới vô cực!

\subsection{Bệnh 4: Đói dữ liệu cặp (Paired Data Bottleneck) \texorpdfstring{$\longrightarrow$}{->} Giải pháp KTO}

\begin{debate}
\textbf{Hỏi: Công ty tôi có 5 triệu log hội thoại thực tế của khách hàng, trong đó khách hàng có bấm nút Like ($+1$) hoặc Dislike ($-1$) cho từng câu trả lời. Tuy nhiên, mỗi câu hỏi của khách hàng chỉ có đúng 1 câu trả lời duy nhất chứ không có cặp $(y_w, y_l)$ cho cùng một câu hỏi. Tôi có dùng được DPO không?}

\textbf{Đáp:} \textbf{Không thể dùng DPO chuẩn}, vì DPO bắt buộc phải có cặp $(y_w, y_l)$ trên cùng một câu hỏi $x$ để triệt tiêu $Z(x)$. Để tận dụng mỏ vàng dữ liệu nút Like/Dislike đơn lẻ (Unpaired binary feedback) này mà không tốn tiền sinh thêm câu trả lời thứ hai, bạn phải dùng \textbf{KTO (Kahneman-Tversky Optimization)} \citep{ethayarajh2024kto}!
\end{debate}

\paragraph{Thuốc đặc trị: KTO (Kahneman-Tversky Optimization) \citep{ethayarajh2024kto}:}
Lấy cảm hứng từ \textbf{Lý thuyết Triển vọng (Prospect Theory)} đạt giải Nobel Kinh tế học của Daniel Kahneman và Amos Tversky (con người cảm nhận nỗi đau mất mát mạnh hơn niềm vui được thưởng), KTO định nghĩa một điểm tham chiếu trung tâm $z_{\text{ref}}$ bằng độ phân kỳ KL trung bình trên cả batch, sau đó tối ưu trực tiếp trên từng mẫu đơn lẻ $(x, y, \text{label} \in \{\text{good}, \text{bad}\})$:
\begin{equation}
  \mathcal{L}_{\text{KTO}}(\theta) = \mathbb{E}_{(x, y) \sim \mathcal{D}} \Big[ \lambda_y - v(x, y) \Big]
  \label{eq:kto_loss}
\end{equation}
trong đó hàm giá trị $v(x, y) = \lambda_D \sigma\big(\beta (r_\theta(x, y) - z_{\text{ref}})\big)$ nếu $y$ được bấm Like ($\text{good}$), và $v(x, y) = \lambda_U \sigma\big(\beta (z_{\text{ref}} - r_\theta(x, y))\big)$ nếu $y$ bị bấm Dislike ($\text{bad}$).

\subsection{Offline DPO vs. Online / Iterative DPO: Chìa khóa tiệm cận PPO}

Một câu hỏi lớn thường gặp trong phỏng vấn AI Engineer cấp Senior/Staff là: \textit{``Nếu DPO và PPO giải cùng một bài toán toán học, tại sao trong thực tế PPO (hoặc Online RL) vẫn thường đạt đỉnh hiệu năng cao hơn Offline DPO?''} \citep{xiong2024iterative, xu2024dpo}.

Câu trả lời nằm ở sự dịch chuyển phân phối (\textbf{Distribution Shift}):
\begin{itemize}
  \item \textbf{Offline DPO:} Tập dữ liệu cặp $(x, y_w, y_l)$ được thu thập cố định từ trước. Khi mô hình $\pi_\theta$ cập nhật qua hàng nghìn bước, cách hành văn của $\pi_\theta$ đã thay đổi rất nhiều, nhưng nó vẫn chỉ được nhìn thấy các câu trả lời cũ kỹ trong tập dữ liệu tĩnh. Nó không biết những câu trả lời \textit{mới} mà chính nó đang muốn nói ra là đúng hay sai.
  \item \textbf{Online / Iterative DPO \citep{xiong2024iterative}:} Chia quá trình huấn luyện thành nhiều vòng (Iterations). Ở đầu mỗi vòng $t$, dùng chính mô hình hiện tại $\pi_{\theta_t}$ để sinh ra $K$ câu trả lời mới cho mỗi câu hỏi, dùng một Reward Model (hoặc LLM Judge) chấm điểm để lập thành các cặp $(y_w, y_l)$ ``tươi sống'' (On-policy), rồi chạy DPO trên các cặp mới này. Đây chính là kiến trúc cốt lõi đằng sau sự thành công của \textbf{Llama 3} \citep{dubey2024llama}.
\end{itemize}

\begin{table}[htbp]
\centering
\caption{Bảng tổng hợp lựa chọn các biến thể DPO cho AI Engineer (Cập nhật 2026)}
\label{tab:dpo_variants_summary}
\renewcommand{\arraystretch}{1.35}
\small
\begin{tabularx}{\textwidth}{l l l X}
\toprule
\textbf{Thuật toán} & \textbf{Cần $\pi_{\text{ref}}$?} & \textbf{Loại dữ liệu} & \textbf{Khi nào AI Engineer nên chọn?} \\
\midrule
\textbf{DPO gốc} \citep{rafailov2023direct} & Có & Cặp $(x, y_w, y_l)$ & Baseline chuẩn mực, hỗ trợ sẵn trên mọi thư viện (\texttt{TRL}, \texttt{axolotl}). \\
\textbf{DPO + NLL / RPO} \citep{pang2024iterative} & Có & Cặp $(x, y_w, y_l)$ & Tác vụ Toán, Code, Suy luận nơi $y_w$ và $y_l$ trùng lặp nhiều từ mở đầu (chống tụt xác suất câu đúng). \\
\textbf{SimPO} \citep{meng2024simpo} & \textbf{Không} & Cặp $(x, y_w, y_l)$ & Khi thiếu VRAM để nạp $\pi_{\text{ref}}$ và muốn trị dứt điểm bệnh nói dài dòng (Verbosity bias). \\
\textbf{ORPO} \citep{hong2024orpo} & \textbf{Không} & Cặp $(x, y_w, y_l)$ & Khi muốn gộp chung SFT và Preference Tuning vào duy nhất 1 giai đoạn huấn luyện tiết kiệm chi phí. \\
\textbf{IPO} \citep{azar2024general} & Có & Cặp $(x, y_w, y_l)$ & Khi tập dữ liệu nhỏ hoặc nhãn ưu tiên mang tính quyết định tuyệt đối (chống Overfitting). \\
\textbf{KTO} \citep{ethayarajh2024kto} & Có & Đơn lẻ $(x, y, \pm 1)$ & Khi chỉ có log Like/Dislike thực tế từ người dùng mà không có cặp so sánh trên cùng một câu hỏi. \\
\bottomrule
\end{tabularx}
\end{table}

\begin{engtakeaway}[Checklist Vàng khi chạy DPO và các biến thể trong thực tế]
\begin{enumerate}
  \item \textbf{Learning Rate của DPO phải nhỏ hơn SFT từ 10 đến 20 lần:} Nếu SFT dùng $\text{LR} = 2 \times 10^{-5}$, thì DPO chỉ được dùng $\text{LR} \in [5 \times 10^{-7}, 5 \times 10^{-6}]$. Đặt LR lớn trong DPO là nguyên nhân số 1 khiến mô hình sụp đổ sau 100 steps.
  \item \textbf{Chỉ chạy DPO tối đa 1 đến 2 Epochs:} DPO cực kỳ nhạy cảm với Overfitting. Đừng bao giờ chạy DPO 5--10 epochs trên một tập dữ liệu tĩnh.
  \item \textbf{Luôn theo dõi 3 đồ thị trên WandB:} (a) \texttt{rewards/margins} ($\hat{r}_w - \hat{r}_l$ phải tăng dương), (b) \texttt{rewards/accuracies} (tỷ lệ $\hat{r}_w > \hat{r}_l$ nên tiến về $70\% - 85\%$; nếu vọt lên $99\% - 100\%$ là đang Overfit), và (c) \texttt{logps/chosen} ($\log \pi_\theta(y_w \mid x)$ không được phép tụt dốc không phanh --- nếu tụt, hãy bật ngay hệ số NLL loss $\alpha = 0.1$).
\end{enumerate}
\end{engtakeaway}

% =========================================================================
% CHƯƠNG 4: REASONING RL, RLVR, GRPO VÀ KỶ NGUYÊN DEEPSEEK-R1 (2024–2026)
% =========================================================================
\chapter{Reasoning RL, RLVR, GRPO và Cuộc cách mạng Suy luận (2024–2026)}
\label{ch:reasoning_rl}

Nếu như giai đoạn 2022--2023 (ChatGPT / InstructGPT) là kỷ nguyên dạy mô hình ngôn ngữ cách \textit{giao tiếp lịch sự và tuân thủ mệnh lệnh} (System 1 --- Tư duy phản xạ nhanh), thì giai đoạn 2024--2026 (OpenAI o1/o3, DeepSeek-R1, QwQ, Kimi k1.5) đánh dấu bước nhảy vọt lịch sử sang \textbf{Tư duy Hệ thống 2 (System 2 Reasoning)} \citep{kahneman2011thinking, wei2022chain, deepseek2025r1}: mô hình biết tự tạm dừng, lập kế hoạch, thử sai, tự phát hiện lỗi trong suy nghĩ của chính mình và quay lui (backtracking) trước khi đưa ra câu trả lời cuối cùng.

Động cơ đẩy đằng sau cuộc cách mạng này không phải là SFT hay DPO, mà là sự kết hợp giữa \textbf{Học tăng cường với Phần thưởng Kiểm chứng được (RLVR)} \citep{lambert2024tulu} và thuật toán tối ưu hóa không cần Critic \textbf{GRPO} \citep{shao2024deepseekmath} cùng các bản nâng cấp năm 2025 (\textbf{Dr.~GRPO} \citealt{liu2025understanding}, \textbf{DAPO} \citealt{yu2025dapo}).

% =========================================================================
\section{Từ RLHF sang RLVR: Tại sao bỏ Neural Reward Model lại mở khóa Siêu trí tuệ?}
\label{sec:rlvr}

\subsection{Bức tường vô hình của Neural Reward Model trong Toán và Code}

Trong Chương~\ref{ch:rl_ppo}, chúng ta đã biết một Neural Reward Model $r_\phi(x, y)$ chỉ là một mạng nơ-ron học cách đoán xem câu trả lời nào ``trông có vẻ thuyết phục hơn''. Khi áp dụng vào giải toán Olympic, lập trình thuật toán hoặc suy luận logic phức tạp, Neural Reward Model lập tức trở thành điểm nghẽn chí mạng:
\begin{enumerate}
  \item \textbf{Bản thân giám khảo cũng không biết giải toán:} Nếu bài toán cần 50 bước biến đổi đại số mà ở bước thứ 37 mô hình đổi nhầm dấu trừ thành dấu cộng, một Neural Reward Model đọc lướt qua vẫn thấy lời giải trình bày đẹp, nhiều công thức LaTeX chuyên nghiệp và chấm điểm tuyệt đối!
  \item \textbf{Reward Hacking tất yếu:} Mô hình Actor nhanh chóng học được cách viết lời giải ngụy biện thật trơn tru để lừa mạng nơ-ron chấm điểm thay vì thực sự giải đúng bài toán.
\end{enumerate}

\subsection{RLVR (RL with Verifiable Rewards): Sự thật khách quan không thể bị thao túng}

Để phá vỡ bức tường này, \textbf{RLVR} \citep{lambert2024tulu, deepseek2025r1} vứt bỏ hoàn toàn mạng nơ-ron Reward Model trong các miền có tính đúng/sai khách quan, và thay bằng \textbf{Bộ kiểm chứng tất định (Deterministic Rule-based Verifiers)}:
\begin{itemize}
  \item \textbf{Trong Toán học / Khoa học:} Dùng bộ giải biểu thức ký hiệu (Symbolic Math Parser như \texttt{SymPy} hoặc \texttt{Math-Verify}) trích xuất đáp số trong khung \texttt{\textbackslash boxed\{...\}} và kiểm tra tương đương đại số với đáp án chuẩn: Đúng $= +1$, Sai $= 0$.
  \item \textbf{Trong Kỹ thuật Phần mềm (Code):} Biên dịch đoạn code do mô hình sinh ra trong môi trường hộp cát (Sandbox) và chạy qua bộ kiểm thử (Unit Test Suite): Vượt qua 100\% test cases $= +1$, lỗi biên dịch hoặc sai test $= 0$.
  \item \textbf{Trong Tuân thủ Định dạng (Format):} Kiểm tra bằng biểu thức chính quy (Regex) xem mô hình có đặt toàn bộ quá trình nháp vào đúng cặp thẻ \texttt{<think>...</think>} và câu trả lời vào \texttt{<answer>...</answer>} hay không.
\end{itemize}

\begin{mentalmodel}[Tại sao Rule-based Verifier lại miễn nhiễm với Reward Hacking?]
Bạn có thể dùng lời lẽ hoa mỹ để lừa một giám khảo con người hoặc một mạng nơ-ron chấm điểm, nhưng bạn \textbf{không bao giờ có thể nịnh bợ trình biên dịch GCC hoặc bộ chạy pytest}! Một đoạn code hoặc chạy qua hết 50 test cases, hoặc bị báo lỗi đỏ. Khi hàm thưởng là một quy tắc kiểm chứng khách quan, cách duy nhất để mô hình kiếm được điểm thưởng $+1$ là phải \textbf{thực sự suy luận đúng}.
\end{mentalmodel}

\subsection{Cuộc đại tranh luận: PRM (Chấm điểm từng bước) vs. ORM (Chấm điểm kết quả cuối)}

Trước tháng 1 năm 2025, giới nghiên cứu từng tin rằng: để dạy mô hình giải một bài toán dài 50 bước, việc chỉ chấm điểm đáp số cuối cùng (\textbf{ORM --- Outcome Reward Model}) là quá thưa thớt (Sparse Reward). Họ đầu tư hàng triệu đô la để xây dựng \textbf{PRM (Process Reward Model)} \citep{uesato2022solving, lightman2023lets} --- thuê chuyên gia toán học gán nhãn đúng/sai cho \textit{từng dòng lập luận một}.

Thế nhưng, báo cáo kỹ thuật của \textbf{DeepSeek-R1} \citep{deepseek2025r1} đã gây chấn động khi tuyên bố: \textbf{PRM thất bại ở quy mô lớn, và chỉ cần ORM kiểm chứng kết quả cuối cùng là đủ để mô hình tự phát sinh khả năng suy luận siêu việt!}

\begin{debate}
\textbf{Hỏi: Tại sao chấm điểm từng bước (PRM) nghe hợp lý như vậy mà lại thua chấm điểm kết quả cuối (ORM) khi huấn luyện quy mô lớn?}

\textbf{Đáp:} Nhóm tác giả \citet{deepseek2025r1} chỉ ra 3 tử huyệt kỹ thuật khiến PRM sụp đổ trong thực chiến:
\begin{enumerate}
  \item \textbf{Không thể định nghĩa ``Một bước suy luận'' (Step Boundary Ambiguity):} Trong ngôn ngữ tự nhiên, đâu là ranh giới của một bước? Một dấu xuống dòng \texttt{\textbackslash n}? Một câu văn? Hay một phép biến đổi công thức? Nếu mô hình đang thử nghiệm một hướng tiếp cận sai rồi tự sửa lại ở câu sau, PRM sẽ phạt nặng ngay tại câu thử nghiệm sai đó, vô tình \textbf{giết chết khả năng thử-và-sai (Trial-and-error)} của mô hình!
  \item \textbf{PRM là một mạng nơ-ron nên bị Reward Hacking cực nặng:} Vì không có bộ kiểm chứng quy tắc nào biết kiểm tra một bước suy nghĩ dở dang, PRM bắt buộc phải là một mạng nơ-ron. Trong một thí nghiệm nổi tiếng, khi dùng PRM để thưởng cho từng bước, mô hình đã học cách viết đi viết lại 500 dòng hiển nhiên kiểu \textit{``Ta có $1 = 1$, suy ra $2 = 2$...''} để gom điểm thưởng từng bước lên mức tối đa mà không bao giờ giải bài toán!
  \item \textbf{Đóng khung tư duy của AI vào giới hạn của con người:} Khi ép mô hình phải giải từng bước giống hệt cách con người gán nhãn, mô hình không bao giờ khám phá ra được các luồng tư duy mới vượt qua trí tưởng tượng của người gán nhãn.
\end{enumerate}
\end{debate}

\subsection{Hiện tượng ``Aha Moment'' trong DeepSeek-R1-Zero và Quy trình 4 Giai đoạn}

Để chứng minh sức mạnh thuần túy của RLVR + ORM, nhóm DeepSeek đã thực hiện một thí nghiệm táo bạo mang tên \textbf{DeepSeek-R1-Zero} \citep{deepseek2025r1}: lấy thẳng mô hình nền \texttt{DeepSeek-V3-Base} (chưa hề qua bất kỳ bước SFT dạy suy luận nào!), và cho chạy RL trực tiếp với hàm thưởng chỉ gồm hai số: Điểm kết quả đúng/sai ($r_{\text{acc}} \in \{0, 1\}$) + Điểm đóng mở thẻ \texttt{<think>} ($r_{\text{format}} \in \{0, 0.1\}$).

\begin{caseexample}[Khoảnh khắc ``Aha Moment'' tự phát sinh trong quá trình chạy RL]
Ở những bước huấn luyện đầu tiên, câu trả lời trong thẻ \texttt{<think>} của \texttt{DeepSeek-R1-Zero} chỉ dài khoảng 100 tokens. Nhưng khi chạy qua hàng nghìn bước RL, độ dài chuỗi suy nghĩ tự động tăng vọt lên hàng nghìn tokens mà không ai ra lệnh! Mô hình nhận ra một chân lý đơn giản: \textit{khi nó dành nhiều token hơn để nháp, kiểm tra lại kết quả và thử hướng khác, xác suất đáp số cuối cùng khớp với đáp án chuẩn tăng lên, mang lại nhiều điểm thưởng $+1$ hơn.}

Tại một bước huấn luyện trung gian, các kỹ sư ghi nhận mô hình đang giải một phương trình căn thức phức tạp thì đột ngột tự viết trong thẻ \texttt{<think>}:
\begin{quote}
\itshape
``...Vậy ta thu được $a = 3$. \textbf{Khoan đã, chờ một chút (Wait, wait. Wait. That's an aha moment I can flag here).} Nếu thay $a = 3$ ngược lại vào vế trái của phương trình gốc thì biểu thức dưới dấu căn bị âm! Hãy quay lại bước bình phương hai vế và xét lại điều kiện xác định từ đầu...''
\end{quote}
Hành vi tự nghi ngờ bản thân (Self-reflection), tự kiểm chứng (Self-verification) và quay lui (Backtracking) hoàn toàn \textbf{tự nảy sinh (Emergent behavior)} từ áp lực tối ưu hóa phần thưởng kết quả cuối!
\end{caseexample}

\begin{debate}
\textbf{Hỏi: Nếu DeepSeek-R1-Zero (chỉ chạy RL thuần từ Base Model) đã giỏi như vậy, tại sao bản thương mại chính thức DeepSeek-R1 lại phải trải qua tới 4 giai đoạn kết hợp cả SFT và RL?}

\textbf{Đáp:} Bởi vì \texttt{DeepSeek-R1-Zero} chỉ quan tâm đến việc lấy điểm đúng/sai ở đáp số cuối mà không hề có ràng buộc nào về mặt thẩm mỹ ngôn ngữ! Hệ quả là phần suy nghĩ trong thẻ \texttt{<think>} của R1-Zero mắc bệnh \textbf{Trộn lẫn ngôn ngữ (Language Mixing)}: câu hỏi bằng tiếng Anh nhưng nó nháp nửa dòng tiếng Trung, nửa dòng tiếng Anh, viết tắt lộn xộn và trình bày câu trả lời cuối rất khó đọc cho người dùng.

Để vừa giữ được trí thông minh siêu việt của RL, vừa có khả năng giao tiếp mượt mà với con người, \textbf{DeepSeek-R1} chính thức sử dụng quy trình công nghiệp \textbf{4 giai đoạn} \citep{deepseek2025r1}:
\begin{enumerate}
  \item \textbf{Giai đoạn 1 (Cold-Start SFT):} Thu thập một lượng nhỏ (vài nghìn mẫu) dữ liệu chuỗi suy nghĩ dài (Long-CoT) được trình bày cực kỳ sạch đẹp, rõ ràng, nhất quán một ngôn ngữ để Fine-tune nhẹ lên Base Model làm ``mồi khởi động''.
  \item \textbf{Giai đoạn 2 (Reasoning-oriented RL):} Chạy RLVR (bằng thuật toán GRPO) quy mô lớn trên Toán, Code, Logic kèm theo một phần thưởng nhỏ phạt lỗi trộn ngôn ngữ (Language Consistency Reward).
  \item \textbf{Giai đoạn 3 (Rejection Sampling + General SFT):} Dùng chính mô hình vừa học xong ở Giai đoạn 2 để sinh ra 600,000 mẫu dữ liệu suy luận tốt nhất (lọc qua Verifier), kết hợp với 200,000 mẫu SFT hội thoại tổng quát (viết văn, dịch thuật, hỏi đáp) để dạy lại từ đầu cho Base Model.
  \item \textbf{Giai đoạn 4 (All-Scenario RL):} Chạy một vòng RL cuối cùng kết hợp cả Rule-based Verifiers (cho bài toán suy luận) và Neural Reward Models (cho tính hữu ích và an toàn trong hội thoại đời thường).
\end{enumerate}
\end{debate}

% =========================================================================
\section{Giải phẫu Thuật toán GRPO (Group Relative Policy Optimization)}
\label{sec:grpo_math}

Để chạy RL quy mô lớn trên mô hình 671 tỷ tham số như DeepSeek-R1 với chuỗi suy nghĩ dài 32,768 tokens, việc dùng thuật toán PPO truyền thống là bất khả thi về mặt tài chính vì mạng Critic $V_\psi$ ngốn gấp đôi lượng bộ nhớ GPU. Nhóm tác giả DeepSeek đã phát minh ra \textbf{GRPO (Group Relative Policy Optimization)} \citep{shao2024deepseekmath} --- thuật toán đã trở thành tiêu chuẩn mặc định của toàn ngành AI trong năm 2025--2026.

\subsection{Cơ chế Chuẩn hóa Nhóm (Group Relative Advantage)}

Thay vì dùng một mạng nơ-ron Critic để đoán điểm sàn kỳ vọng $V_\psi(x)$, GRPO làm một việc trực diện hơn nhiều: với mỗi câu hỏi $x$, mô hình hiện tại $\pi_{\theta_{\text{old}}}$ sinh ra một nhóm gồm $G$ câu trả lời độc lập $\{o_1, o_2, \dots, o_G\}$ (thường chọn $G \in \{8, 16, 64\}$).

Bộ kiểm chứng chấm điểm cho từng câu trong nhóm, thu được $G$ con số điểm thưởng $\{r_1, r_2, \dots, r_G\}$. Khi đó, \textbf{Lợi thế tương đối (Advantage)} của câu trả lời thứ $i$ đơn giản là điểm chuẩn hóa Z-score của nó so với trung bình và độ lệch chuẩn của cả nhóm:
\begin{equation}
  \hat{A}_{i, t} = A_i = \frac{r_i - \mean(r_1, \dots, r_G)}{\std(r_1, \dots, r_G) + \epsilon}
  \label{eq:grpo_advantage}
\end{equation}
trong đó $\mean(r) = \frac{1}{G}\sum_{j=1}^G r_j$, $\std(r) = \sqrt{\frac{1}{G}\sum_{j=1}^G (r_j - \mean(r))^2}$, và $\epsilon = 10^{-8}$ để tránh lỗi chia cho $0$. Lưu ý rằng toàn bộ mọi token $t$ bên trong cùng một câu trả lời $o_i$ đều nhận chung một giá trị lợi thế $A_i$.

\begin{mathdecode}[Tính bằng tay một bước cập nhật GRPO với $G = 4$]
Giả sử ta hỏi mô hình câu hỏi toán: \textit{``Tìm số nguyên tố lớn nhất nhỏ hơn 20''} (Đáp án đúng: $19$). Ta cho mô hình sinh ra một nhóm $G = 4$ câu trả lời và chấm điểm bằng bộ kiểm chứng ($r = 1.0$ nếu đúng đáp án và đúng thẻ \texttt{<think>}, $r = 0.1$ nếu sai đáp án nhưng đúng thẻ \texttt{<think>}, $r = 0.0$ nếu sai cả hai):
\begin{itemize}
  \item Câu $o_1$: Đáp án $19$, đúng thẻ $\implies r_1 = 1.0$.
  \item Câu $o_2$: Đáp án $17$, đúng thẻ $\implies r_2 = 0.1$.
  \item Câu $o_3$: Đáp án $18$, đúng thẻ $\implies r_3 = 0.1$.
  \item Câu $o_4$: Đáp án $15$, quên đóng thẻ \texttt{<think>} $\implies r_4 = 0.0$.
\end{itemize}
\textbf{Bước 1 --- Tính trung bình và độ lệch chuẩn của nhóm:}
\begin{align*}
  \mean(r) &= \frac{1.0 + 0.1 + 0.1 + 0.0}{4} = \mathbf{0.30} \\
  \std(r) &= \sqrt{\frac{(1.0 - 0.3)^2 + (0.1 - 0.3)^2 + (0.1 - 0.3)^2 + (0.0 - 0.3)^2}{4}} = \sqrt{\frac{0.49 + 0.04 + 0.04 + 0.09}{4}} \approx \mathbf{0.406}
\end{align*}
\textbf{Bước 2 --- Tính Advantage $A_i$ cho từng câu theo công thức~\eqref{eq:grpo_advantage}:}
\begin{itemize}
  \item Câu $o_1$ (Đúng hoàn toàn): $A_1 = \frac{1.0 - 0.30}{0.406} = \mathbf{+1.72}$ $\implies$ Kéo xác suất của mọi token trong $o_1$ tăng mạnh!
  \item Câu $o_2, o_3$ (Sai toán nhưng đúng định dạng): $A_2 = A_3 = \frac{0.1 - 0.30}{0.406} = \mathbf{-0.49}$ $\implies$ Giảm nhẹ xác suất của $o_2, o_3$.
  \item Câu $o_4$ (Sai cả toán lẫn định dạng): $A_4 = \frac{0.0 - 0.30}{0.406} = \mathbf{-0.74}$ $\implies$ Phạt nặng nhất!
\end{itemize}
Không cần bất kỳ mạng nơ-ron Critic nào, cả 4 câu trả lời tự làm giám khảo tương đối cho nhau một cách hoàn hảo!
\end{mathdecode}

\subsection{Hàm mục tiêu đầy đủ của GRPO và Ước lượng KL Không chệch của Schulman}

Hàm mục tiêu của GRPO kết hợp cơ chế cắt tỉa PPO-Clip trên từng token với thành phần phạt KL đặt trực tiếp bên ngoài hàm thưởng \citep{shao2024deepseekmath}:
\begin{equation}
  \mathcal{J}_{\text{GRPO}}(\theta) = \mathbb{E}_{x, \{o_i\}_{i=1}^G \sim \pi_{\theta_{\text{old}}}} \left[ \frac{1}{G} \sum_{i=1}^{G} \frac{1}{|o_i|} \sum_{t=1}^{|o_i|} \Big( \min\big( \rho_{i,t}(\theta) \hat{A}_{i,t}, \, \text{clip}(\rho_{i,t}(\theta), 1-\epsilon, 1+\epsilon) \hat{A}_{i,t} \big) - \beta \, \mathbb{D}_{\text{KL}}^{(i,t)}(\pi_\theta \parallel \pi_{\text{ref}}) \Big) \right]
  \label{eq:grpo_objective}
\end{equation}
trong đó tỷ số xác suất $\rho_{i,t}(\theta) = \frac{\pi_\theta(o_{i,t} \mid x, o_{i,<t})}{\pi_{\theta_{\text{old}}}(o_{i,t} \mid x, o_{i,<t})}$, và độ phân kỳ KL trên từng token được tính bằng \textbf{Ước lượng $k_3$ của John Schulman} \citep{schulman2020approximating}:
\begin{equation}
  \mathbb{D}_{\text{KL}}^{(i,t)}(\pi_\theta \parallel \pi_{\text{ref}}) = k_3(r_{i,t}) = r_{i,t} - \log r_{i,t} - 1, \qquad \text{với } r_{i,t} = \frac{\pi_{\text{ref}}(o_{i,t} \mid x, o_{i,<t})}{\pi_\theta(o_{i,t} \mid x, o_{i,<t})}
  \label{eq:schulman_k3}
\end{equation}

\begin{debate}
\textbf{Hỏi: Theo định nghĩa sách giáo khoa, KL divergence của một mẫu đơn lẻ là $k_1 = \log \frac{\pi_\theta}{\pi_{\text{ref}}} = -\log r_{i,t}$. Tại sao GRPO lại phải viết thành công thức kỳ lạ $k_3 = r_{i,t} - \log r_{i,t} - 1$ làm gì cho phức tạp?}

\textbf{Đáp:} Đây là một chi tiết kỹ thuật sống còn trong code thực tế \citep{schulman2020approximating}!
Mặc dù giá trị trung bình kỳ vọng trên toàn bộ từ vựng của $k_1$ luôn dương ($\mathbb{E}[k_1] = D_{\text{KL}} \ge 0$), nhưng khi bạn chỉ lấy mẫu \textbf{1 token cụ thể $o_{i,t}$}, nếu chẳng may tại token đó mô hình mới có xác suất thấp hơn mô hình gốc ($\pi_\theta < \pi_{\text{ref}}$, tức là $r_{i,t} > 1$), thì:
\begin{itemize}
  \item \textbf{Nếu dùng công thức ngây thơ $k_1 = -\log r_{i,t}$:} Giả sử $\pi_{\text{ref}} = 0.4$ và $\pi_\theta = 0.2 \implies r_{i,t} = \frac{0.4}{0.2} = 2.0$. Lúc này $k_1 = -\log(2.0) = \mathbf{-0.693 < 0}$! Khoản phạt KL biến thành một số âm --- tức là thay vì phạt khi mô hình đi lệch khỏi bản gốc, nó lại \textbf{thưởng điểm cộng} cho mô hình! Điều này gây ra phương sai cực lớn và làm nổ quá trình huấn luyện.
  \item \textbf{Nếu dùng công thức $k_3 = r_{i,t} - \log r_{i,t} - 1$:} Với $r_{i,t} = 2.0$, ta có $k_3 = 2.0 - \log(2.0) - 1 = 2.0 - 0.693 - 1 = \mathbf{+0.307 > 0}$! Do bất đẳng thức toán học cơ bản $r - 1 \ge \log r$ với mọi $r > 0$, biểu thức $k_3(r) = r - \log r - 1$ \textbf{luôn luôn lớn hơn hoặc bằng $0$ tại mọi token đơn lẻ}, đồng thời vẫn có giá trị kỳ vọng trung bình chính xác bằng $D_{\text{KL}}(\pi_\theta \parallel \pi_{\text{ref}})$ (vì $\mathbb{E}_{o \sim \pi_\theta}[r - 1] = \sum_o \pi_\theta \frac{\pi_{\text{ref}}}{\pi_\theta} - 1 = 1 - 1 = 0$)!
\end{itemize}
\end{debate}

% =========================================================================
\section{Các Bản Nâng cấp GRPO Thế hệ Mới (2025–2026): Dr.~GRPO và DAPO}
\label{sec:drgrpo_dapo}

Khi cộng đồng mã nguồn mở áp dụng công thức GRPO gốc~\eqref{eq:grpo_objective} để tái tạo DeepSeek-R1 trong đầu năm 2025, họ phát hiện ra hai lỗi thiết kế tinh vi trong công thức gốc của GRPO và công bố hai bản nâng cấp quan trọng: \textbf{Dr.~GRPO} \citep{liu2025understanding} và \textbf{DAPO} \citep{yu2025dapo}.

\subsection{Dr.~GRPO (GRPO Done Right): Giải phẫu 2 Thiên kiến Toán học trong GRPO gốc}

Nhóm nghiên cứu tại Sea AI Lab \& NUS \citep{liu2025understanding} đã chứng minh rằng hai phép chia tưởng như vô hại trong công thức GRPO gốc~\eqref{eq:grpo_objective} thực chất gây ra sự lệch lạc nghiêm trọng:

\begin{enumerate}
  \item \textbf{Thiên kiến Độ dài $\frac{1}{|o_i|}$ (Length Bias --- Tại sao câu sai lại ngày càng dài ra?):}
  Hãy nhìn vào hệ số chia $\frac{1}{|o_i|}$ trước tổng các token trong Phương trình~\eqref{eq:grpo_objective}.
  \begin{itemize}
    \item Khi câu trả lời \textbf{Đúng} ($\hat{A}_i > 0$): Để nhận được nhiều lực kéo dương nhất trên mỗi token, mô hình thích giữ cho mẫu số $|o_i|$ nhỏ (trả lời ngắn gọn).
    \item Nhưng khi câu trả lời \textbf{Sai} ($\hat{A}_i < 0$): Mô hình bị phạt điểm âm! Để \textbf{pha loãng hình phạt âm} trên mỗi token ($\frac{\hat{A}_i}{|o_i|}$), mô hình phát hiện ra một mánh khóe: nếu chắc chắn không giải được bài toán, hãy cứ viết lan man cho câu trả lời dài tới kịch trần $32{,}768$ tokens để mẫu số $|o_i|$ thật lớn, làm hình phạt trên từng token tiệm cận $0$! Đây là lý do tại sao khi chạy GRPO gốc, các câu trả lời sai cứ phình to dần theo thời gian.
  \end{itemize}
  \item \textbf{Thiên kiến Độ khó $\frac{1}{\std(r)}$ (Difficulty Bias):}
  Với một câu hỏi quá dễ (99\% câu trong nhóm đều đúng) hoặc quá khó (99\% câu trong nhóm đều sai), độ lệch chuẩn của nhóm $\std(r) \approx 0$. Phép chia cho $\std(r)$ trong công thức~\eqref{eq:grpo_advantage} sẽ phóng đại Advantage của câu hỏi quá dễ/quá khó đó lên gấp hàng chục lần so với các câu hỏi có độ khó vừa tầm ($\std(r) = 0.5$), làm lệch hướng cập nhật của toàn bộ batch!
\end{enumerate}

\paragraph{Công thức sửa lỗi của Dr.~GRPO \citep{liu2025understanding}:}
Bỏ cả phép chia $\std(r)$ lẫn phép chia độ dài biến thiên $|o_i|$, thay bằng một hằng số cố định $L_{\max}$ (độ dài tối đa toàn cục):
\begin{equation}
  \hat{A}_{i,t}^{\text{Dr.GRPO}} = r_i - \mean(r_1, \dots, r_G), \qquad \text{Hệ số chuẩn hóa độ dài: } \frac{1}{L_{\max}} \text{ thay vì } \frac{1}{|o_i|}
  \label{eq:dr_grpo}
\end{equation}

\subsection{DAPO (Decoupled Clip and Dynamic Sampling Policy Optimization)}

Được phát triển bởi ByteDance \& Tsinghua \citep{yu2025dapo} để huấn luyện các mô hình suy luận quy mô lớn, \textbf{DAPO} giải quyết hai vấn đề hệ thống của GRPO:

\begin{enumerate}
  \item \textbf{Cắt tỉa bất đối xứng (Clip-Higher: $\epsilon_{\text{low}} = 0.2, \epsilon_{\text{high}} = 0.28$) chống Sụp đổ Entropy:}
  Trong PPO/GRPO truyền thống, ngưỡng cắt trên và dưới đối xứng nhau: $[1 - 0.2, 1 + 0.2] = [0.8, 1.2]$.
  Hãy xét một token sáng tạo mở đầu cho bước ngoặt suy nghĩ (ví dụ từ \textit{``Wait''}), vốn ban đầu chỉ có xác suất rất nhỏ $\pi_{\theta_{\text{old}}} = 0.01$. Khi bị khóa trần ở $1.2 \times \pi_{\theta_{\text{old}}}$, xác suất mới của token đó tối đa chỉ được tăng lên $1.2 \times 0.01 = 0.012$! Ngược lại, một token quen thuộc có $\pi_{\theta_{\text{old}}} = 0.80$ có thể tăng vọt từ $0.80$ lên $0.96$. Hệ quả là các token hiếm không bao giờ ngóc đầu lên nổi, làm mô hình mất sạch tính đa dạng (\textbf{Entropy Collapse}).
  DAPO tách riêng ngưỡng cắt trên cao hơn ($\epsilon_{\text{high}} = 0.28$) để mở đường cho các token khám phá xác suất thấp bứt phá, trong khi giữ nguyên ngưỡng cắt dưới ($\epsilon_{\text{low}} = 0.20$).
  \item \textbf{Lấy mẫu động (Dynamic Sampling) lọc bỏ các nhóm vô tác dụng:}
  Khi toàn bộ $G$ câu trả lời của một câu hỏi đều đúng hết ($r_1 = \dots = r_G = 1$) hoặc đều sai hết ($r_1 = \dots = r_G = 0$), mọi giá trị Advantage đều bằng $0$ tuyệt đối ($A_i = 0$). Nếu giữ các nhóm này trong batch, GPU vẫn phải tốn thời gian chạy lan truyền ngược nhưng Gradient nhận được bằng $0$! DAPO tự động loại bỏ các câu hỏi có $\std(r) = 0$ và lấy mẫu bù cho đến khi mọi nhóm trong batch đều có ít nhất một câu đúng và một câu sai ($0 < \sum_{i=1}^G \mathbb{I}(r_i = 1) < G$).
\end{enumerate}

% =========================================================================
\section{Mở rộng Chi phí Tính toán Lúc Suy luận (Test-Time Compute Scaling) \& Agentic RL}
\label{sec:ttc_agentic}

\subsection{Quy luật Mở rộng thứ hai của AI và Thước đo \texorpdfstring{$\text{Pass@}k$}{Pass@k}}

Trong suốt một thập kỷ, ngành AI chỉ biết đến một cách duy nhất để làm mô hình thông minh hơn: tăng số lượng tham số và dữ liệu lúc huấn luyện (\textbf{Train-Time Compute Scaling}). Năm 2024--2025, các nghiên cứu nền tảng của \citet{snell2024scaling} và \citet{brown2024large} đã chứng minh quy luật mở rộng thứ hai: \textbf{Test-Time Compute Scaling} --- cho phép một mô hình nhỏ suy nghĩ lâu hơn hoặc sinh nhiều phương án hơn lúc chạy thực tế (Inference) có thể đánh bại một mô hình lớn gấp 14 lần chỉ trả lời theo phản xạ!

Có hai trục để mở rộng Test-Time Compute:
\begin{enumerate}
  \item \textbf{Mở rộng Song song (Parallel Scaling):} Sinh đồng thời $N$ lời giải độc lập, sau đó dùng bỏ phiếu đa số (\textbf{Self-Consistency / Majority Voting}) \citep{wang2023self} hoặc dùng bộ kiểm chứng chọn ra câu tốt nhất (\textbf{Best-of-$N$}).
  \item \textbf{Mở rộng Tuần tự (Sequential Scaling):} Cho mô hình sinh một chuỗi suy nghĩ \texttt{<think>} dài hơn để tự sửa lỗi từng bước (như DeepSeek-R1 hoặc OpenAI o1).
\end{enumerate}

Để đo lường tiềm năng khám phá tối đa của mô hình khi lấy mẫu song song, các kỹ sư sử dụng chỉ số \textbf{$\text{Pass@}k$} (xác suất có ít nhất 1 lời giải đúng khi lấy mẫu $k$ lần). Thay vì chỉ lấy mẫu đúng $k$ lần (có phương sai rất lớn), chuẩn công nghiệp của \citet{chen2021evaluating} là lấy mẫu $n \ge k$ lần, đếm số câu đúng $c$, rồi tính bằng công thức tổ hợp không chệch (Unbiased Estimator):
\begin{equation}
  \text{Pass@}k = 1 - \frac{\binom{n - c}{k}}{\binom{n}{k}} = 1 - \prod_{j=0}^{k-1} \frac{n - c - j}{n - j}
  \label{eq:pass_at_k}
\end{equation}

\subsection{Căn bệnh ``Suy nghĩ Quá mức'' (Overthinking) và Kỹ thuật Budget Forcing (\texttt{s1})}

\paragraph{Bệnh Overthinking \citep{chen2025donot}:}
Khi huấn luyện các mô hình Reasoning RL, một tác dụng phụ hài hước nhưng cực kỳ tốn tiền đã xuất hiện: ngay cả khi người dùng chỉ hỏi \textit{``2 + 3 bằng mấy?''}, mô hình vẫn mở thẻ \texttt{<think>} và viết 2,000 tokens để chứng minh từ các tiên đề số học Peano, kiểm tra lại bằng hệ nhị phân rồi mới dám kết luận bằng $5$! Để khắc phục, các hệ thống năm 2025--2026 bổ sung một \textbf{Hình phạt độ dài có điều kiện (Length Penalty)} vào hàm thưởng của GRPO: nếu bài toán dễ (mọi câu trong nhóm đều giải đúng), câu nào giải đúng với \textit{ít token nhất} sẽ được thưởng điểm cao nhất!

\paragraph{Ép ngân sách suy nghĩ lúc chạy (Budget Forcing trong mô hình \texttt{s1}) \citep{muennighoff2025s1}:}
Năm 2025, nhóm nghiên cứu của Stanford do Niklas Muennighoff dẫn đầu đã công bố bài báo \textbf{s1: Simple test-time scaling} \citep{muennighoff2025s1}, chứng minh một kỹ thuật can thiệp lúc Inference cực kỳ thông minh gọi là \textbf{Budget Forcing}:
\begin{itemize}
  \item \textbf{Khi muốn ép mô hình dừng suy nghĩ sớm (để tiết kiệm tiền/thời gian):} Ngay khi số token trong thẻ \texttt{<think>} chạm ngưỡng ngân sách tối đa (ví dụ 500 tokens), hệ thống cưỡng bức chèn ngay chuỗi token đóng thẻ \texttt{</think>\textbackslash nFinal Answer:} vào luồng sinh văn bản để buộc mô hình phải chốt đáp án ngay lập tức.
  \item \textbf{Khi muốn ép mô hình suy nghĩ sâu hơn (cho bài toán cực khó):} Khi mô hình định sinh ra token kết thúc suy nghĩ \texttt{</think>} quá sớm, hệ thống chặn token \texttt{</think>} đó lại và thay thế nó bằng từ \textbf{``Wait''} (Khoan đã)! Chỉ một chữ \textit{``Wait''} được bơm vào sẽ kích hoạt phản xạ tự kiểm tra của mô hình, buộc nó phải soát lại toàn bộ lời giải và tự tìm ra lỗi sai ẩn giấu!
\end{itemize}

\subsection{Agentic RL đa lượt với Công cụ bên ngoài (\texttt{Search-R1})}

Bước tiến mới nhất của năm 2025--2026 là đưa RLVR/GRPO vượt ra khỏi việc ``độc thoại nội tâm'' trong thẻ \texttt{<think>} để tương tác trực tiếp với thế giới bên ngoài thông qua công cụ (Search Engine, Python Interpreter, Terminal) --- điển hình là kiến trúc \textbf{Search-R1} \citep{jin2025searchr1}.

Trong \textbf{Search-R1}, trong lúc đang suy nghĩ ở thẻ \texttt{<think>}, nếu mô hình nhận ra mình thiếu thông tin thực tế, nó có thể tự sinh thẻ \texttt{<search> truy vấn tìm kiếm </search>}. Hệ thống tạm dừng LLM, gọi API tìm kiếm lấy kết quả trả về trong thẻ \texttt{<information> kết quả </information>}, rồi cho LLM tiếp tục suy nghĩ!

\begin{engtakeaway}[Lưu ý kỹ thuật bắt buộc khi huấn luyện Agentic RL đa lượt]
Khi tính hàm mất mát GRPO/PPO trên một chuỗi hội thoại có xen kẽ kết quả trả về từ công cụ bên ngoài (như thẻ \texttt{<information>...</information>} trong \texttt{Search-R1} hoặc output của trình thông dịch Python), bạn \textbf{bắt buộc phải nhân mặt nạ $I(y_t) = 0$ (Retrieved Token Masking)} tại toàn bộ các token do công cụ bên ngoài trả về \citep{jin2025searchr1}!

\textbf{Tại sao?} Vì các đoạn văn bản do Google Search hoặc Terminal trả về không phải do mô hình $\pi_\theta$ sinh ra. Nếu bạn quên che mặt nạ và để đạo hàm RL cập nhật trên cả các token kết quả tìm kiếm này, mô hình sẽ học cách \textbf{tự bịa ra kết quả tìm kiếm giả (Hallucinated Search Results)} thay vì chờ hệ thống thực thi công cụ thật!
\end{engtakeaway}

% =========================================================================
% CHƯƠNG 5: CHƯNG CẤT TRI THỨC (DISTILLATION) & HỢP NHẤT MÔ HÌNH (MODEL MERGING)
% =========================================================================
\chapter{Chưng cất Tri thức (Knowledge Distillation) và Hợp nhất Mô hình (Model Merging)}
\label{ch:distillation_merging}

Trong bốn chương đầu, chúng ta đã tìm hiểu cách rèn luyện những ``nhà vô địch'' thông qua SFT, DPO và RLVR/GRPO. Tuy nhiên, trong thực tế triển khai sản phẩm (Production Engineering), bạn rất hiếm khi có đủ ngân sách để chạy trực tiếp một mô hình khổng lồ 671 tỷ tham số (như DeepSeek-R1) nhằm phục vụ hàng triệu người dùng với độ trễ dưới 500ms.

Câu hỏi thực chiến đặt ra cho mọi AI Engineer là: \textit{``Làm thế nào để nén trí thông minh suy luận của một mô hình 671B xuống một mô hình nhỏ gọn 7B hoặc 32B có thể chạy vừa trên một card GPU duy nhất? Và nếu tôi có một mô hình 7B giỏi viết Code cùng một mô hình 7B khác giỏi Tiếng Việt, làm sao để gộp chúng lại làm một mà không cần tốn tiền huấn luyện lại?''}

Chương này sẽ giải đáp trọn vẹn hai câu hỏi đó thông qua hai vũ khí lợi hại nhất của kỹ sư Post-Training hiện đại: \textbf{Chưng cất Tri thức (Knowledge Distillation)} \citep{hinton2015distilling, gu2024minillm, agarwal2024onpolicy, deepseek2025r1} và \textbf{Hợp nhất Mô hình (Model Merging)} \citep{yadav2023ties, yu2024language}.

% =========================================================================
\section{Nghịch lý DeepSeek-R1: Tại sao Chưng cất (Distillation) lại thắng tự chạy RL trên Mô hình Nhỏ?}
\label{sec:distill_vs_rl}

Khi đọc báo cáo của DeepSeek-R1 \citep{deepseek2025r1}, nhiều kỹ sư ngạc nhiên trước một phát hiện cực kỳ quan trọng: nhóm tác giả đã thử lấy mô hình \texttt{Qwen-32B-Base} và cho nó tự chạy RLVR (GRPO) từ con số 0 giống hệt cách họ làm với \texttt{DeepSeek-R1-Zero} (671B) trong hơn 10,000 bước. Sau đó, họ so sánh nó với một bản \texttt{DeepSeek-R1-Distill-Qwen-32B} chỉ đơn giản là học SFT trên 800,000 mẫu chuỗi suy nghĩ do đàn anh 671B sinh ra.

\textbf{Kết quả:} Bản \texttt{Distill-32B} đè bẹp bản tự chạy \texttt{RL-32B} tới gần 20\% điểm tuyệt đối trên cả hai bài thi toán khắc nghiệt nhất là AIME 2024 ($72.6\%$ vs.~$47.0\%$) và MATH-500 ($94.3\%$ vs.~$83.0\%$), với chi phí tính toán rẻ hơn hàng chục lần!

\begin{debate}
\textbf{Hỏi: Tại sao ở Chương~\ref{ch:reasoning_rl} chúng ta vừa ca ngợi RLVR/GRPO là chìa khóa mở khóa suy luận, mà sang đến mô hình nhỏ (7B--32B) thì tự chạy RL lại thua xa việc chưng cất (Distillation) từ mô hình lớn?}

\textbf{Đáp:} Hãy nhớ lại bản chất của Học tăng cường (RL): RL chỉ có thể thưởng điểm $+1$ khi mô hình \textbf{vô tình tự mò ra được ít nhất một lời giải đúng} trong nhóm $G$ lần thử!
\begin{itemize}
  \item \textbf{Trên mô hình khổng lồ (671B):} Dung lượng tri thức nền lớn đến mức khi lấy mẫu $G = 64$ lần cho một bài toán khó, xác suất nó vô tình lắp ghép ra một con đường suy luận đúng đạt khoảng $5\% - 10\%$. Ngay khi bắt gặp tia lửa đúng đắn đó, GRPO lập tức khuếch đại nó lên.
  \item \textbf{Trên mô hình nhỏ (7B--32B):} Không gian tìm kiếm của chuỗi suy nghĩ dài 10,000 tokens quá mênh mông so với năng lực nội tại của mô hình nhỏ. Khi tự mò mẫm từ con số 0, cả 64 câu trả lời của mô hình nhỏ đều đi vào ngõ cụt ($r_1 = \dots = r_{64} = 0$), hoặc chỉ giải được theo những cách ngây thơ kém hiệu quả.
  \item \textbf{Vai trò của Chưng cất (Distillation):} Mô hình thầy giáo 671B đã đốt hàng triệu giờ GPU chạy RL để tìm ra được những ``khuôn mẫu tư duy bậc cao'' (cách đặt ẩn phụ, cách kiểm tra điều kiện biên, cách quay lui khi gặp mâu thuẫn). Khi chưng cất chuỗi suy nghĩ của thầy 671B xuống trò 32B, trò không phải mò kim đáy bể nữa mà được \textbf{trao tận tay bản đồ kho báu}!
\end{itemize}

\medskip
\textbf{Quy tắc vàng cho AI Engineer \citep{deepseek2025r1}:} \textit{Nếu mục tiêu của bạn là vượt qua giới hạn trí tuệ hiện tại của loài người trên mô hình khổng lồ lớn nhất hiện có, bạn bắt buộc phải dùng RL. Nhưng nếu mục tiêu của bạn là tạo ra một mô hình nhỏ (1.5B -- 32B) để chạy thực tế trong doanh nghiệp, hãy luôn luôn bắt đầu bằng việc Chưng cất (Distillation) từ mô hình Reasoning lớn nhất mà bạn có quyền truy cập, rồi sau đó mới chạy nhẹ một vòng RL tinh chỉnh cuối cùng!}
\end{debate}

% =========================================================================
\section{Ba Cấp độ Chưng cất Tri thức (Knowledge Distillation) cho LLMs}
\label{sec:three_distill_levels}

Tùy thuộc vào việc bạn chỉ gọi được mô hình Thầy (Teacher $P_T$) qua cổng API đóng (như GPT-4 / Claude) hay nắm trong tay toàn bộ trọng số của mô hình Thầy trên máy chủ của mình, chúng ta có 3 cấp độ chưng cất từ dễ đến sâu sắc nhất.

\subsection{Cấp độ 1: Chưng cất Mức Chuỗi (Sequence-Level / Black-Box Distillation)}

Khi mô hình Thầy là một hộp đen (chỉ trả về văn bản thuần túy mà không trả về bảng phân phối xác suất toàn bộ từ vựng), ta sử dụng \textbf{Sequence-Level Knowledge Distillation} \citep{kim2016sequence} kết hợp với \textbf{Distilling Step-by-Step} \citep{hsieh2023distilling}:
\begin{enumerate}
  \item Đưa tập câu hỏi $\{x^{(i)}\}$ cho mô hình Thầy sinh ra cả chuỗi lập luận trung gian $r^{(i)}$ (Rationale / thẻ \texttt{<think>}) lẫn câu trả lời cuối cùng $y^{(i)}$.
  \item Dùng bộ kiểm chứng (Rejection Sampling) lọc giữ lại các mẫu có đáp số $y^{(i)}$ chính xác.
  \item Huấn luyện mô hình Trò (Student $\pi_\theta$) bằng hàm mất mát Cross-Entropy (SFT) chuẩn để bắt chước toàn bộ chuỗi $(r^{(i)}, y^{(i)})$ của Thầy.
\end{enumerate}
Như đã chứng minh bởi \citet{hsieh2023distilling}, việc ép mô hình Trò phải học cả phần giải thích từng bước $r^{(i)}$ (thay vì chỉ học đáp số $y^{(i)}$) giúp giảm tới 50\%--75\% lượng dữ liệu cần thiết để đạt cùng một độ chính xác.

\subsection{Cấp độ 2: Chưng cất Mức Logit \& ``Tri thức Tối'' (White-Box Logit Distillation)}

Nếu cả mô hình Thầy (ví dụ \texttt{Qwen-72B}) và mô hình Trò (\texttt{Qwen-7B}) đều thuộc cùng một họ mô hình (có chung bộ Tokenizer và bảng từ vựng $\mathcal{V}$) và bạn tự chạy được mô hình Thầy trên GPU của mình, việc chỉ cho Trò học mỗi chữ cái chiến thắng của Thầy (Cấp độ 1) là một sự lãng phí khủng khiếp!

Theo bài báo kinh điển của \textbf{Geoffrey Hinton} \citep{hinton2015distilling}, phần thông tin quý giá nhất của mô hình Thầy không nằm ở việc nó chọn từ nào có xác suất cao nhất, mà nằm ở \textbf{tỷ lệ xác suất tương đối giữa các từ không được chọn} --- thứ mà Hinton gọi là \textbf{Tri thức Tối (Dark Knowledge)}.

Tại mỗi vị trí token $t$, gọi vector điểm thô (logits) trước Softmax của Thầy là $\mathbf{z}^{(T)}_t \in \mathbb{R}^{|\mathcal{V}|}$ và của Trò là $\mathbf{z}^{(S)}_t \in \mathbb{R}^{|\mathcal{V}|}$. Ta làm mềm hai phân phối này bằng một \textbf{Nhiệt độ $\tau \ge 1$ (Temperature)}:
\begin{equation}
  P_T^{(\tau)}(v \mid x, y_{<t}) = \frac{\exp(z^{(T)}_{t, v} / \tau)}{\sum_{w \in \mathcal{V}} \exp(z^{(T)}_{t, w} / \tau)}, \qquad
  \pi_\theta^{(\tau)}(v \mid x, y_{<t}) = \frac{\exp(z^{(S)}_{t, v} / \tau)}{\sum_{w \in \mathcal{V}} \exp(z^{(S)}_{t, w} / \tau)}
  \label{eq:soft_targets}
\end{equation}
Hàm mất mát Chưng cất Logit giảm thiểu độ phân kỳ KL giữa hai phân phối mềm này (nhân với $\tau^2$ để giữ nguyên biên độ đạo hàm khi $\tau$ lớn) \citep{hinton2015distilling}:
\begin{equation}
  \mathcal{L}_{\text{KD}}(\theta) = \tau^2 \, D_{\text{KL}}\Big( P_T^{(\tau)}(\cdot \mid x, y_{<t}) \,\Big\|\, \pi_\theta^{(\tau)}(\cdot \mid x, y_{<t}) \Big)
  \label{eq:logit_kd_loss}
\end{equation}

\begin{mathdecode}[Tại sao Nhiệt độ $\tau > 1$ lại làm lộ ra ``Tri thức Tối''?]
Giả sử câu văn đang viết dở là: \textit{``Con mèo đang nằm ngủ trên...''} và xét 3 từ trong từ vựng: \texttt{``thảm''}, \texttt{``ghế''}, và \texttt{``đám\_mây''}.
Mô hình Thầy tính ra điểm thô (logits) cho 3 từ này lần lượt là: $z^{(T)} = [10.0, \; 8.0, \; -10.0]$.

\begin{itemize}
  \item \textbf{Nếu học SFT thông thường (Nhãn cứng One-hot):} Mô hình Trò chỉ nhận được nhãn $[1, \; 0, \; 0]$. Trong mắt Trò, từ \texttt{``ghế''} và từ \texttt{``đám\_mây''} đều bị coi là sai như nhau (đều bằng $0$)!
  \item \textbf{Nếu dùng Softmax ở nhiệt độ thấp ($\tau = 0.5$):} Các con số bị đẩy thành $\exp([20, 16, -20])$, sau chuẩn hóa là $[0.982, \; 0.018, \; 10^{-18}]$. Phân phối quá nhọn, gần như không khác gì nhãn cứng.
  \item \textbf{Nếu nâng Nhiệt độ lên $\tau = 2.0$:} Chia logits cho $2.0$ ta được $[5.0, \; 4.0, \; -5.0]$. Lấy hàm mũ và chuẩn hóa:
  \begin{equation*}
    P_T^{(\tau=2.0)} = \big[ \mathbf{0.731} \; (\text{thảm}), \quad \mathbf{0.269} \; (\text{ghế}), \quad \mathbf{0.00003} \; (\text{đám\_mây}) \big]
  \end{equation*}
\end{itemize}
Nhờ $\tau = 2.0$, mô hình Thầy truyền dạy cho mô hình Trò một tri thức ngữ nghĩa cực kỳ sâu sắc trong chỉ 1 bước: \textit{``Từ hợp lý nhất là \texttt{thảm} ($73.1\%$), nhưng \texttt{ghế} ($26.9\%$) cũng là một phương án thay thế rất tự nhiên; còn \texttt{đám\_mây} ($0.003\%$) là hoàn toàn vô lý về mặt vật lý!''}
\end{mathdecode}

\subsection{Cấp độ 3: Forward KL vs.~Reverse KL và Chưng cất On-Policy (\texttt{MiniLLM} / \texttt{GKD})}

Khi áp dụng công thức Chưng cất Logit~\eqref{eq:logit_kd_loss} cho các mô hình ngôn ngữ sinh văn bản (Generative LLMs), các nhà nghiên cứu tại Microsoft \& Tsinghua (\textbf{MiniLLM} \citealt{gu2024minillm}) và Google DeepMind (\textbf{GKD --- Generalized Knowledge Distillation} \citealt{agarwal2024onpolicy}) đã phát hiện ra hai sai lầm kinh điển của phương pháp chưng cất truyền thống:

\subsubsection{Sai lầm 1: Dùng Forward KL khiến mô hình nhỏ bị ``Bắt cá hai tay'' (Mean-Seeking)}

Độ phân kỳ Kullback-Leibler không có tính đối xứng: $D_{\text{KL}}(P_T \parallel \pi_\theta) \neq D_{\text{KL}}(\pi_\theta \parallel P_T)$.
\begin{itemize}
  \item \textbf{Forward KL (Chuẩn truyền thống):} $D_{\text{KL}}(P_T \parallel \pi_\theta) = \sum_{y} P_T(y \mid x) \log \frac{P_T(y \mid x)}{\pi_\theta(y \mid x)}$.
  Hãy nhìn vào công thức này: bất cứ chỗ nào mô hình Thầy gán xác suất $P_T(y \mid x) > 0$, mô hình Trò \textbf{bắt buộc phải gán $\pi_\theta(y \mid x) > 0$} (vì nếu $\pi_\theta(y \mid x) \to 0$, thì $\log \frac{P_T}{0} \to +\infty$ làm nổ hàm mất mát!).
  Với một câu hỏi mở, mô hình Thầy 671B rất uyên bác nên nó biết tới 10 cách trả lời đúng khác nhau (phân phối có 10 đỉnh --- 10 modes). Mô hình Trò 7B có dung lượng quá nhỏ không thể nhớ hết cả 10 cách, nhưng bị Forward KL ép phải phủ hết cả 10 đỉnh! Kết quả là mô hình Trò phải kéo giãn phân phối của nó ra nằm chính giữa các đỉnh (\textbf{Zero-avoiding / Mean-seeking}) --- tức là nó gán xác suất cao vào đúng vùng thung lũng \textbf{giữa các câu đúng} (nơi mà $P_T = 0$, đồng nghĩa với việc sinh ra câu trả lời nửa nạc nửa mỡ, bịa đặt!).
  \item \textbf{Reverse KL (Đảo ngược chiều KL --- Chuẩn \texttt{MiniLLM} \citealt{gu2024minillm}):}
  \begin{equation}
    D_{\text{KL}}(\pi_\theta \parallel P_T) = \sum_{y} \pi_\theta(y \mid x) \log \frac{\pi_\theta(y \mid x)}{P_T(y \mid x)}
    \label{eq:reverse_kl_distill}
  \end{equation}
  Với Reverse KL, nếu tại một câu trả lời $y$ nào đó mô hình Thầy gán $P_T(y \mid x) = 0$ (câu sai), mô hình Trò bắt buộc phải gán $\pi_\theta(y \mid x) = 0$ để không bị phạt vô cực. Ngược lại, trong số 10 cách trả lời đúng của Thầy, nếu Trò chỉ chọn học thật xuất sắc \textbf{1 cách duy nhất} và bỏ qua 9 cách còn lại ($\pi_\theta = 0$ ở 9 đỉnh kia), thì $0 \log(0/P_T) = 0$ --- Trò hoàn toàn không bị phạt! Đặc tính này gọi là \textbf{Mode-seeking (Tập trung vào một đỉnh chất lượng cao)}, giúp mô hình nhỏ trả lời cực kỳ sắc bén và không bị ảo giác!
\end{itemize}

\subsubsection{Sai lầm 2: Chỉ học trên câu của Thầy mà không sửa lỗi trên câu của Trò (On-Policy GKD)}

Trong chưng cất tĩnh (Offline Distillation), ta chỉ tính KL trên các câu văn do Thầy viết sẵn. Giống hệt bệnh Exposure Bias của SFT (Chương~\ref{ch:rl_ppo}), khi Trò tự viết văn lúc chạy thực tế và lỡ bước chệch hướng, nó không biết phải quay lại như thế nào.

Giải pháp \textbf{On-Policy GKD (Generalized Knowledge Distillation)} của Google DeepMind \citep{agarwal2024onpolicy} đảo ngược quy trình:
\begin{enumerate}
  \item Cho chính \textbf{mô hình Trò $\pi_\theta$ tự sinh ra câu trả lời} $y \sim \pi_\theta(\cdot \mid x)$ (On-policy rollout).
  \item Đưa câu trả lời $y$ mà Trò vừa tự viết sang cho mô hình Thầy $P_T$ đọc (chỉ tốn 1 lượt Forward pass cực nhanh của Thầy, không cần Thầy phải sinh văn bản!).
  \item Tại từng vị trí token $t$ trên chính câu văn của Trò, Thầy đưa ra phân phối xác suất chuẩn $P_T(\cdot \mid x, y_{<t})$ để dạy cho Trò biết: \textit{``Tại đúng cái chỗ mà em vừa viết nhầm này, lẽ ra từ tiếp theo em phải chọn là từ này mới đúng!''}
\end{enumerate}

% =========================================================================
\section{Hợp nhất Mô hình (Model Merging): ``Hợp thể'' Đa kỹ năng Không cần GPU}
\label{sec:model_merging}

Hãy tưởng tượng một tình huống rất phổ biến tại các công ty AI Việt Nam:
\begin{itemize}
  \item Nhóm A lấy mô hình nền \texttt{Qwen2.5-7B-Base} đi huấn luyện trên dữ liệu tiếng Việt, tạo ra mô hình $\theta_{\text{VN}}$ rất giỏi văn phong Việt Nam.
  \item Nhóm B cũng lấy đúng mô hình nền \texttt{Qwen2.5-7B-Base} đó đi huấn luyện trên dữ liệu lập trình, tạo ra mô hình $\theta_{\text{Code}}$ rất giỏi viết Python/SQL.
  \item Nhóm C lại lấy \texttt{Qwen2.5-7B-Base} đi huấn luyện giải Toán, tạo ra $\theta_{\text{Math}}$.
\end{itemize}
Bây giờ Giám đốc Sản phẩm yêu cầu một mô hình duy nhất vừa giỏi Tiếng Việt, vừa giỏi Code, vừa giỏi Toán. Chẳng lẽ bạn phải trộn lại toàn bộ 3 tập dữ liệu và thuê cụm GPU huấn luyện lại từ đầu tốn hàng chục nghìn USD?

Câu trả lời của giai đoạn 2024--2026 là \textbf{KHÔNG CẦN}! Bạn có thể gộp trực tiếp ma trận trọng số của cả 3 mô hình trên CPU chỉ trong vòng 2 phút bằng kỹ thuật \textbf{Model Merging}.

\subsection{Vector Tác vụ (Task Vectors) và Tại sao Cộng Trung bình Ngây thơ lại Thất bại?}

Vì cả $K$ mô hình chuyên biệt $\theta_1, \theta_2, \dots, \theta_K$ đều xuất phát từ cùng một mô hình gốc $\theta_{\text{base}}$, chúng nằm trong cùng một bồn trũng không gian tham số (Linear Mode Connectivity). Ta định nghĩa \textbf{Vector tác vụ (Task Vector / Delta Parameter)} $\tau_k$ của mô hình thứ $k$ là phần độ lệch trọng số mà nó đã học thêm được so với mô hình gốc:
\begin{equation}
  \tau_k = \theta_k - \theta_{\text{base}} \quad \text{với } k \in \{1, \dots, K\}
  \label{eq:task_vector}
\end{equation}

Cách gộp ngây thơ nhất (Task Arithmetic / Linear Averaging) là cộng trung bình các vector tác vụ này lại: $\theta_{\text{merged}} = \theta_{\text{base}} + \frac{1}{K}\sum_{k=1}^K \tau_k$. Tuy nhiên, khi cộng trực tiếp như vậy, điểm số của mô hình gộp thường bị tụt giảm nghiêm trọng so với từng mô hình thành phần. Tại sao?

Nghiên cứu đột phá tại NeurIPS 2023 của \citet{yadav2023ties} đã chỉ ra \textbf{2 nguyên nhân gây nhiễu tham số (Parameter Interference)}:
\begin{enumerate}
  \item \textbf{Nhiễu từ 80\%--90\% các tham số dư thừa (Redundant Parameters):} Khi Fine-tune một mô hình 7 tỷ tham số, thực chất chỉ có khoảng $10\% - 20\%$ các trọng số thay đổi mạnh là mang tri thức cốt lõi của tác vụ mới; $80\%$ trọng số còn lại chỉ dao động nhẹ quanh số $0$ (nhiễu ngẫu nhiên do SGD). Khi cộng trung bình nhiều mô hình, đám nhiễu 80\% này chồng lấn lên nhau và làm loãng đi 20\% trọng số quan trọng của mô hình kia!
  \item \textbf{Xung đột Dấu (Sign Disagreement):} Tại cùng một vị trí trọng số thứ $i$, mô hình Tiếng Việt muốn tăng nó lên ($\tau_{1,i} = +0.05$) nhưng mô hình Code lại muốn giảm nó xuống ($\tau_{2,i} = -0.04$). Khi cộng lại, chúng triệt tiêu nhau về $0.005$, làm hỏng cả hai kỹ năng!
\end{enumerate}

\subsection{Giải pháp 1: TIES-Merging (Trim, Elect Sign, \& Merge) \citep{yadav2023ties}}

Để giải quyết triệt để hai xung đột trên, thuật toán \textbf{TIES-Merging} \citep{yadav2023ties} thực hiện quy trình 3 bước cực kỳ gọn gàng trên các vector tác vụ $\tau_k = \theta_k - \theta_{\text{base}}$:

\begin{enumerate}
  \item \textbf{Bước 1 --- TRIM (Cắt tỉa nhiễu):} Với mỗi vector tác vụ $\tau_k$, chỉ giữ lại top-$p\%$ (thường chọn $20\%$) các phần tử có độ lớn tuyệt đối $|\tau_{k,i}|$ lớn nhất; toàn bộ $80\%$ phần tử nhỏ còn lại bị gán thẳng về $0$, thu được vector thưa $\hat{\tau}_k$.
  \item \textbf{Bước 2 --- ELECT SIGN (Bầu cử chọn Dấu chung cuộc):} Tại mỗi vị trí trọng số thứ $i$, tính tổng toàn bộ các giá trị $\sum_{k=1}^K \hat{\tau}_{k,i}$. Nếu tổng này dương, dấu chiến thắng chung cuộc tại vị trí $i$ là $\gamma_i = +1$; nếu tổng âm, dấu chiến thắng là $\gamma_i = -1$:
  \begin{equation}
    \gamma_i = \text{sign}\left( \sum_{k=1}^{K} \hat{\tau}_{k, i} \right)
  \end{equation}
  \item \textbf{Bước 3 --- DISJOINT MERGE (Cộng trung bình không xung đột):} Tại vị trí thứ $i$, loại bỏ toàn bộ những mô hình nào có dấu ngược với dấu chiến thắng $\gamma_i$, rồi chỉ tính trung bình cộng trên tập các mô hình cùng phe $\mathcal{A}_i = \{k \mid \text{sign}(\hat{\tau}_{k,i}) = \gamma_i\}$:
  \begin{equation}
    \tau_{\text{TIES}, i} = \frac{1}{|\mathcal{A}_i|} \sum_{k \in \mathcal{A}_i} \hat{\tau}_{k, i}, \qquad \theta_{\text{merged}} = \theta_{\text{base}} + \lambda \, \tau_{\text{TIES}}
    \label{eq:ties_merge}
  \end{equation}
\end{enumerate}

\begin{mathdecode}[Chạy thử bằng tay thuật toán TIES-Merging trên 3 trọng số]
Giả sử ta gộp 3 mô hình ($\tau_1, \tau_2, \tau_3$) tại một trọng số $i$ duy nhất sau bước Trim:
\begin{equation*}
  \hat{\tau}_{1, i} = +0.06, \qquad \hat{\tau}_{2, i} = +0.04, \qquad \hat{\tau}_{3, i} = -0.05
\end{equation*}
\begin{itemize}
  \item \textbf{Nếu cộng trung bình ngây thơ:} $\frac{+0.06 + 0.04 - 0.05}{3} = \frac{+0.05}{3} = \mathbf{+0.0167}$ (tín hiệu bị kéo tụt thê thảm xuống gần $0$ do mô hình 3 kéo ngược chiều!).
  \item \textbf{Nếu dùng TIES-Merging:}
  \begin{enumerate}
    \item \textit{Bầu cử dấu:} Tổng khối lượng phe Dương là $+0.06 + 0.04 = +0.10$; phe Âm là $-0.05$. Tổng chung cuộc $= +0.05 > 0 \implies$ Phe Dương thắng cử ($\gamma_i = +1$).
    \item \textit{Cộng cùng phe:} Loại bỏ $\hat{\tau}_{3,i} = -0.05$ vì ngược dấu. Chỉ tính trung bình trên tập cùng phe Dương $\mathcal{A}_i = \{1, 2\}$:
    \begin{equation*}
      \tau_{\text{TIES}, i} = \frac{+0.06 + 0.04}{2} = \mathbf{+0.050}
    \end{equation*}
  \end{enumerate}
\end{itemize}
Cường độ tín hiệu $+0.050$ được bảo toàn nguyên vẹn, cao gấp 3 lần so với cách cộng ngây thơ!
\end{mathdecode}

\subsection{Giải pháp 2: DARE (Drop And REscale) \citep{yu2024language}}

Năm 2024, nhóm nghiên cứu của \citet{yu2024language} tại ICML 2024 đã công bố một phát hiện còn gây kinh ngạc hơn nữa mang tên \textbf{DARE (Drop And REscale)}:

Trong các mô hình ngôn ngữ lớn đã qua SFT, các vector độ lệch $\tau_k = \theta_k - \theta_{\text{base}}$ có mức độ dư thừa cực đoan đến mức: bạn có thể \textbf{xóa ngẫu nhiên tới 90\%--99\% các phần tử của $\tau_k$ về $0$}, miễn là bạn \textbf{nhân bù (Rescale)} $10\%$ phần tử sống sót còn lại với hệ số $\frac{1}{1 - p}$ (tức là nhân gấp $10$ lần khi xóa $p = 0.9$), thì năng lực của mô hình gần như \textbf{không suy giảm một chút nào}!

\begin{equation}
  m_{k, i} \sim \text{Bernoulli}(1 - p), \qquad \tilde{\tau}_{k, i} = \frac{m_{k, i} \cdot \tau_{k, i}}{1 - p}
  \label{eq:dare_formula}
\end{equation}

\begin{mentalmodel}[Tại sao xóa 90\% rồi nhân 10 phần còn lại của DARE lại hoạt động kỳ diệu?]
Hãy nhớ lại cách một lớp Linear trong mạng nơ-ron hoạt động: đầu ra là tổng tích vô hướng $y = \sum_{i=1}^{d} (\theta_{\text{base}, i} + \tau_i) x_i = \sum \theta_{\text{base}, i} x_i + \sum_{i=1}^{d} \tau_i x_i$.
Vì chiều ẩn $d$ của LLM rất lớn ($d = 4{,}096$ hoặc $8{,}192$), theo Luật Số Lớn trong xác suất, khi bạn giữ lại ngẫu nhiên $10\%$ các số hạng trong tổng $\sum_{i=1}^{4096} \tau_i x_i$ rồi nhân mỗi số hạng sống sót lên $10$ lần, \textbf{giá trị kỳ vọng của tổng đầu ra hoàn toàn không đổi}:
\begin{equation*}
  \mathbb{E}\left[ \sum_{i=1}^{d} \tilde{\tau}_i x_i \right] = \sum_{i=1}^{d} \mathbb{E}\left[\frac{m_i}{1 - p}\right] \tau_i x_i = \sum_{i=1}^{d} \frac{1 - p}{1 - p} \tau_i x_i = \sum_{i=1}^{d} \tau_i x_i
\end{equation*}
Nhưng điều tuyệt vời nhất xảy ra khi bạn mang hai mô hình đã qua DARE (mỗi mô hình đã xóa trống 90\% vị trí) đi gộp vào nhau: xác suất để hai phần tử khác $0$ của hai mô hình đâm sầm vào cùng một vị trí chỉ còn là $0.1 \times 0.1 = \mathbf{1\%}$! $99\%$ các vị trí trên ma trận hoàn toàn không hề va chạm nhau!
\end{mentalmodel}

Trong thực tế năm 2025--2026, tổ hợp mạnh nhất trên thư viện mã nguồn mở \texttt{mergekit} chính là \textbf{DARE-TIES}: dùng DARE để làm thưa và nhân bù các vector tác vụ trước, sau đó dùng quy tắc bầu cử dấu của TIES để xử lý $1\%$ các vị trí va chạm còn lại!

% =========================================================================
% CHƯƠNG 6: TỐI ƯU HỆ THỐNG, ĐÁNH GIÁ MÔ HÌNH & PLAYBOOK GỠ LỖI THỰC CHIẾN
% =========================================================================
\chapter{Tối ưu Hệ thống, Đánh giá Mô hình và Playbook Gỡ lỗi Thực chiến}
\label{ch:systems_eval}

Một AI Engineer giỏi không chỉ hiểu công thức của SFT, DPO hay GRPO, mà còn phải biết tính nhẩm chính xác lượng bộ nhớ GPU (VRAM) trước khi bấm nút chạy, biết lựa chọn giữa LoRA và Full Fine-Tuning, biết thiết kế hệ thống chấm điểm không bị lừa bởi thiên kiến độ dài, và quan trọng nhất: khi nhìn vào đồ thị hàm mất mát bị nổ hoặc mô hình sinh ra toàn ký tự rác lúc 2 giờ sáng, phải chẩn đoán được ngay nguyên nhân gốc rễ.

Chương cuối cùng này trang bị cho bạn toàn bộ bộ công cụ thực chiến đó.

% =========================================================================
\section{Toán học Tính nhẩm Bộ nhớ GPU (Back-of-the-Envelope VRAM Math)}
\label{sec:vram_math}

\subsection{Quy tắc \texorpdfstring{$16\Psi$}{16 Psi}: Tại sao Mô hình 7B lại ngốn hơn 112 GB VRAM khi Huấn luyện?}

Một câu hỏi kinh điển của các kỹ sư mới vào nghề: \textit{``Mô hình 7 tỷ tham số ($\Psi = 7 \times 10^9$) lưu ở định dạng 16-bit (BF16, mỗi tham số tốn 2 bytes) chỉ nặng có $7 \times 2 = 14\text{ GB}$. Tại sao tôi cắm vào card RTX 4090 (24 GB) hay thậm chí A100 (80 GB) để chạy Full Fine-Tuning lại bị báo lỗi tràn bộ nhớ (Out-Of-Memory --- OOM) ngay lập tức?''}

Câu trả lời nằm ở kiến trúc \textbf{Mixed-Precision AdamW} \citep{micikevicius2018mixed, loshchilov2019decoupled, rajbhandari2020zero}. Để huấn luyện ổn định mà không bị mất mát các bước cập nhật siêu nhỏ (Underflow), mỗi tham số của mô hình đòi hỏi GPU phải lưu trữ đồng thời \textbf{4 thành phần trạng thái tĩnh}:

\begin{enumerate}
  \item \textbf{Trọng số mô hình (Weights --- BF16):} $2\text{ bytes} \times \Psi$.
  \item \textbf{Đạo hàm (Gradients --- BF16):} $2\text{ bytes} \times \Psi$.
  \item \textbf{Bản sao Trọng số Chủ (FP32 Master Weights):} $4\text{ bytes} \times \Psi$.
  (Tại sao phải có bản sao 32-bit này? Vì ở định dạng 16-bit BF16, khoảng cách nhỏ nhất giữa hai số liên tiếp quanh số $1.0$ là $\approx 0.0078$. Trong khi đó, bước cập nhật $\text{LR} \times g_t \approx 10^{-5} \times 0.1 = 10^{-6}$ nhỏ hơn hàng nghìn lần! Nếu cộng trực tiếp $1.0 + 10^{-6}$ ở BF16, máy tính sẽ làm tròn kết quả về đúng $1.0$ --- mô hình đứng im không học được gì cả! Do đó, bắt buộc phải cộng dồn trên một bản sao 32-bit FP32).
  \item \textbf{Hai trạng thái lịch sử của bộ tối ưu hóa AdamW (Momentum $m_t$ và Variance $v_t$ --- FP32):} $(4 + 4)\text{ bytes} \times \Psi = 8\text{ bytes} \times \Psi$.
\end{enumerate}

Tổng cộng, chỉ riêng phần bộ nhớ tĩnh (chưa tính dữ liệu đầu vào!) đã tốn:
\begin{equation}
  M_{\text{static}} = (2 + 2 + 4 + 4 + 4) \times \Psi = \mathbf{16\Psi \text{ bytes}}
  \label{eq:16psi_rule}
\end{equation}
Với mô hình $\Psi = 7\text{ tỷ tham số}$, $M_{\text{static}} = 16 \times 7 = \mathbf{112\text{ GB VRAM}}$! Để nhét vừa 112 GB này vào một cụm 4 card GPU 40GB, chúng ta phải dùng kỹ thuật \textbf{FSDP (Fully Sharded Data Parallel)} \citep{zhao2023pytorch} hoặc \textbf{DeepSpeed ZeRO-3} \citep{rajbhandari2020zero} để cắt nhỏ (shard) toàn bộ $16\Psi$ bytes này chia đều cho $N_{\text{GPU}}$ card (mỗi card chỉ còn giữ $\frac{112}{4} = 28\text{ GB}$).

\subsection{``Sát thủ giấu mặt'' làm tràn VRAM: Ma trận Logits ở Đầu ra Từ vựng (LM Head)}

Nhiều kỹ sư đã bật ZeRO-3, bật cả Gradient Checkpointing và chỉ chạy với Batch Size cực nhỏ ($B = 2$), nhưng khi nâng chiều dài chuỗi lên $T = 8{,}192$ tokens trên các mô hình đời mới (như Llama 3 hay Qwen 2.5), GPU vẫn nổ OOM ngay tại bước tính hàm mất mát cuối cùng! Thủ phạm là gì?

\begin{mathdecode}[Tính nhẩm kích thước Ma trận Logits tại lớp LM Head]
Các mô hình thế hệ 2024--2026 có bảng từ vựng $|\mathcal{V}|$ cực lớn để hỗ trợ đa ngôn ngữ và code (ví dụ Llama 3 có $|\mathcal{V}| = 128{,}256$; Qwen 2.5 có $|\mathcal{V}| = 152{,}064$).
Ở lớp cuối cùng của mạng Transformer, để tính hàm mất mát Cross-Entropy (trong SFT) hoặc $\log \pi_\theta$ (trong DPO/GRPO), mô hình phải tạo ra tensor \texttt{logits} có kích thước $(B, T, |\mathcal{V}|)$ và ép kiểu lên FP32 (4 bytes) để tính hàm \texttt{log\_softmax} chính xác:
\begin{itemize}
  \item Với $B = 2$, $T = 8{,}192$, $|\mathcal{V}| = 128{,}256$:
  \begin{equation*}
    M_{\text{logits}} = 2 \times 8{,}192 \times 128{,}256 \times 4\text{ bytes} \approx \mathbf{8.4\text{ GB}}
  \end{equation*}
  \item Khi gọi \texttt{.backward()}, PyTorch tạo thêm một tensor đạo hàm có kích thước y hệt ($8.4\text{ GB}$ nữa).
  \item $\implies$ Chỉ riêng 1 dòng tính hàm mất mát ở lớp cuối cùng đã ngốn đột ngột $\mathbf{16.8\text{ GB VRAM}}$ trên mỗi card GPU!
\end{itemize}
\textbf{Giải pháp kỹ thuật hiện đại (\texttt{Liger-Kernel} / Fused Linear Cross-Entropy):} Thay vì nhân ra toàn bộ ma trận khổng lồ $(B \times T, |\mathcal{V}|)$ cùng một lúc, kỹ thuật \textit{Fused Linear Cross-Entropy} cắt nhỏ chiều thời gian $T$ thành từng khúc nhỏ (ví dụ 512 tokens), tính \texttt{log\_softmax} và đạo hàm ngay trong bộ nhớ đệm SRAM của GPU rồi xóa ngay trước khi sang khúc tiếp theo. Nhờ đó, bộ nhớ đỉnh của lớp LM Head giảm từ $16.8\text{ GB}$ xuống chỉ còn dưới $\mathbf{200\text{ MB}}$!
\end{mathdecode}

% =========================================================================
\section{Kiến trúc Hạ tầng RL Hiện đại (\texttt{HybridFlow} / \texttt{verl})}
\label{sec:hybridflow_verl}

Khi chạy các thuật toán Online RL như PPO hay GRPO (Chương~\ref{ch:rl_ppo} và Chương~\ref{ch:reasoning_rl}), mỗi vòng lặp huấn luyện đều bị chia thành hai pha có bản chất vật lý trái ngược hoàn toàn nhau:

\begin{enumerate}
  \item \textbf{Pha Sinh văn bản (Rollout / Autoregressive Decoding):} Mô hình sinh ra từng chữ một cho $G$ câu trả lời dài hàng nghìn tokens. Ở pha này, để sinh ra 1 token mới, GPU phải đọc toàn bộ hàng chục Gigabyte trọng số mô hình từ bộ nhớ HBM vào chip xử lý nhưng chỉ thực hiện vài phép cộng nhân vector nhỏ lẻ. Do đó, pha này bị \textbf{nghẽn cổ chai ở băng thông đọc bộ nhớ (Memory-Bandwidth Bound)}, chiếm tới $70\% - 80\%$ tổng thời gian chạy RL! Để chạy nhanh pha này, bắt buộc phải dùng các engine suy luận chuyên dụng như \textbf{vLLM} hoặc \textbf{SGLang} \citep{kwon2023efficient} với kiểu chia mô hình \textbf{Tensor Parallelism (TP)}.
  \item \textbf{Pha Cập nhật Đạo hàm (Training / Forward \& Backward):} Toàn bộ $G$ câu trả lời đã có sẵn nên GPU nhân hai ma trận khổng lồ cùng lúc. Pha này bị \textbf{nghẽn ở năng lực tính toán FLOPs (Compute Bound)} và cần kiểu chia bộ nhớ \textbf{FSDP / ZeRO} \citep{zhao2023pytorch, rajbhandari2020zero}.
\end{enumerate}

\begin{mentalmodel}[Cách thư viện \texttt{verl} (HybridFlow) giải cứu hạ tầng RL]
Trước năm 2024, để kết hợp vLLM (cho pha Rollout) và FSDP (cho pha Training), các kỹ sư phải lưu trọng số ra ổ cứng hoặc truyền qua mạng TCP rất chậm chạp ở mỗi bước RL.
Kiến trúc \textbf{HybridFlow (\texttt{verl})} \citep{sheng2024hybridflow} (được sử dụng rộng rãi để huấn luyện các mô hình Reasoning năm 2025--2026) đã giải quyết triệt để bằng cơ chế \textbf{3D-HybridEngine}:
\begin{itemize}
  \item Trong pha Training, trọng số nằm ở định dạng cắt nhỏ FSDP.
  \item Ngay khi chuyển sang pha Rollout, \texttt{verl} thực hiện một phép gom nhóm trực tiếp trên đường truyền siêu tốc NVLink giữa các GPU (\texttt{All-Gather}) để biến đổi bố cục bộ nhớ từ FSDP sang đúng bố cục Tensor Parallelism của vLLM chỉ trong vài phần nghìn giây mà không chạm vào ổ cứng!
  \item Sau khi vLLM sinh xong nhóm $G$ câu trả lời, bộ nhớ KV-Cache của vLLM lập tức được giải phóng để nhường chỗ cho các tensor đạo hàm của pha Training.
\end{itemize}
\end{mentalmodel}

% =========================================================================
\section{Phản biện Kỹ thuật: Khi nào dùng LoRA / QLoRA, Khi nào bắt buộc Full Fine-Tuning?}
\label{sec:lora_vs_full}

Khi ngân sách GPU hạn hẹp, \textbf{LoRA (Low-Rank Adaptation)} \citep{hu2022lora} và \textbf{QLoRA} \citep{dettmers2023qlora} là phao cứu sinh của mọi kỹ sư: thay vì cập nhật toàn bộ ma trận trọng số gốc $\mathbf{W}_0 \in \mathbb{R}^{d \times k}$, ta đóng băng $\mathbf{W}_0$ (thậm chí nén $\mathbf{W}_0$ xuống 4-bit NormalFloat4 trong QLoRA) và chỉ gắn thêm hai ma trận hạng thấp mỏng dính $\mathbf{A} \in \mathbb{R}^{r \times k}$ và $\mathbf{B} \in \mathbb{R}^{d \times r}$ với $r \ll \min(d, k)$ (thường chọn $r \in \{16, 32, 64\}$):
\begin{equation}
  \mathbf{W} = \mathbf{W}_0 + \frac{\alpha}{r} \mathbf{B}\mathbf{A}
  \label{eq:lora_formula}
\end{equation}
Vì số lượng tham số cần huấn luyện trong $\mathbf{A}$ và $\mathbf{B}$ chỉ chiếm chưa tới $1\%$ mô hình gốc, toàn bộ $12\Psi$ bytes của bộ tối ưu hóa AdamW và Master Weights biến mất! Bạn có thể huấn luyện mô hình 7B chỉ với 12 GB VRAM, hoặc 70B trên một card 48 GB duy nhất.

Nhưng một AI Engineer chuyên nghiệp không bao giờ dùng công cụ một cách mù quáng. Hãy giải đáp 3 câu hỏi hóc búa nhất về LoRA:

\begin{debate}
\textbf{Câu hỏi 1: ``LoRA có tốt bằng Full Fine-Tuning không? Khi nào LoRA thất bại?''}

\textbf{Đáp:} Quy luật cốt lõi được giới nghiên cứu đúc kết là: \textbf{``LoRA học được ít hơn nhưng cũng quên ít hơn'' (LoRA learns less and forgets less)}!
\begin{itemize}
  \item \textbf{Khi nào LoRA đạt 99\% hiệu quả của Full Fine-Tuning (Thậm chí tốt hơn):} Khi tác vụ của bạn thuộc nhóm \textit{Định hình hành vi / Phong cách / Định dạng} (như SFT hội thoại, DPO căn chỉnh sở thích, hoặc chạy GRPO trên một miền hẹp). Lúc này mô hình không cần nạp thêm quá nhiều tri thức mới; chiều nội tại của sự thay đổi rất nhỏ \citep{aghajanyan2021intrinsic}. Hơn nữa, vì $\mathbf{W}_0$ bị khóa cứng, LoRA đóng vai trò như một bộ điều chuẩn tự nhiên giúp mô hình ít bị quên kiến thức cũ (Catastrophic Forgetting) hơn hẳn Full Fine-Tuning!
  \item \textbf{Khi nào LoRA thất bại nặng nề trước Full Fine-Tuning:} Khi bạn muốn \textit{bơm một khối lượng tri thức hoàn toàn mới khổng lồ} vào mô hình --- ví dụ: Continual Pre-training dạy thêm một ngôn ngữ mới mà Base Model chưa biết, hoặc Chưng cất (Distillation) 800,000 chuỗi suy nghĩ Toán/Code phức tạp từ DeepSeek-R1 xuống một mô hình nhỏ. Các ma trận hạng thấp $r = 16$ hay $64$ bị quá tải dung lượng (Capacity Bottleneck) và không thể biểu diễn hết sự dịch chuyển trọng số đa chiều này!
\end{itemize}

\medskip
\textbf{Câu hỏi 2: ``Tại sao trong bài báo LoRA gốc năm 2021 người ta khuyên chỉ cần gắn LoRA vào 2 ma trận Attention là \texttt{q\_proj} và \texttt{v\_proj}, nhưng đến năm 2025--2026 mọi người đều khuyên phải gắn vào TẤT CẢ các lớp Linear (\texttt{all-linear})?''}

\textbf{Đáp:} Bài báo \citet{dettmers2023qlora} đã chứng minh thực nghiệm rằng: trong kiến trúc Transformer hiện đại, các khối mạng truyền thẳng \textbf{MLP / SwiGLU} (\texttt{gate\_proj}, \texttt{up\_proj}, \texttt{down\_proj}) mới là nơi lưu trữ phần lớn tri thức và năng lực suy luận của mô hình! Nếu bạn chỉ gắn LoRA vào \texttt{q\_proj, v\_proj} (dù có tăng hạng $r$ lên $128$), chất lượng vẫn thua xa việc gắn LoRA hạng nhỏ $r = 16$ nhưng phủ đều lên \textbf{toàn bộ 7 ma trận Linear} (\texttt{q\_proj, k\_proj, v\_proj, o\_proj, gate\_proj, up\_proj, down\_proj})!

\medskip
\textbf{Câu hỏi 3: ``Tại sao Learning Rate của LoRA luôn phải đặt lớn gấp 10 lần so với Full Fine-Tuning (ví dụ $2 \times 10^{-4}$ thay vì $2 \times 10^{-5}$)?''}

\textbf{Đáp:} Bởi vì trong LoRA, ma trận $\mathbf{B}$ được khởi tạo bằng $\mathbf{0}$ hoàn toàn ($\mathbf{B}_0 = \mathbf{0}$) để tại bước đầu tiên $\mathbf{B}\mathbf{A} = \mathbf{0}$ không làm xáo trộn mô hình gốc. Do bắt đầu từ điểm yên ngựa $\mathbf{B} = \mathbf{0}$ và chỉ cập nhật trên một không gian con rất nhỏ, các ma trận $\mathbf{A}, \mathbf{B}$ cần một tốc độ học lớn gấp $\approx 10$ lần để tạo ra cùng một biên độ thay đổi đầu ra $\Delta \mathbf{h} = \frac{\alpha}{r}\mathbf{B}\mathbf{A}\mathbf{x}$ như khi cập nhật đồng loạt hàng tỷ tham số của toàn bộ ma trận $\mathbf{W}$!
\end{debate}

% =========================================================================
\section{Đánh giá Khoa học (Evaluation) và Khử Thiên kiến LLM-as-a-Judge}
\label{sec:evaluation}

Làm sao bạn biết mô hình sau khi chạy Post-Training của mình thực sự thông minh hơn hay chỉ đang ``học vẹt và nói dài dòng hơn''?

\subsection{Thuế Căn chỉnh (Alignment Tax) \citep{ouyang2022training, bai2022training}}

Khi bạn ép mô hình phải từ chối các câu hỏi nhạy cảm hoặc phải trả lời theo một khuôn mẫu cố định, năng lực giải quyết các bài toán kiến thức nền tảng của nó có thể bị sụt giảm --- khoản sụt giảm này được gọi là \textbf{Thuế Căn chỉnh (Alignment Tax)}. Do đó, một quy trình đánh giá chuẩn bắt buộc phải đo đồng thời hai nhóm thước đo:
\begin{enumerate}
  \item \textbf{Thước đo Năng lực Cốt lõi \& Chống Học tủ (Contamination-Free Verifiable Benchmarks):}
  \begin{itemize}
    \item \textbf{IFEval (Instruction Following Evaluation)} \citep{zhou2023instruction}: Kiểm tra bằng code (không dùng AI chấm) xem mô hình có tuân thủ chính xác các ràng buộc cứng như \textit{``Viết đúng 3 đoạn văn, không dùng chữ cái 'e', kết thúc bằng câu hỏi''} hay không.
    \item \textbf{LiveCodeBench} \citep{jain2024livecodebench}: Liên tục cập nhật các đề thi lập trình mới từ LeetCode, AtCoder và Codeforces \textit{sau ngày phát hành mô hình} để đảm bảo không mô hình nào có thể gian lận bằng cách học thuộc đề thi trong tập Pre-training.
    \item \textbf{SWE-bench Verified} \citep{jimenez2024swebench}: Giao cho mô hình một kho mã nguồn GitHub thực tế bị lỗi và yêu cầu viết bản vá (git patch) vượt qua toàn bộ unit tests của dự án.
    \item \textbf{MATH-500 \& AIME 2024/2025} \citep{hendrycks2021measuring}: Thước đo vàng cho các mô hình suy luận Toán học.
  \end{itemize}
  \item \textbf{Thước đo Sở thích Mở: Chatbot Arena Elo \citep{chiang2024chatbot} và LLM-as-a-Judge \citep{zheng2023judging}:}
  Với các câu hỏi mở (viết email, tóm tắt, tư vấn), hệ thống \textbf{Chatbot Arena} \citep{chiang2024chatbot} cho hai mô hình ẩn danh đấu cặp trực tiếp trước người dùng thật và tính điểm \textbf{Elo} dựa trên hàm hợp lý cực đại của mô hình Bradley-Terry:
  \begin{equation}
    P(\text{Model } A \succ \text{Model } B) = \frac{1}{1 + 10^{(R_B - R_A) / 400}}
    \label{eq:elo_prob}
  \end{equation}
  Theo công thức~\eqref{eq:elo_prob}, nếu mô hình $A$ cao hơn mô hình $B$ đúng $100$ điểm Elo ($R_A - R_B = 100$), tỷ lệ thắng đối đầu của $A$ trước $B$ là $\frac{1}{1 + 10^{-100/400}} \approx \mathbf{64\%}$.
\end{enumerate}

\subsection{Ba Thiên kiến Chí mạng của LLM-as-a-Judge và Thuốc giải LC-AlpacaEval}

Trong quá trình phát triển hàng ngày, không ai có thể đợi người dùng thật trên Chatbot Arena chấm điểm. Chúng ta phải dùng một mô hình lớn (như GPT-4o hoặc Claude 3.5/3.7 Sonnet) làm trọng tài tự động (\textbf{LLM-as-a-Judge}) \citep{zheng2023judging}. Tuy nhiên, trọng tài AI mắc \textbf{3 căn bệnh thiên kiến nặng nề}:

\begin{enumerate}
  \item \textbf{Thiên kiến Vị trí (Position Bias):} Trọng tài AI có xu hướng thích câu trả lời được đặt ở vị trí đầu tiên (Response A) hơn vị trí thứ hai (Response B).
  \begin{itemize}
    \item \textit{Cách khử bắt buộc (Position Swapping):} Mọi cặp câu trả lời phải được chấm \textbf{2 lần} với thứ tự đảo ngược $(A, B)$ và $(B, A)$. Chỉ công nhận chiến thắng khi cả 2 lần chấm đều chỉ về cùng một mô hình; nếu 2 lần chấm mâu thuẫn nhau thì tính là Hòa (Tie)!
  \end{itemize}
  \item \textbf{Thiên kiến Tự yêu mình (Self-Enhancement Bias):} Trọng tài họ GPT thường chấm điểm cao hơn cho văn phong của các mô hình họ GPT; trọng tài họ Claude thích văn phong của Claude.
  \begin{itemize}
    \item \textit{Cách khử:} Sử dụng hội đồng giám khảo đa họ mô hình (Jury of Judges) hoặc dùng mô hình trọng tài khác họ với mô hình đang được huấn luyện.
  \end{itemize}
  \item \textbf{Thiên kiến Độ dài (Verbosity / Length Bias):} Đây là căn bệnh nguy hiểm nhất! Trọng tài AI cực kỳ dễ bị lừa bởi các câu trả lời dài dòng, nhiều gạch đầu dòng và chữ in đậm, ngay cả khi nội dung rỗng tuếch. Trên bảng xếp hạng AlpacaEval đời đầu, một mô hình chỉ cần viết dài gấp đôi là tỷ lệ thắng nhảy vọt từ $30\%$ lên $55\%$!
\end{enumerate}

\paragraph{Thuốc giải chuẩn công nghiệp: Length-Controlled AlpacaEval (LC-AlpacaEval) \citep{dubois2024length}:}
Để ngăn chặn việc gian lận bằng độ dài câu, \citet{dubois2024length} khớp một mô hình hồi quy Logistic tổng quát tách biệt tác động thực sự của chất lượng mô hình ($\theta_m$) khỏi tác động của chênh lệch độ dài câu ($\text{len}(y_m) - \text{len}(y_b)$):
\begin{equation}
  \hat{P}(y_m \succ y_b \mid x) = \sigma\Big( \underbrace{\theta_m - \theta_b}_{\text{Chất lượng thực}} + \underbrace{w_x \cdot \psi(x)}_{\text{Độ khó câu hỏi}} + \underbrace{\beta_{\text{len}} \cdot \tanh\big(\alpha (\text{len}(y_m) - \text{len}(y_b))\big)}_{\text{Phần điểm ăn gian nhờ viết dài}} \Big)
  \label{eq:lc_alpacaeval}
\end{equation}
Sau khi học xong các hệ số từ dữ liệu chấm điểm, \textbf{LC-AlpacaEval} tính tỷ lệ thắng công bằng bằng cách \textbf{gán cưỡng bức chênh lệch độ dài bằng 0} ($\text{len}(y_m) - \text{len}(y_b) = 0$), trả lời chính xác câu hỏi phản thực tế: \textit{``Nếu bắt mô hình của bạn phải viết ngắn gọn đúng bằng mô hình chuẩn, nó sẽ thắng được bao nhiêu phần trăm?''}

% =========================================================================
\section{Playbook Gỡ lỗi Thực chiến (10 Tình huống Cấp cứu cho AI Engineer)}
\label{sec:debugging_playbook}

Để khép lại cuốn cẩm nang này, Bảng~\ref{tab:troubleshooting_playbook} tổng hợp \textbf{10 sự cố thực tế phổ biến nhất} mà mọi kỹ sư Post-Training đều sẽ gặp phải ít nhất một lần trong đời, kèm theo chẩn đoán nguyên nhân toán học/hệ thống và phác đồ xử lý tức thì.

\begin{table}[htbp]
\centering
\caption{Bảng Cấp cứu \& Gỡ lỗi Hệ thống Post-Training (SFT, DPO, PPO, GRPO)}
\label{tab:troubleshooting_playbook}
\renewcommand{\arraystretch}{1.38}
\footnotesize
\begin{tabularx}{\textwidth}{p{3.4cm} p{4.5cm} X}
\toprule
\textbf{Triệu chứng lâm sàng} & \textbf{Nguyên nhân gốc rễ (Toán / Hệ thống)} & \textbf{Phác đồ xử lý tức thì cho AI Engineer} \\
\midrule
\textbf{1. [SFT] Mô hình trả lời xong không chịu dừng}, cứ tự sinh tiếp \texttt{User: ... Assistant: ...} đến kịch trần độ dài. & Khi làm Prompt Masking (Chương~\ref{ch:sft}), bạn đã vô tình gán nhãn \texttt{-100} đè lên cả token kết thúc câu \texttt{<|im\_end|>} hoặc \texttt{<|eot\_id|>}. & Kiểm tra lại mảng \texttt{labels}: bắt buộc vị trí token đóng lượt của Assistant phải mang ID thật (khác \texttt{-100}) để mô hình học được điểm dừng. \\
\addlinespace
\textbf{2. [SFT] Loss nổ thành \texttt{NaN}} ngay ở bước 1--5 sau khi thêm Special Tokens mới vào Tokenizer. & Vector embedding của các token mới được khởi tạo ngẫu nhiên với phương sai lớn, làm nổ tích vô hướng ở lớp cuối (Chương~\ref{ch:sft}). & Khởi tạo embedding của các token mới bằng trung bình cộng (\texttt{mean}) của toàn bộ ma trận embedding cũ (Phương trình~\ref{eq:mean_init_eng}). \\
\addlinespace
\textbf{3. [SFT Packing] Mô hình trả lời râu ông nọ cắm cằm bà kia} (trộn nội dung câu trước vào câu sau). & Ghép nối nhiều hội thoại (Sequence Packing) nhưng vẫn dùng mặt nạ tam giác chuẩn và quên reset \texttt{position\_ids}. & Bật \texttt{flash\_attn\_varlen\_func} với mảng \texttt{cu\_seqlens} (Block-Diagonal Mask) và reset \texttt{position\_ids} về \texttt{0} ở đầu mỗi mẫu. \\
\addlinespace
\textbf{4. [DPO] Khoảng cách thưởng $\hat{r}_w - \hat{r}_l$ tăng cao nhưng mô hình sinh ra toàn ký tự rác.} & Bệnh \textbf{Likelihood Displacement} (Chương~\ref{ch:dpo}): $y_w$ và $y_l$ trùng nhiều từ đầu câu khiến cả $\log \pi_\theta(y_w)$ lẫn $\log \pi_\theta(y_l)$ cùng tụt xuống âm vô cùng. & Thêm thành phần neo SFT (NLL Loss) trên câu thắng $y_w$ với hệ số $\alpha = 0.1$ (Phương trình~\ref{eq:dpo_nll_loss} --- RPO) và giảm Learning Rate xuống $\le 10^{-6}$. \\
\addlinespace
\textbf{5. [DPO] Mô hình sau khi căn chỉnh nói dài dòng gấp 3 lần} bình thường. & Thiên kiến độ dài do DPO cộng tổng log-prob trên toàn câu mà không chia trung bình theo số lượng token $|y|$. & Chuyển từ DPO chuẩn sang \textbf{SimPO} \citep{meng2024simpo} (chia trung bình theo $\frac{1}{|y|}$ kèm biên $\gamma$, Phương trình~\ref{eq:simpo_loss}). \\
\addlinespace
\textbf{6. [PPO / GRPO] Khoản phạt KL nhảy ra số âm} trên log và quá trình huấn luyện sập. & Bạn đang tính KL trên từng mẫu bằng công thức ngây thơ $k_1 = \log(\pi_\theta / \pi_{\text{ref}})$, vốn nhận giá trị âm khi $\pi_\theta < \pi_{\text{ref}}$. & Thay ngay bằng ước lượng không chệch luôn dương $k_3 = r - \log r - 1$ của Schulman (Phương trình~\ref{eq:schulman_k3}). \\
\addlinespace
\textbf{7. [GRPO] Các câu trả lời SAI ngày càng phình dài ra} tới kịch trần $L_{\max}$, trong khi câu đúng ngắn lại. & Phép chia $\frac{1}{|o_i|}$ trong GRPO gốc khiến mô hình cố tình kéo dài câu sai để pha loãng hình phạt âm $\hat{A}_i < 0$ trên mỗi token. & Áp dụng công thức \textbf{Dr.~GRPO} \citep{liu2025understanding}: thay mẫu số biến thiên $|o_i|$ bằng hằng số cố định $L_{\max}$ và bỏ chia $\std(r)$ (Phương trình~\ref{eq:dr_grpo}). \\
\addlinespace
\textbf{8. [GRPO] Entropy tụt về 0 quá sớm}, mọi câu trong nhóm $G$ giống hệt nhau, điểm số ngừng tăng. & Ngưỡng cắt trên đối xứng $\epsilon = 0.2$ khóa chặt các token khám phá xác suất thấp; nhiều nhóm có $\std(r) = 0$ làm mất đạo hàm. & Áp dụng \textbf{DAPO} \citep{yu2025dapo}: nâng ngưỡng cắt trên bất đối xứng $\epsilon_{\text{high}} = 0.28$ (Clip-Higher) và bật Dynamic Sampling lọc bỏ các nhóm có $\std(r) = 0$. \\
\addlinespace
\textbf{9. [Reasoning RL] Hỏi ``1+1 bằng mấy'' cũng nháp 2,000 chữ} trong thẻ \texttt{<think>}. & Bệnh \textbf{Overthinking} \citep{chen2025donot} do hàm thưởng chỉ chấm đúng/sai mà không phạt chi phí token trên các câu hỏi dễ. & Thêm hình phạt độ dài có điều kiện cho các nhóm giải đúng 100\%, hoặc áp dụng \textbf{Budget Forcing} \citep{muennighoff2025s1} lúc Inference. \\
\addlinespace
\textbf{10. [All] Tràn VRAM (OOM) ở bước tính Loss} dù đã bật ZeRO-3 và Batch Size $= 1$. & Ma trận Logits kích thước $(B, T, |\mathcal{V}|)$ ở lớp đầu ra từ vựng (LM Head) phình lên hàng chục GB khi $T \ge 8{,}192$ và $|\mathcal{V}| \ge 128{,}000$. & Cài đặt và kích hoạt \texttt{Liger-Kernel} (\texttt{FusedLinearCrossEntropy} / Chunked Logits) để cắt nhỏ chiều thời gian $T$ trước khi nhân với bảng từ vựng. \\
\bottomrule
\end{tabularx}
\end{table}


% ============================================================================
% TÀI LIỆU THAM KHẢO
% ============================================================================
\bibliographystyle{plainnat}
\bibliography{references}

\end{document}
